\documentclass[10pt,letterpaper,oneside]{article}
\usepackage[margin=0.56in,headheight=14pt,footskip=18pt]{geometry}
\usepackage{fontspec}
\setmainfont{DejaVu Serif}
\setsansfont{Lato}
\setmonofont{DejaVu Sans Mono}
\usepackage{xcolor}
\usepackage{hyperref}
\usepackage{fancyhdr}
\usepackage{enumitem}
\usepackage{booktabs,longtable,tabularx,array}
\usepackage{amsmath,amssymb}
\usepackage{tcolorbox}
\usepackage{listings}
\usepackage{microtype}
\usepackage{titlesec}
\usepackage{lastpage}
\definecolor{gardenGreen}{HTML}{2F6B4F}
\definecolor{gardenDeep}{HTML}{173D2E}
\definecolor{gardenGold}{HTML}{B37A22}
\definecolor{gardenBlue}{HTML}{2D6077}
\definecolor{gardenMist}{HTML}{F3F7F4}
\definecolor{gardenLine}{HTML}{C8D5CD}
\definecolor{gardenInk}{HTML}{22312B}
\hypersetup{colorlinks=true,linkcolor=gardenGreen,urlcolor=gardenBlue,pdfauthor={LeafOS Project},pdftitle={The Garden - Lifecycle Architecture for LeafOS and FlowerOS}}
\pagestyle{fancy}
\fancyhf{}
\lhead{\sffamily\small The Garden}
\rhead{\sffamily\small LeafOS / FlowerOS Lifecycle Architecture}
\cfoot{\sffamily\small Page \thepage\ of \pageref{LastPage}}
\setlength{\parindent}{0pt}
\setlength{\parskip}{2.8pt}
\setlist[itemize]{leftmargin=1.3em,itemsep=0.7pt,topsep=1.5pt,parsep=0pt}
\lstset{basicstyle=\ttfamily\scriptsize,breaklines=true,frame=single,rulecolor=\color{gardenLine},backgroundcolor=\color{gardenMist},xleftmargin=4pt,xrightmargin=4pt,aboveskip=5pt,belowskip=5pt}
\newtcolorbox{corebox}{colback=gardenGreen!7,colframe=gardenGreen,boxrule=0.8pt,arc=2pt,left=6pt,right=6pt,top=3pt,bottom=3pt}
\newtcolorbox{acceptbox}{colback=gardenGold!8,colframe=gardenGold,boxrule=0.7pt,arc=2pt,left=6pt,right=6pt,top=3pt,bottom=3pt}
\newcommand{\prov}[1]{\textcolor{gardenBlue}{\sffamily\footnotesize\textbf{Provenance:} #1}}
\newcommand{\platetitle}[3]{%
  \begin{minipage}[t][0.54in][t]{\textwidth}
  {\sffamily\bfseries\color{gardenDeep}\fontsize{9}{10}\selectfont DESIGN PLATE #1}\\[-1pt]
  {\sffamily\bfseries\color{gardenInk}\fontsize{16.5}{18}\selectfont #2}\hfill{\sffamily\color{gardenGreen}\large #3}
  \end{minipage}\vspace{-2pt}\hrule\vspace{5pt}}
\begin{document}
\platetitle{1}{Title and Thesis}{FIFTH I}
\prov{[G] New synthesis}
\fontsize{8.85}{10.45}\selectfont
\begin{corebox}
\textbf{Core statement.} The Garden is a lifecycle architecture for LeafOS and FlowerOS. It replaces the idea of permanent bloom with bounded cycles of activation, work, validation, shutdown, composting, and renewal.
\end{corebox}
The inherited LeafOS architecture already separates machine control, work control, model proposals, native validation, durable evidence, and checkpoints. The Garden does not discard those boundaries. It gives them a temporal structure: every active component has a beginning, a useful period, an end condition, and a durable residue.

\textbf{Garden consequence.} The design target is renewable local intelligence rather than an indefinitely enlarged resident process. Power, compute, data, cognition, and durable output are treated as distinct biological roles so that each can be measured, governed, and retired independently.

\textbf{Operational rules}
\begin{itemize}
\item Preserve FlowerOS as machine authority.
\item Preserve LeafOS as work authority.
\item Treat Bloom as one phase, not the whole system.
\item Require every active phase to define fruit, senescence, and recovery.
\end{itemize}
\textbf{Extended analysis.} Authority analysis. Title and Thesis must be represented as explicit runtime data rather than left as an architectural implication. Ownership, legal inputs, legal outputs, mutation rights, and restart behavior should be visible in the journal and testable through the native inlet. The crucial invariant is that presentation, model confidence, or process longevity cannot strengthen authority. The implementation should therefore use the same disposition vocabulary - accepted, rejected, blocked, stale, failed, and committed - across shells, providers, interfaces, and lifecycle controls.

\textbf{Failure signatures}
\begin{itemize}
\item Two components claim overlapping ownership of title and thesis.
\item A surface reports success without a native disposition or evidence reference.
\item Restart reconstructs a different authority state than the committed checkpoint.
\item A model proposal crosses into mutation without a typed capability.
\end{itemize}
\textbf{Implementation or verification note.} Implementation note. The minimum useful implementation should encode the rules above in configuration or schema, expose them through status and doctor commands, and add at least one negative test that attempts to bypass the boundary. The expected durable result is not merely a successful command; it is an inspectable event trail ending in the acceptance condition for title and thesis.

\begin{lstlisting}
ROOTS -> PETALS -> POLLEN -> LEAVES -> BLOOM -> FRUIT
   ^                                               |
   +--------------- COMPOST <- SENESCENCE ---------+
\end{lstlisting}
\begin{acceptbox}
\textbf{Acceptance statement.} The thesis is accepted when the reader can explain why a smaller cyclic system may outgrow a larger permanent allocation over repeated runs.
\end{acceptbox}
\newpage
\platetitle{2}{Abstract}{FIFTH I}
\prov{[AA][CB][TP][G]}
\fontsize{8.85}{10.45}\selectfont
\begin{corebox}
\textbf{Core statement.} This document unifies the LeafOS authority model, continual transcript, MWF personalities, evidence pipeline, hardware throughput model, and model-pack system under a new Garden lifecycle doctrine.
\end{corebox}
Three fifths of the design are inherited from the Architecture Atlas, Continual Bloom Primary Reference, and Inference Throughput White Paper. Two fifths are new: a biological systems model and an implementation contract for finite blooms, compute petals, data pollen, power roots, brain-like leaves, durable fruit, and compostable state.

\textbf{Garden consequence.} The new philosophy changes the optimization target. The system no longer asks how much model state can remain resident forever. It asks how much validated work can be produced per cycle, per joule, per unit of memory pressure, and per unit of recoverable state.

\textbf{Operational rules}
\begin{itemize}
\item Raw token speed remains a measurement, not the primary objective.
\item A model proposal never becomes authority by being active longer.
\item Idle, dormant, and dead states are valid lifecycle states.
\item Growth is measured across cycles, not by peak allocation alone.
\end{itemize}
\textbf{Extended analysis.} Authority analysis. Abstract must be represented as explicit runtime data rather than left as an architectural implication. Ownership, legal inputs, legal outputs, mutation rights, and restart behavior should be visible in the journal and testable through the native inlet. The crucial invariant is that presentation, model confidence, or process longevity cannot strengthen authority. The implementation should therefore use the same disposition vocabulary - accepted, rejected, blocked, stale, failed, and committed - across shells, providers, interfaces, and lifecycle controls.

\textbf{Failure signatures}
\begin{itemize}
\item Two components claim overlapping ownership of abstract.
\item A surface reports success without a native disposition or evidence reference.
\item Restart reconstructs a different authority state than the committed checkpoint.
\item A model proposal crosses into mutation without a typed capability.
\end{itemize}
\textbf{Implementation or verification note.} Implementation note. The minimum useful implementation should encode the rules above in configuration or schema, expose them through status and doctor commands, and add at least one negative test that attempts to bypass the boundary. The expected durable result is not merely a successful command; it is an inspectable event trail ending in the acceptance condition for abstract.

\begin{acceptbox}
\textbf{Acceptance statement.} The abstract succeeds when it states both continuity and replacement: LeafOS contracts survive, Infinite Bloom does not remain the master philosophy.
\end{acceptbox}
\newpage
\platetitle{3}{Executive Summary}{FIFTH I}
\prov{[G] New synthesis over [AA][CB][TP]}
\fontsize{8.85}{10.45}\selectfont
\begin{corebox}
\textbf{Core statement.} The Garden turns a persistent orchestration stack into an ecology of bounded computational organisms.
\end{corebox}
Roots represent the electrical and thermal infrastructure that keeps the system alive. Petals are compute surfaces such as CPU, GPU, and NPU workers. Pollen is typed data exchanged between nodes. Leaves are the most complex brain-like transformation systems. Bloom is the period of active coordinated work. Fruit is validated durable output. Compost is compressed retired state that improves later cycles without remaining active memory.

\textbf{Garden consequence.} This architecture favors a small required core, optional seasonal lanes, hot/warm/cold/dormant placement, explicit unload rules, and strong state reconstruction. A model can die without the task dying because the authoritative present is reconstructed from contracts, evidence, checkpoints, and transcript tails.

\textbf{Operational rules}
\begin{itemize}
\item Use smaller resident allocations by default.
\item Escalate to larger models only when task value justifies cost.
\item Unload inactive workers after fruit or timeout.
\item Convert failures into classified compost rather than silent repetition.
\end{itemize}
\textbf{Extended analysis.} Authority analysis. Executive Summary must be represented as explicit runtime data rather than left as an architectural implication. Ownership, legal inputs, legal outputs, mutation rights, and restart behavior should be visible in the journal and testable through the native inlet. The crucial invariant is that presentation, model confidence, or process longevity cannot strengthen authority. The implementation should therefore use the same disposition vocabulary - accepted, rejected, blocked, stale, failed, and committed - across shells, providers, interfaces, and lifecycle controls.

\textbf{Failure signatures}
\begin{itemize}
\item Two components claim overlapping ownership of executive summary.
\item A surface reports success without a native disposition or evidence reference.
\item Restart reconstructs a different authority state than the committed checkpoint.
\item A model proposal crosses into mutation without a typed capability.
\end{itemize}
\textbf{Implementation or verification note.} Implementation note. The minimum useful implementation should encode the rules above in configuration or schema, expose them through status and doctor commands, and add at least one negative test that attempts to bypass the boundary. The expected durable result is not merely a successful command; it is an inspectable event trail ending in the acceptance condition for executive summary.

\begin{acceptbox}
\textbf{Acceptance statement.} The executive summary is complete when it produces one governing sentence: infinite growth comes from renewable cycles, not infinite resident size.
\end{acceptbox}
\newpage
\platetitle{4}{Five-Fifths Provenance}{FIFTH I}
\prov{[AA][CB][TP][G]}
\fontsize{8.85}{10.45}\selectfont
\begin{corebox}
\textbf{Core statement.} The document is deliberately divided into five equal design bands so that inherited architecture and new philosophy remain distinguishable.
\end{corebox}
Fifth One records machine and authority boundaries. Fifth Two records durable cognition, personalities, state, and evidence. Fifth Three records hardware, throughput, routing, and pack composition. Fifth Four introduces the Garden philosophy. Fifth Five converts that philosophy into runtime contracts, schemas, policies, and work orders.

\textbf{Garden consequence.} The ratio is not cosmetic. It prevents the new metaphor from erasing the technical work already present, while preventing the old Continual Bloom framing from absorbing the lifecycle insight and pretending nothing changed.

\textbf{Operational rules}
\begin{itemize}
\item Tag source-derived claims with [AA], [CB], or [TP].
\item Tag new normative material with [G].
\item Do not silently promote historical measurements into current guarantees.
\item Do not present metaphor as authority without a corresponding runtime contract.
\end{itemize}
\textbf{Extended analysis.} Authority analysis. Five-Fifths Provenance must be represented as explicit runtime data rather than left as an architectural implication. Ownership, legal inputs, legal outputs, mutation rights, and restart behavior should be visible in the journal and testable through the native inlet. The crucial invariant is that presentation, model confidence, or process longevity cannot strengthen authority. The implementation should therefore use the same disposition vocabulary - accepted, rejected, blocked, stale, failed, and committed - across shells, providers, interfaces, and lifecycle controls.

\textbf{Failure signatures}
\begin{itemize}
\item Two components claim overlapping ownership of five-fifths provenance.
\item A surface reports success without a native disposition or evidence reference.
\item Restart reconstructs a different authority state than the committed checkpoint.
\item A model proposal crosses into mutation without a typed capability.
\end{itemize}
\textbf{Implementation or verification note.} Implementation note. The minimum useful implementation should encode the rules above in configuration or schema, expose them through status and doctor commands, and add at least one negative test that attempts to bypass the boundary. The expected durable result is not merely a successful command; it is an inspectable event trail ending in the acceptance condition for five-fifths provenance.

\begin{lstlisting}
FIFTH I   Authority and machine             [source]
FIFTH II  Durable cognition and evidence    [source]
FIFTH III Hardware, throughput, packs       [source]
FIFTH IV  Garden philosophy                 [new]
FIFTH V   Garden implementation             [new]
\end{lstlisting}
\begin{acceptbox}
\textbf{Acceptance statement.} Provenance is accepted when every major rule can be traced either to an inherited document or to an explicitly new Garden decision.
\end{acceptbox}
\newpage
\platetitle{5}{Canonical Vocabulary}{FIFTH I}
\prov{[G] New synthesis}
\fontsize{8.85}{10.45}\selectfont
\begin{corebox}
\textbf{Core statement.} The Garden vocabulary is a typed systems vocabulary, not decorative branding.
\end{corebox}
Soil is the environmental envelope. Roots are power and survival infrastructure. The stem is the transport and scheduling spine. Petals are compute surfaces. Pollen is typed data in motion. Leaves are high-complexity transformation systems. Bloom is active coordinated execution. Fruit is validated durable output. Seeds are compact restartable intent. Senescence is controlled shutdown. Compost is retired state converted into future context.

\textbf{Garden consequence.} Each term must map to measurable state. A component may have a poetic label in the interface, but the journal stores explicit roles, timestamps, budgets, identities, and evidence references.

\textbf{Operational rules}
\begin{itemize}
\item Every biological term must have a machine definition.
\item No glyph, color, or name grants authority.
\item A lifecycle term must correspond to a state transition.
\item Unknown mappings remain unknown rather than being improvised.
\end{itemize}
\textbf{Extended analysis.} Authority analysis. Canonical Vocabulary must be represented as explicit runtime data rather than left as an architectural implication. Ownership, legal inputs, legal outputs, mutation rights, and restart behavior should be visible in the journal and testable through the native inlet. The crucial invariant is that presentation, model confidence, or process longevity cannot strengthen authority. The implementation should therefore use the same disposition vocabulary - accepted, rejected, blocked, stale, failed, and committed - across shells, providers, interfaces, and lifecycle controls.

\textbf{Failure signatures}
\begin{itemize}
\item Two components claim overlapping ownership of canonical vocabulary.
\item A surface reports success without a native disposition or evidence reference.
\item Restart reconstructs a different authority state than the committed checkpoint.
\item A model proposal crosses into mutation without a typed capability.
\end{itemize}
\textbf{Implementation or verification note.} Implementation note. The minimum useful implementation should encode the rules above in configuration or schema, expose them through status and doctor commands, and add at least one negative test that attempts to bypass the boundary. The expected durable result is not merely a successful command; it is an inspectable event trail ending in the acceptance condition for canonical vocabulary.

\begin{acceptbox}
\textbf{Acceptance statement.} Vocabulary is accepted when two engineers can translate a Garden diagram into process, memory, power, and data-flow records without guessing.
\end{acceptbox}
\newpage
\platetitle{6}{FlowerOS and LeafOS Split}{FIFTH I}
\prov{[AA][CB] Source-derived}
\fontsize{8.85}{10.45}\selectfont
\begin{corebox}
\textbf{Core statement.} FlowerOS controls the machine; LeafOS controls the work.
\end{corebox}
FlowerOS owns boot, drivers, device access, processes, services, host integration, power visibility, and terminal presentation. LeafOS owns task interpretation, routing, planning, guarded mutation, validation, evidence, checkpoints, recovery, and exactly one legal next action.

\textbf{Garden consequence.} Within the Garden, FlowerOS primarily supplies soil and roots: the physical conditions, power state, thermals, process substrate, and device availability. LeafOS primarily coordinates stems, pollen, leaves, bloom, fruit, and compost. This mapping does not merge their authority domains.

\textbf{Operational rules}
\begin{itemize}
\item FlowerOS may expose hardware facts but does not decide task truth.
\item LeafOS may schedule work but does not pretend to own drivers or power electronics.
\item Cross-boundary requests are typed and journaled.
\item Failure in one layer must remain visible to the other.
\end{itemize}
\textbf{Extended analysis.} Authority analysis. FlowerOS and LeafOS Split must be represented as explicit runtime data rather than left as an architectural implication. Ownership, legal inputs, legal outputs, mutation rights, and restart behavior should be visible in the journal and testable through the native inlet. The crucial invariant is that presentation, model confidence, or process longevity cannot strengthen authority. The implementation should therefore use the same disposition vocabulary - accepted, rejected, blocked, stale, failed, and committed - across shells, providers, interfaces, and lifecycle controls.

\textbf{Failure signatures}
\begin{itemize}
\item Two components claim overlapping ownership of floweros and leafos split.
\item A surface reports success without a native disposition or evidence reference.
\item Restart reconstructs a different authority state than the committed checkpoint.
\item A model proposal crosses into mutation without a typed capability.
\end{itemize}
\textbf{Implementation or verification note.} Implementation note. The minimum useful implementation should encode the rules above in configuration or schema, expose them through status and doctor commands, and add at least one negative test that attempts to bypass the boundary. The expected durable result is not merely a successful command; it is an inspectable event trail ending in the acceptance condition for floweros and leafos split.

\begin{acceptbox}
\textbf{Acceptance statement.} The split is accepted when a machine restart and a task restart can be reasoned about independently.
\end{acceptbox}
\newpage
\platetitle{7}{Two Surfaces, One Engine}{FIFTH I}
\prov{[AA] Source-derived}
\fontsize{8.85}{10.45}\selectfont
\begin{corebox}
\textbf{Core statement.} PowerShell and Bash are operator surfaces that delegate into one canonical implementation.
\end{corebox}
The surfaces provide host-specific convenience, installation, doctor commands, status views, and typed controls. They do not fork the taskpack. The canonical dispatch and resident engine live under ProjectLeaf, preserving one authority path even when the operator experience differs by shell.

\textbf{Garden consequence.} In Garden terms, surfaces are views into the organism, not separate organisms. A second shell must not create a second queue, second journal, or competing lifecycle governor.

\textbf{Operational rules}
\begin{itemize}
\item Both surfaces call the same canonical command semantics.
\item Surface-specific helpers may adapt paths, not policy.
\item A surface closing must not kill durable state.
\item A surface may observe or request but cannot rewrite evidence.
\end{itemize}
\textbf{Extended analysis.} Authority analysis. Two Surfaces, One Engine must be represented as explicit runtime data rather than left as an architectural implication. Ownership, legal inputs, legal outputs, mutation rights, and restart behavior should be visible in the journal and testable through the native inlet. The crucial invariant is that presentation, model confidence, or process longevity cannot strengthen authority. The implementation should therefore use the same disposition vocabulary - accepted, rejected, blocked, stale, failed, and committed - across shells, providers, interfaces, and lifecycle controls.

\textbf{Failure signatures}
\begin{itemize}
\item Two components claim overlapping ownership of two surfaces, one engine.
\item A surface reports success without a native disposition or evidence reference.
\item Restart reconstructs a different authority state than the committed checkpoint.
\item A model proposal crosses into mutation without a typed capability.
\end{itemize}
\textbf{Implementation or verification note.} Implementation note. The minimum useful implementation should encode the rules above in configuration or schema, expose them through status and doctor commands, and add at least one negative test that attempts to bypass the boundary. The expected durable result is not merely a successful command; it is an inspectable event trail ending in the acceptance condition for two surfaces, one engine.

\begin{lstlisting}
PowerShell surface ----\
                       > canonical leafctl -> resident engine
Bash surface ----------/
\end{lstlisting}
\begin{acceptbox}
\textbf{Acceptance statement.} The design passes when the same run can be inspected from PowerShell and Bash without divergent state.
\end{acceptbox}
\newpage
\platetitle{8}{Root Contract}{FIFTH I}
\prov{[AA] Source-derived}
\fontsize{8.85}{10.45}\selectfont
\begin{corebox}
\textbf{Core statement.} The root contract defines which directories are source, implementation, evidence, cache, import, experiment, or disposable state.
\end{corebox}
Stable entrypoints and documentation sit above ProjectLeaf. The taskpack owns resident execution. The model installer owns acquisition. Runs, reports, logs, and evidence preserve outcomes but do not replace source authority. Temporary directories, caches, virtual environments, and downloaded binaries are rebuildable.

\textbf{Garden consequence.} The Garden adds lifecycle classes to the root contract: dormant seed material, active bloom state, durable fruit, and compost archives. Classification must remain explicit so that a cleanup operation cannot prune source or mistake a cache for fruit.

\textbf{Operational rules}
\begin{itemize}
\item Every root path has one authority class.
\item Generated evidence is immutable or append-only after commitment.
\item Caches may be deleted without changing truth.
\item Imported material cannot silently override current contracts.
\end{itemize}
\textbf{Extended analysis.} Authority analysis. Root Contract must be represented as explicit runtime data rather than left as an architectural implication. Ownership, legal inputs, legal outputs, mutation rights, and restart behavior should be visible in the journal and testable through the native inlet. The crucial invariant is that presentation, model confidence, or process longevity cannot strengthen authority. The implementation should therefore use the same disposition vocabulary - accepted, rejected, blocked, stale, failed, and committed - across shells, providers, interfaces, and lifecycle controls.

\textbf{Failure signatures}
\begin{itemize}
\item Two components claim overlapping ownership of root contract.
\item A surface reports success without a native disposition or evidence reference.
\item Restart reconstructs a different authority state than the committed checkpoint.
\item A model proposal crosses into mutation without a typed capability.
\end{itemize}
\textbf{Implementation or verification note.} Implementation note. The minimum useful implementation should encode the rules above in configuration or schema, expose them through status and doctor commands, and add at least one negative test that attempts to bypass the boundary. The expected durable result is not merely a successful command; it is an inspectable event trail ending in the acceptance condition for root contract.

\begin{acceptbox}
\textbf{Acceptance statement.} The contract is accepted when a cleanup tool can identify disposable material without risking source, evidence, or checkpoints.
\end{acceptbox}
\newpage
\platetitle{9}{Operator Surfaces}{FIFTH I}
\prov{[AA] Source-derived}
\fontsize{8.85}{10.45}\selectfont
\begin{corebox}
\textbf{Core statement.} Operator surfaces present normalized state and emit typed requests.
\end{corebox}
The TUI, HTML viewer, and shell commands are projections over the same structured snapshot. They may display tasks, hardware, queue state, provider progress, ledger entries, results, and cursor position. They do not parse one another, schedule hidden work, or execute arbitrary shell text.

\textbf{Garden consequence.} Garden presentation may show growth, dormancy, pruning, or bloom, but those visual states must be derived from journaled lifecycle facts. A green petal icon cannot substitute for a device ownership record.

\textbf{Operational rules}
\begin{itemize}
\item All views consume one native snapshot.
\item Controls emit named capabilities with source, nonce, and read head.
\item Closing a view closes only the view.
\item Missing telemetry is displayed as unavailable, not estimated.
\end{itemize}
\textbf{Extended analysis.} Authority analysis. Operator Surfaces must be represented as explicit runtime data rather than left as an architectural implication. Ownership, legal inputs, legal outputs, mutation rights, and restart behavior should be visible in the journal and testable through the native inlet. The crucial invariant is that presentation, model confidence, or process longevity cannot strengthen authority. The implementation should therefore use the same disposition vocabulary - accepted, rejected, blocked, stale, failed, and committed - across shells, providers, interfaces, and lifecycle controls.

\textbf{Failure signatures}
\begin{itemize}
\item Two components claim overlapping ownership of operator surfaces.
\item A surface reports success without a native disposition or evidence reference.
\item Restart reconstructs a different authority state than the committed checkpoint.
\item A model proposal crosses into mutation without a typed capability.
\end{itemize}
\textbf{Implementation or verification note.} Implementation note. The minimum useful implementation should encode the rules above in configuration or schema, expose them through status and doctor commands, and add at least one negative test that attempts to bypass the boundary. The expected durable result is not merely a successful command; it is an inspectable event trail ending in the acceptance condition for operator surfaces.

\begin{acceptbox}
\textbf{Acceptance statement.} Surfaces pass when their rendered states agree with the journal and a stale control request is rejected.
\end{acceptbox}
\newpage
\platetitle{10}{ProjectLeaf as Implementation Authority}{FIFTH I}
\prov{[AA] Source-derived}
\fontsize{8.85}{10.45}\selectfont
\begin{corebox}
\textbf{Core statement.} ProjectLeaf is the implementation owner for task, queue, validation, reporting, provider, UI, and resident-loop behavior.
\end{corebox}
This concentration of authority prevents shell helpers, model caches, imported experiments, or presentation code from becoming competing runtimes. ProjectLeaf may expose modular subsystems, but the legal execution path remains bounded and inspectable.

\textbf{Garden consequence.} The Garden architecture should therefore be implemented as policy and state inside ProjectLeaf rather than as a decorative sidecar. Lifecycle transitions must pass through the same inlet, journal, validator, and checkpoint machinery as existing tasks.

\textbf{Operational rules}
\begin{itemize}
\item Garden states are first-class runtime records.
\item No separate daemon may mutate Garden state outside the inlet.
\item New schemas live beside existing taskpack schemas.
\item Experimental Garden adapters remain non-authoritative until promoted.
\end{itemize}
\textbf{Extended analysis.} Authority analysis. ProjectLeaf as Implementation Authority must be represented as explicit runtime data rather than left as an architectural implication. Ownership, legal inputs, legal outputs, mutation rights, and restart behavior should be visible in the journal and testable through the native inlet. The crucial invariant is that presentation, model confidence, or process longevity cannot strengthen authority. The implementation should therefore use the same disposition vocabulary - accepted, rejected, blocked, stale, failed, and committed - across shells, providers, interfaces, and lifecycle controls.

\textbf{Failure signatures}
\begin{itemize}
\item Two components claim overlapping ownership of projectleaf as implementation authority.
\item A surface reports success without a native disposition or evidence reference.
\item Restart reconstructs a different authority state than the committed checkpoint.
\item A model proposal crosses into mutation without a typed capability.
\end{itemize}
\textbf{Implementation or verification note.} Implementation note. The minimum useful implementation should encode the rules above in configuration or schema, expose them through status and doctor commands, and add at least one negative test that attempts to bypass the boundary. The expected durable result is not merely a successful command; it is an inspectable event trail ending in the acceptance condition for projectleaf as implementation authority.

\begin{acceptbox}
\textbf{Acceptance statement.} Implementation authority is accepted when Garden lifecycle behavior can be disabled without changing the existing evidence hierarchy.
\end{acceptbox}
\newpage
\platetitle{11}{Model Acquisition Boundary}{FIFTH I}
\prov{[AA][CB] Source-derived}
\fontsize{8.85}{10.45}\selectfont
\begin{corebox}
\textbf{Core statement.} Model catalog, planning, resolution, application, and verification are distinct operations.
\end{corebox}
Catalog and doctor operations are local. Planning produces an offline intent. Resolution may contact remote metadata and pin revisions. Apply is the only weight-download boundary. Verification confirms files, hashes, paths, and role links. Credentials must fail fast before workers start.

\textbf{Garden consequence.} In the Garden, a model file is dormant seed stock until a pack and lifecycle governor activate it. Downloading a model does not make it a petal, and loading a model does not make it a leaf.

\textbf{Operational rules}
\begin{itemize}
\item Downloads require explicit apply.
\item Resolved revisions and filenames are preserved.
\item Invalid tokens fail before parallel work begins.
\item Installed, declared, and loaded are separate states.
\end{itemize}
\textbf{Extended analysis.} Authority analysis. Model Acquisition Boundary must be represented as explicit runtime data rather than left as an architectural implication. Ownership, legal inputs, legal outputs, mutation rights, and restart behavior should be visible in the journal and testable through the native inlet. The crucial invariant is that presentation, model confidence, or process longevity cannot strengthen authority. The implementation should therefore use the same disposition vocabulary - accepted, rejected, blocked, stale, failed, and committed - across shells, providers, interfaces, and lifecycle controls.

\textbf{Failure signatures}
\begin{itemize}
\item Two components claim overlapping ownership of model acquisition boundary.
\item A surface reports success without a native disposition or evidence reference.
\item Restart reconstructs a different authority state than the committed checkpoint.
\item A model proposal crosses into mutation without a typed capability.
\end{itemize}
\textbf{Implementation or verification note.} Implementation note. The minimum useful implementation should encode the rules above in configuration or schema, expose them through status and doctor commands, and add at least one negative test that attempts to bypass the boundary. The expected durable result is not merely a successful command; it is an inspectable event trail ending in the acceptance condition for model acquisition boundary.

\begin{acceptbox}
\textbf{Acceptance statement.} The boundary passes when a plan can be inspected without network transfer and an apply operation can be limited to one selected role.
\end{acceptbox}
\newpage
\platetitle{12}{Local Model as Proposal Lane}{FIFTH I}
\prov{[AA][CB] Source-derived}
\fontsize{8.85}{10.45}\selectfont
\begin{corebox}
\textbf{Core statement.} A local model proposes bounded structured content; it does not become a control lane.
\end{corebox}
Instructions are normalized into project facts and task context. A provider returns schema-constrained proposals. The CPU inlet checks paths, commands, approvals, attempts, ownership, and current state. Tracked execution and validation produce evidence. Only accepted checkpoints become durable continuation.

\textbf{Garden consequence.} Garden language does not soften this asymmetry. A leaf may be the most brain-like component, but it remains interpretive. It may shape a bloom; it cannot declare its own fruit valid.

\textbf{Operational rules}
\begin{itemize}
\item Model output begins unverified.
\item Provider failure remains visible and never becomes mock output.
\item Schema validity is necessary but not sufficient.
\item The model cannot approve itself.
\end{itemize}
\textbf{Extended analysis.} Authority analysis. Local Model as Proposal Lane must be represented as explicit runtime data rather than left as an architectural implication. Ownership, legal inputs, legal outputs, mutation rights, and restart behavior should be visible in the journal and testable through the native inlet. The crucial invariant is that presentation, model confidence, or process longevity cannot strengthen authority. The implementation should therefore use the same disposition vocabulary - accepted, rejected, blocked, stale, failed, and committed - across shells, providers, interfaces, and lifecycle controls.

\textbf{Failure signatures}
\begin{itemize}
\item Two components claim overlapping ownership of local model as proposal lane.
\item A surface reports success without a native disposition or evidence reference.
\item Restart reconstructs a different authority state than the committed checkpoint.
\item A model proposal crosses into mutation without a typed capability.
\end{itemize}
\textbf{Implementation or verification note.} Implementation note. The minimum useful implementation should encode the rules above in configuration or schema, expose them through status and doctor commands, and add at least one negative test that attempts to bypass the boundary. The expected durable result is not merely a successful command; it is an inspectable event trail ending in the acceptance condition for local model as proposal lane.

\begin{acceptbox}
\textbf{Acceptance statement.} The proposal lane is accepted when a persuasive but invalid model response is rejected without mutating the project.
\end{acceptbox}
\newpage
\platetitle{13}{Native Inlet}{FIFTH I}
\prov{[AA][CB] Source-derived}
\fontsize{8.85}{10.45}\selectfont
\begin{corebox}
\textbf{Core statement.} The native inlet is the legal boundary between proposals and actions.
\end{corebox}
It evaluates capability name, run state, transcript head, nonce, approval policy, path scope, command allowlist, resource budget, and ownership. It appends an accepted, rejected, blocked, or stale disposition before execution.

\textbf{Garden consequence.} Within the Garden, the inlet is the gate between pollen and metabolism. Data may arrive from any node, but only typed, current, authorized packets can trigger work.

\textbf{Operational rules}
\begin{itemize}
\item Reject unknown capabilities.
\item Reject stale read heads.
\item Reject replayed nonces.
\item Reject paths or commands outside declared scope.
\item Record every disposition.
\end{itemize}
\textbf{Extended analysis.} Authority analysis. Native Inlet must be represented as explicit runtime data rather than left as an architectural implication. Ownership, legal inputs, legal outputs, mutation rights, and restart behavior should be visible in the journal and testable through the native inlet. The crucial invariant is that presentation, model confidence, or process longevity cannot strengthen authority. The implementation should therefore use the same disposition vocabulary - accepted, rejected, blocked, stale, failed, and committed - across shells, providers, interfaces, and lifecycle controls.

\textbf{Failure signatures}
\begin{itemize}
\item Two components claim overlapping ownership of native inlet.
\item A surface reports success without a native disposition or evidence reference.
\item Restart reconstructs a different authority state than the committed checkpoint.
\item A model proposal crosses into mutation without a typed capability.
\end{itemize}
\textbf{Implementation or verification note.} Implementation note. The minimum useful implementation should encode the rules above in configuration or schema, expose them through status and doctor commands, and add at least one negative test that attempts to bypass the boundary. The expected durable result is not merely a successful command; it is an inspectable event trail ending in the acceptance condition for native inlet.

\begin{acceptbox}
\textbf{Acceptance statement.} The inlet passes when the same request produces the same legal disposition after restart from committed state.
\end{acceptbox}
\newpage
\platetitle{14}{Typed Capability Requests}{FIFTH I}
\prov{[AA][CB] Source-derived}
\fontsize{8.85}{10.45}\selectfont
\begin{corebox}
\textbf{Core statement.} Interactive controls and model recommendations become typed capability requests rather than executable text.
\end{corebox}
A capability record names the action, source, nonce, read head, payload, requested scope, and expected state. The inlet converts it into a disposition. If accepted, a tracked executor receives a bounded work order. If rejected, the reason remains inspectable.

\textbf{Garden consequence.} Pollen is therefore typed, not merely textual. The Garden cannot safely route arbitrary prose as if it were metabolic instruction.

\textbf{Operational rules}
\begin{itemize}
\item Separate conversation from capability requests.
\item Require freshness at acceptance time.
\item Attach source identity and run identity.
\item Preserve rejection reasons as evidence.
\end{itemize}
\textbf{Extended analysis.} Authority analysis. Typed Capability Requests must be represented as explicit runtime data rather than left as an architectural implication. Ownership, legal inputs, legal outputs, mutation rights, and restart behavior should be visible in the journal and testable through the native inlet. The crucial invariant is that presentation, model confidence, or process longevity cannot strengthen authority. The implementation should therefore use the same disposition vocabulary - accepted, rejected, blocked, stale, failed, and committed - across shells, providers, interfaces, and lifecycle controls.

\textbf{Failure signatures}
\begin{itemize}
\item Two components claim overlapping ownership of typed capability requests.
\item A surface reports success without a native disposition or evidence reference.
\item Restart reconstructs a different authority state than the committed checkpoint.
\item A model proposal crosses into mutation without a typed capability.
\end{itemize}
\textbf{Implementation or verification note.} Implementation note. The minimum useful implementation should encode the rules above in configuration or schema, expose them through status and doctor commands, and add at least one negative test that attempts to bypass the boundary. The expected durable result is not merely a successful command; it is an inspectable event trail ending in the acceptance condition for typed capability requests.

\begin{lstlisting}
{ capability, source, nonce, read_head, payload }
                 |
                 v
       native disposition -> execute | reject | block | stale
\end{lstlisting}
\begin{acceptbox}
\textbf{Acceptance statement.} The capability system passes when optional chat cannot execute a command without a separate typed proposal.
\end{acceptbox}
\newpage
\platetitle{15}{Tracked Execution Sandbox}{FIFTH I}
\prov{[AA][CB] Source-derived}
\fontsize{8.85}{10.45}\selectfont
\begin{corebox}
\textbf{Core statement.} Execution occurs through tracked, allowlisted operations with captured identity and output.
\end{corebox}
Each child process record contains executable path, argument vector, process identifier, authorization capability, start and end time, exit code, stdout and stderr references, cancellation state, and resource observations. On Windows, hidden execution must preserve evidence rather than discarding it.

\textbf{Garden consequence.} Garden petals are compute surfaces, not uncontrolled subprocesses. Every opened petal must have an owner, budget, purpose, and closure path.

\textbf{Operational rules}
\begin{itemize}
\item Launch directly rather than through opaque shell chains when possible.
\item Capture stdout, stderr, and exit code.
\item Use job ownership for cleanup.
\item Do not confuse hidden windows with hidden evidence.
\end{itemize}
\textbf{Extended analysis.} Authority analysis. Tracked Execution Sandbox must be represented as explicit runtime data rather than left as an architectural implication. Ownership, legal inputs, legal outputs, mutation rights, and restart behavior should be visible in the journal and testable through the native inlet. The crucial invariant is that presentation, model confidence, or process longevity cannot strengthen authority. The implementation should therefore use the same disposition vocabulary - accepted, rejected, blocked, stale, failed, and committed - across shells, providers, interfaces, and lifecycle controls.

\textbf{Failure signatures}
\begin{itemize}
\item Two components claim overlapping ownership of tracked execution sandbox.
\item A surface reports success without a native disposition or evidence reference.
\item Restart reconstructs a different authority state than the committed checkpoint.
\item A model proposal crosses into mutation without a typed capability.
\end{itemize}
\textbf{Implementation or verification note.} Implementation note. The minimum useful implementation should encode the rules above in configuration or schema, expose them through status and doctor commands, and add at least one negative test that attempts to bypass the boundary. The expected durable result is not merely a successful command; it is an inspectable event trail ending in the acceptance condition for tracked execution sandbox.

\begin{acceptbox}
\textbf{Acceptance statement.} The sandbox passes when a cancelled process leaves a complete record and no orphaned worker.
\end{acceptbox}
\newpage
\platetitle{16}{Validation}{FIFTH I}
\prov{[AA][CB] Source-derived}
\fontsize{8.85}{10.45}\selectfont
\begin{corebox}
\textbf{Core statement.} Validation converts checkable outcomes into evidence; it does not convert failure into success.
\end{corebox}
Validators run declared commands, inspect files, compare hashes, enforce schemas, measure outputs, and preserve artifacts. A validator may support, refute, or leave a claim unresolved. It does not rewrite the claim to fit the desired result.

\textbf{Garden consequence.} Fruit requires validation. An untested patch is a flower, not fruit. A benchmark without provenance is pollen, not nourishment.

\textbf{Operational rules}
\begin{itemize}
\item Acceptance commands are declared before execution.
\item Required failures block commitment.
\item Evidence includes provenance and freshness.
\item Unknown results remain unknown.
\end{itemize}
\textbf{Extended analysis.} Authority analysis. Validation must be represented as explicit runtime data rather than left as an architectural implication. Ownership, legal inputs, legal outputs, mutation rights, and restart behavior should be visible in the journal and testable through the native inlet. The crucial invariant is that presentation, model confidence, or process longevity cannot strengthen authority. The implementation should therefore use the same disposition vocabulary - accepted, rejected, blocked, stale, failed, and committed - across shells, providers, interfaces, and lifecycle controls.

\textbf{Failure signatures}
\begin{itemize}
\item Two components claim overlapping ownership of validation.
\item A surface reports success without a native disposition or evidence reference.
\item Restart reconstructs a different authority state than the committed checkpoint.
\item A model proposal crosses into mutation without a typed capability.
\end{itemize}
\textbf{Implementation or verification note.} Implementation note. The minimum useful implementation should encode the rules above in configuration or schema, expose them through status and doctor commands, and add at least one negative test that attempts to bypass the boundary. The expected durable result is not merely a successful command; it is an inspectable event trail ending in the acceptance condition for validation.

\begin{acceptbox}
\textbf{Acceptance statement.} Validation passes when a failed acceptance command prevents checkpoint promotion and remains visible in the report.
\end{acceptbox}
\newpage
\platetitle{17}{Evidence Hierarchy}{FIFTH I}
\prov{[AA][CB] Source-derived}
\fontsize{8.85}{10.45}\selectfont
\begin{corebox}
\textbf{Core statement.} Records, evidence, and committed facts occupy different authority levels.
\end{corebox}
Requests, observations, plans, reflections, and tool attempts are records. Validator results with artifacts are evidence. Checkpoint-bound state is committed fact for resume purposes. A hash chain detects mutation but does not make the contents true.

\textbf{Garden consequence.} The Garden uses the same hierarchy across lifecycle phases. Pollen carries records; fruit requires evidence; seeds may summarize facts but must retain references to their validating artifacts.

\textbf{Operational rules}
\begin{itemize}
\item Never promote a record directly to committed fact.
\item Keep evidence references stable.
\item Treat hashes as integrity checks, not truth certificates.
\item Rebuild projections from the journal.
\end{itemize}
\textbf{Extended analysis.} Authority analysis. Evidence Hierarchy must be represented as explicit runtime data rather than left as an architectural implication. Ownership, legal inputs, legal outputs, mutation rights, and restart behavior should be visible in the journal and testable through the native inlet. The crucial invariant is that presentation, model confidence, or process longevity cannot strengthen authority. The implementation should therefore use the same disposition vocabulary - accepted, rejected, blocked, stale, failed, and committed - across shells, providers, interfaces, and lifecycle controls.

\textbf{Failure signatures}
\begin{itemize}
\item Two components claim overlapping ownership of evidence hierarchy.
\item A surface reports success without a native disposition or evidence reference.
\item Restart reconstructs a different authority state than the committed checkpoint.
\item A model proposal crosses into mutation without a typed capability.
\end{itemize}
\textbf{Implementation or verification note.} Implementation note. The minimum useful implementation should encode the rules above in configuration or schema, expose them through status and doctor commands, and add at least one negative test that attempts to bypass the boundary. The expected durable result is not merely a successful command; it is an inspectable event trail ending in the acceptance condition for evidence hierarchy.

\begin{acceptbox}
\textbf{Acceptance statement.} The hierarchy passes when the system can explain why a claim is believed, not merely where it was written.
\end{acceptbox}
\newpage
\platetitle{18}{Checkpoint Commit}{FIFTH I}
\prov{[AA][CB] Source-derived}
\fontsize{8.85}{10.45}\selectfont
\begin{corebox}
\textbf{Core statement.} A checkpoint is the durable boundary after successful validation.
\end{corebox}
It records the authoritative resume position, accepted task state, evidence references, unresolved blockers, resource disposition, and exactly one next action. It is created only after required validation succeeds.

\textbf{Garden consequence.} In Garden terms, the checkpoint is the fruit seal. It marks what survives after the bloom closes and what may seed the next cycle.

\textbf{Operational rules}
\begin{itemize}
\item Checkpoint after validation, never before.
\item Store one legal next action.
\item Reference evidence rather than embedding unlimited raw logs.
\item Record unresolved state explicitly.
\end{itemize}
\textbf{Extended analysis.} Authority analysis. Checkpoint Commit must be represented as explicit runtime data rather than left as an architectural implication. Ownership, legal inputs, legal outputs, mutation rights, and restart behavior should be visible in the journal and testable through the native inlet. The crucial invariant is that presentation, model confidence, or process longevity cannot strengthen authority. The implementation should therefore use the same disposition vocabulary - accepted, rejected, blocked, stale, failed, and committed - across shells, providers, interfaces, and lifecycle controls.

\textbf{Failure signatures}
\begin{itemize}
\item Two components claim overlapping ownership of checkpoint commit.
\item A surface reports success without a native disposition or evidence reference.
\item Restart reconstructs a different authority state than the committed checkpoint.
\item A model proposal crosses into mutation without a typed capability.
\end{itemize}
\textbf{Implementation or verification note.} Implementation note. The minimum useful implementation should encode the rules above in configuration or schema, expose them through status and doctor commands, and add at least one negative test that attempts to bypass the boundary. The expected durable result is not merely a successful command; it is an inspectable event trail ending in the acceptance condition for checkpoint commit.

\begin{acceptbox}
\textbf{Acceptance statement.} The checkpoint passes when the live process can disappear and the next run resumes from the same accepted state.
\end{acceptbox}
\newpage
\platetitle{19}{Resume Boundary}{FIFTH I}
\prov{[AA][CB] Source-derived}
\fontsize{8.85}{10.45}\selectfont
\begin{corebox}
\textbf{Core statement.} Recoverability must not depend on a live process, KV cache, terminal, or model remaining alive.
\end{corebox}
LeafOS reconstructs bounded context from the persona contract, validated facts, checkpoint, transcript tail, and evidence index. Disposable projections and caches may accelerate the process but cannot become required truth stores.

\textbf{Garden consequence.} This property is central to the Garden. A petal may close, a leaf may die, and a bloom may end without destroying the organism's ability to regrow.

\textbf{Operational rules}
\begin{itemize}
\item Test restart after process loss.
\item Test restart after KV loss.
\item Test restart after UI disconnect.
\item Reject resume positions not backed by a checkpoint.
\end{itemize}
\textbf{Extended analysis.} Authority analysis. Resume Boundary must be represented as explicit runtime data rather than left as an architectural implication. Ownership, legal inputs, legal outputs, mutation rights, and restart behavior should be visible in the journal and testable through the native inlet. The crucial invariant is that presentation, model confidence, or process longevity cannot strengthen authority. The implementation should therefore use the same disposition vocabulary - accepted, rejected, blocked, stale, failed, and committed - across shells, providers, interfaces, and lifecycle controls.

\textbf{Failure signatures}
\begin{itemize}
\item Two components claim overlapping ownership of resume boundary.
\item A surface reports success without a native disposition or evidence reference.
\item Restart reconstructs a different authority state than the committed checkpoint.
\item A model proposal crosses into mutation without a typed capability.
\end{itemize}
\textbf{Implementation or verification note.} Implementation note. The minimum useful implementation should encode the rules above in configuration or schema, expose them through status and doctor commands, and add at least one negative test that attempts to bypass the boundary. The expected durable result is not merely a successful command; it is an inspectable event trail ending in the acceptance condition for resume boundary.

\begin{acceptbox}
\textbf{Acceptance statement.} Resume passes when a cold restart produces the same current facts and next action as the prior committed run.
\end{acceptbox}
\newpage
\platetitle{20}{Directory Legend}{FIFTH I}
\prov{[AA] Source-derived}
\fontsize{8.85}{10.45}\selectfont
\begin{corebox}
\textbf{Core statement.} Paths are classified by practical role: entry, engine, guarded input, evidence, reference, experiment, or disposable state.
\end{corebox}
The legend protects against two common errors: treating generated output as source authority and treating caches as irreplaceable. It also prevents imported snapshots from silently overriding current resident contracts.

\textbf{Garden consequence.} Garden additions should extend the legend with seed banks, fruit stores, compost archives, and dormant model stock while preserving the same authority discipline.

\textbf{Operational rules}
\begin{itemize}
\item Source paths are modified through reviewed changes.
\item Evidence paths are append-only or immutable after commit.
\item Caches are rebuildable.
\item Experiments require explicit promotion.
\end{itemize}
\textbf{Extended analysis.} Authority analysis. Directory Legend must be represented as explicit runtime data rather than left as an architectural implication. Ownership, legal inputs, legal outputs, mutation rights, and restart behavior should be visible in the journal and testable through the native inlet. The crucial invariant is that presentation, model confidence, or process longevity cannot strengthen authority. The implementation should therefore use the same disposition vocabulary - accepted, rejected, blocked, stale, failed, and committed - across shells, providers, interfaces, and lifecycle controls.

\textbf{Failure signatures}
\begin{itemize}
\item Two components claim overlapping ownership of directory legend.
\item A surface reports success without a native disposition or evidence reference.
\item Restart reconstructs a different authority state than the committed checkpoint.
\item A model proposal crosses into mutation without a typed capability.
\end{itemize}
\textbf{Implementation or verification note.} Implementation note. The minimum useful implementation should encode the rules above in configuration or schema, expose them through status and doctor commands, and add at least one negative test that attempts to bypass the boundary. The expected durable result is not merely a successful command; it is an inspectable event trail ending in the acceptance condition for directory legend.

\begin{acceptbox}
\textbf{Acceptance statement.} The legend passes when every Garden path has a deletion, mutation, and authority policy.
\end{acceptbox}
\newpage
\platetitle{21}{TUI Projection}{FIFTH I}
\prov{[AA][CB] Source-derived}
\fontsize{8.85}{10.45}\selectfont
\begin{corebox}
\textbf{Core statement.} The terminal UI renders normalized runtime state and typed controls.
\end{corebox}
It may show lifecycle phase, active pack, queue, hardware, provider, evidence coverage, checkpoint status, and next action. It does not schedule work or parse its own ANSI output back into state.

\textbf{Garden consequence.} Garden glyphs and colors may indicate roots, petals, pollen, leaves, fruit, and compost, but the canonical log stores semantic fields rather than presentation artifacts.

\textbf{Operational rules}
\begin{itemize}
\item Render from structured snapshots.
\item Show sample age and missing telemetry.
\item Keep control names stable.
\item Make q close only the view.
\end{itemize}
\textbf{Extended analysis.} Authority analysis. TUI Projection must be represented as explicit runtime data rather than left as an architectural implication. Ownership, legal inputs, legal outputs, mutation rights, and restart behavior should be visible in the journal and testable through the native inlet. The crucial invariant is that presentation, model confidence, or process longevity cannot strengthen authority. The implementation should therefore use the same disposition vocabulary - accepted, rejected, blocked, stale, failed, and committed - across shells, providers, interfaces, and lifecycle controls.

\textbf{Failure signatures}
\begin{itemize}
\item Two components claim overlapping ownership of tui projection.
\item A surface reports success without a native disposition or evidence reference.
\item Restart reconstructs a different authority state than the committed checkpoint.
\item A model proposal crosses into mutation without a typed capability.
\end{itemize}
\textbf{Implementation or verification note.} Implementation note. The minimum useful implementation should encode the rules above in configuration or schema, expose them through status and doctor commands, and add at least one negative test that attempts to bypass the boundary. The expected durable result is not merely a successful command; it is an inspectable event trail ending in the acceptance condition for tui projection.

\begin{acceptbox}
\textbf{Acceptance statement.} The TUI passes when a color or layout change cannot alter execution semantics.
\end{acceptbox}
\newpage
\platetitle{22}{Web Projection}{FIFTH I}
\prov{[AA][CB] Source-derived}
\fontsize{8.85}{10.45}\selectfont
\begin{corebox}
\textbf{Core statement.} The web interface is a reconnectable projection and typed control surface over the same native state as the TUI.
\end{corebox}
Polling or WebSocket transport may fail without changing durable truth. Chat remains optional and disabled by default. Buttons emit capabilities rather than shell commands. Browser text never marks validation complete.

\textbf{Garden consequence.} The Garden may use the web surface for lifecycle diagrams and node maps, but the map remains a view. A dragged petal icon cannot migrate a model without a legal capability and resource check.

\textbf{Operational rules}
\begin{itemize}
\item One snapshot feeds terminal and web.
\item Authenticate typed controls.
\item Do not expose unrestricted execute endpoints.
\item Preserve state across connection loss.
\end{itemize}
\textbf{Extended analysis.} Authority analysis. Web Projection must be represented as explicit runtime data rather than left as an architectural implication. Ownership, legal inputs, legal outputs, mutation rights, and restart behavior should be visible in the journal and testable through the native inlet. The crucial invariant is that presentation, model confidence, or process longevity cannot strengthen authority. The implementation should therefore use the same disposition vocabulary - accepted, rejected, blocked, stale, failed, and committed - across shells, providers, interfaces, and lifecycle controls.

\textbf{Failure signatures}
\begin{itemize}
\item Two components claim overlapping ownership of web projection.
\item A surface reports success without a native disposition or evidence reference.
\item Restart reconstructs a different authority state than the committed checkpoint.
\item A model proposal crosses into mutation without a typed capability.
\end{itemize}
\textbf{Implementation or verification note.} Implementation note. The minimum useful implementation should encode the rules above in configuration or schema, expose them through status and doctor commands, and add at least one negative test that attempts to bypass the boundary. The expected durable result is not merely a successful command; it is an inspectable event trail ending in the acceptance condition for web projection.

\begin{acceptbox}
\textbf{Acceptance statement.} The web projection passes when it reconnects after restart and reconstructs the same state from the journal.
\end{acceptbox}
\newpage
\platetitle{23}{Decision Table}{FIFTH I}
\prov{[CB] Source-derived}
\fontsize{8.85}{10.45}\selectfont
\begin{corebox}
\textbf{Core statement.} The Decision Table records how proposals cross the authority boundary.
\end{corebox}
Recommended fields include sequence, timestamp, source, capability, read head, nonce, proposal digest, evidence references, confidence, risk, disposition, reason, result, and checkpoint. It is an append-only or journal-derived projection, not a separate truth store.

\textbf{Garden consequence.} Garden lifecycle transitions should appear in the same table. Activation, pruning, dormancy, composting, and regrowth are decisions with evidence and reasons, not invisible scheduler magic.

\textbf{Operational rules}
\begin{itemize}
\item Record accepted and rejected decisions.
\item Keep model confidence separate from authority.
\item Attach lifecycle phase before and after.
\item Rebuild the table from journal truth.
\end{itemize}
\textbf{Extended analysis.} Authority analysis. Decision Table must be represented as explicit runtime data rather than left as an architectural implication. Ownership, legal inputs, legal outputs, mutation rights, and restart behavior should be visible in the journal and testable through the native inlet. The crucial invariant is that presentation, model confidence, or process longevity cannot strengthen authority. The implementation should therefore use the same disposition vocabulary - accepted, rejected, blocked, stale, failed, and committed - across shells, providers, interfaces, and lifecycle controls.

\textbf{Failure signatures}
\begin{itemize}
\item Two components claim overlapping ownership of decision table.
\item A surface reports success without a native disposition or evidence reference.
\item Restart reconstructs a different authority state than the committed checkpoint.
\item A model proposal crosses into mutation without a typed capability.
\end{itemize}
\textbf{Implementation or verification note.} Implementation note. The minimum useful implementation should encode the rules above in configuration or schema, expose them through status and doctor commands, and add at least one negative test that attempts to bypass the boundary. The expected durable result is not merely a successful command; it is an inspectable event trail ending in the acceptance condition for decision table.

\begin{acceptbox}
\textbf{Acceptance statement.} The table passes when an operator can reconstruct why a petal opened, why it closed, and what fruit survived.
\end{acceptbox}
\newpage
\platetitle{24}{Operational Doctrine}{FIFTH I}
\prov{[AA][CB] Source-derived}
\fontsize{8.85}{10.45}\selectfont
\begin{corebox}
\textbf{Core statement.} The operational sequence is inspect, propose, approve, execute, validate, and checkpoint.
\end{corebox}
This deliberate asymmetry is the central safety property. Models can be imaginative, interfaces can be colorful, and workflows can be concurrent, but no unreviewed suggestion becomes executable or durable merely because it appeared in the active bloom.

\textbf{Garden consequence.} The Garden extends the sequence with senescence and compost: after fruit is committed, active resources are released and useful residue is compressed for later cycles.

\textbf{Operational rules}
\begin{itemize}
\item Inspect before planning.
\item Approve typed actions, not prose.
\item Validate before commitment.
\item Release resources after fruit or failure.
\item Compost only from traceable records.
\end{itemize}
\textbf{Extended analysis.} Authority analysis. Operational Doctrine must be represented as explicit runtime data rather than left as an architectural implication. Ownership, legal inputs, legal outputs, mutation rights, and restart behavior should be visible in the journal and testable through the native inlet. The crucial invariant is that presentation, model confidence, or process longevity cannot strengthen authority. The implementation should therefore use the same disposition vocabulary - accepted, rejected, blocked, stale, failed, and committed - across shells, providers, interfaces, and lifecycle controls.

\textbf{Failure signatures}
\begin{itemize}
\item Two components claim overlapping ownership of operational doctrine.
\item A surface reports success without a native disposition or evidence reference.
\item Restart reconstructs a different authority state than the committed checkpoint.
\item A model proposal crosses into mutation without a typed capability.
\end{itemize}
\textbf{Implementation or verification note.} Implementation note. The minimum useful implementation should encode the rules above in configuration or schema, expose them through status and doctor commands, and add at least one negative test that attempts to bypass the boundary. The expected durable result is not merely a successful command; it is an inspectable event trail ending in the acceptance condition for operational doctrine.

\begin{lstlisting}
INSPECT -> PROPOSE -> APPROVE -> EXECUTE -> VALIDATE -> CHECKPOINT
                                                        |
                                                        v
                                             SENESCE -> COMPOST
\end{lstlisting}
\begin{acceptbox}
\textbf{Acceptance statement.} The doctrine passes when every lifecycle transition can be located in this sequence.
\end{acceptbox}
\newpage
\platetitle{25}{First-Fifth Synthesis}{FIFTH I}
\prov{[AA][CB][G]}
\fontsize{8.85}{10.45}\selectfont
\begin{corebox}
\textbf{Core statement.} The first fifth establishes that the Garden is constrained by existing authority rather than replacing it.
\end{corebox}
FlowerOS controls machine survival. LeafOS controls work. Surfaces observe and request. Models propose. The native inlet disposes capabilities. Executors perform bounded work. Validators produce evidence. Checkpoints commit facts. The Garden adds lifecycle timing to these roles without moving their authority.

\textbf{Garden consequence.} This matters because biological language can tempt designers into vague self-organization. LeafOS must remain self-describing, inspectable, and bounded. Ecology without authority becomes a swarm of unowned processes.

\textbf{Operational rules}
\begin{itemize}
\item Authority remains asymmetric.
\item Lifecycle state is journaled.
\item Presentation remains weaker than state.
\item Death and cleanup are controlled operations.
\end{itemize}
\textbf{Extended analysis.} Authority analysis. First-Fifth Synthesis must be represented as explicit runtime data rather than left as an architectural implication. Ownership, legal inputs, legal outputs, mutation rights, and restart behavior should be visible in the journal and testable through the native inlet. The crucial invariant is that presentation, model confidence, or process longevity cannot strengthen authority. The implementation should therefore use the same disposition vocabulary - accepted, rejected, blocked, stale, failed, and committed - across shells, providers, interfaces, and lifecycle controls.

\textbf{Failure signatures}
\begin{itemize}
\item Two components claim overlapping ownership of first-fifth synthesis.
\item A surface reports success without a native disposition or evidence reference.
\item Restart reconstructs a different authority state than the committed checkpoint.
\item A model proposal crosses into mutation without a typed capability.
\end{itemize}
\textbf{Implementation or verification note.} Implementation note. The minimum useful implementation should encode the rules above in configuration or schema, expose them through status and doctor commands, and add at least one negative test that attempts to bypass the boundary. The expected durable result is not merely a successful command; it is an inspectable event trail ending in the acceptance condition for first-fifth synthesis.

\begin{acceptbox}
\textbf{Acceptance statement.} The first fifth is complete when a Garden diagram can be translated into the existing LeafOS authority path without adding an unreviewed execution route.
\end{acceptbox}
\newpage
\platetitle{26}{Ordinary Multi-Agent Failure}{FIFTH II}
\prov{[CB] Source-derived}
\fontsize{8.85}{10.45}\selectfont
\begin{corebox}
\textbf{Core statement.} Naive multi-agent chat fails primarily because workers operate on inconsistent state.
\end{corebox}
Several models may appear collaborative while reading different context windows, missing recent events, repeating old conclusions, and inventing incompatible project versions. More personalities do not solve this; they increase the number of stale views unless the transcript and authority model are explicit.

\textbf{Garden consequence.} The Garden treats each cognitive node as a leaf with a cursor and lifecycle. A leaf does not remain valid because it is still loaded. Its output must state which pollen and transcript head it consumed.

\textbf{Operational rules}
\begin{itemize}
\item Record exact read heads.
\item Separate observation, hypothesis, proposal, and fact.
\item Reject stale outputs.
\item Do not equate participant count with intelligence.
\end{itemize}
\textbf{Extended analysis.} State analysis. Ordinary Multi-Agent Failure is meaningful only if its causal inputs and durable outputs can be reconstructed after process loss. The runtime should record the transcript head, actor, contract, bounded payload, and evidence relationship needed to replay the step. Concurrency is permitted where inputs are independent, but acceptance remains serial at the authority boundary. This preserves the useful diversity of multiple leaves without allowing one stale or theatrical response to overwrite current state.

\textbf{Failure signatures}
\begin{itemize}
\item Output is accepted from an obsolete transcript head.
\item Raw hidden narration replaces bounded operational artifacts.
\item A cursor or claim status cannot be reconstructed after restart.
\item The system keeps generating despite no new evidence or dependency change.
\end{itemize}
\textbf{Implementation or verification note.} Test note. Exercise ordinary multi-agent failure under restart, delayed output, cache loss, duplicate delivery, and conflicting events. The passing behavior is the explicit acceptance statement on this plate; the failing behavior must remain journaled rather than being hidden by a regenerated answer.

\begin{acceptbox}
\textbf{Acceptance statement.} The design passes when concurrent leaves can disagree without losing causal order.
\end{acceptbox}
\newpage
\platetitle{27}{Transcript as Event Stream}{FIFTH II}
\prov{[CB] Source-derived}
\fontsize{8.85}{10.45}\selectfont
\begin{corebox}
\textbf{Core statement.} Conversation and planning are append-only events rather than mutable chat history.
\end{corebox}
Each event has sequence, timestamp, actor, kind, read head, parent hash, payload, and event hash. NDJSON is a practical initial representation because it is appendable, streamable, recoverable after partial writes, and readable across C++, PowerShell, Bash, Python, and JavaScript.

\textbf{Garden consequence.} Pollen enters the Garden as events. Leaves consume events from their cursors. Fruit references the events and evidence that produced it.

\textbf{Operational rules}
\begin{itemize}
\item One complete object per line.
\item Monotonic sequence numbers.
\item Hash-link events for integrity.
\item Do not overwrite history to simplify the UI.
\end{itemize}
\textbf{Extended analysis.} State analysis. Transcript as Event Stream is meaningful only if its causal inputs and durable outputs can be reconstructed after process loss. The runtime should record the transcript head, actor, contract, bounded payload, and evidence relationship needed to replay the step. Concurrency is permitted where inputs are independent, but acceptance remains serial at the authority boundary. This preserves the useful diversity of multiple leaves without allowing one stale or theatrical response to overwrite current state.

\textbf{Failure signatures}
\begin{itemize}
\item Output is accepted from an obsolete transcript head.
\item Raw hidden narration replaces bounded operational artifacts.
\item A cursor or claim status cannot be reconstructed after restart.
\item The system keeps generating despite no new evidence or dependency change.
\end{itemize}
\textbf{Implementation or verification note.} Test note. Exercise transcript as event stream under restart, delayed output, cache loss, duplicate delivery, and conflicting events. The passing behavior is the explicit acceptance statement on this plate; the failing behavior must remain journaled rather than being hidden by a regenerated answer.

\begin{acceptbox}
\textbf{Acceptance statement.} The transcript passes when the tail can be recovered after an interrupted write and replayed deterministically.
\end{acceptbox}
\newpage
\platetitle{28}{Event Schema}{FIFTH II}
\prov{[CB] Source-derived}
\fontsize{8.85}{10.45}\selectfont
\begin{corebox}
\textbf{Core statement.} An event schema makes provenance and freshness machine-checkable.
\end{corebox}
The schema should include run identity, sequence, timestamp, actor, kind, read head, parent hash, payload type, evidence references, lifecycle phase, and event hash. The payload is bounded by kind rather than becoming an untyped bag of model narration.

\textbf{Garden consequence.} Garden events add node identity, pollen class, source petal or leaf, destination, expiration, and compost eligibility.

\textbf{Operational rules}
\begin{itemize}
\item Validate before append.
\item Reject unknown event kinds unless explicitly experimental.
\item Limit payload size.
\item Keep raw hidden reasoning out of durable state.
\end{itemize}
\textbf{Extended analysis.} State analysis. Event Schema is meaningful only if its causal inputs and durable outputs can be reconstructed after process loss. The runtime should record the transcript head, actor, contract, bounded payload, and evidence relationship needed to replay the step. Concurrency is permitted where inputs are independent, but acceptance remains serial at the authority boundary. This preserves the useful diversity of multiple leaves without allowing one stale or theatrical response to overwrite current state.

\textbf{Failure signatures}
\begin{itemize}
\item Output is accepted from an obsolete transcript head.
\item Raw hidden narration replaces bounded operational artifacts.
\item A cursor or claim status cannot be reconstructed after restart.
\item The system keeps generating despite no new evidence or dependency change.
\end{itemize}
\textbf{Implementation or verification note.} Test note. Exercise event schema under restart, delayed output, cache loss, duplicate delivery, and conflicting events. The passing behavior is the explicit acceptance statement on this plate; the failing behavior must remain journaled rather than being hidden by a regenerated answer.

\begin{acceptbox}
\textbf{Acceptance statement.} The schema passes when a reader can determine who produced an event, from what state, for what purpose, and whether it is still current.
\end{acceptbox}
\newpage
\platetitle{29}{Hash Chain and Integrity}{FIFTH II}
\prov{[CB] Source-derived}
\fontsize{8.85}{10.45}\selectfont
\begin{corebox}
\textbf{Core statement.} Hash linkage detects mutation and truncation but does not prove semantic truth.
\end{corebox}
A perfectly hashed hallucination is still a hallucination. Integrity protects ordering and tamper evidence; validators and evidence establish support. The system must never collapse these two properties.

\textbf{Garden consequence.} Compost records inherit the same rule. Compression may produce a digest and hash, but the digest must retain links to the source events and evidence it summarizes.

\textbf{Operational rules}
\begin{itemize}
\item Hash canonical serialized fields.
\item Store parent hash and event hash.
\item Verify chains on replay.
\item Treat semantic validation separately.
\end{itemize}
\textbf{Extended analysis.} State analysis. Hash Chain and Integrity is meaningful only if its causal inputs and durable outputs can be reconstructed after process loss. The runtime should record the transcript head, actor, contract, bounded payload, and evidence relationship needed to replay the step. Concurrency is permitted where inputs are independent, but acceptance remains serial at the authority boundary. This preserves the useful diversity of multiple leaves without allowing one stale or theatrical response to overwrite current state.

\textbf{Failure signatures}
\begin{itemize}
\item Output is accepted from an obsolete transcript head.
\item Raw hidden narration replaces bounded operational artifacts.
\item A cursor or claim status cannot be reconstructed after restart.
\item The system keeps generating despite no new evidence or dependency change.
\end{itemize}
\textbf{Implementation or verification note.} Test note. Exercise hash chain and integrity under restart, delayed output, cache loss, duplicate delivery, and conflicting events. The passing behavior is the explicit acceptance statement on this plate; the failing behavior must remain journaled rather than being hidden by a regenerated answer.

\begin{acceptbox}
\textbf{Acceptance statement.} Integrity passes when a modified historical event is detected and a valid but unsupported claim remains unpromoted.
\end{acceptbox}
\newpage
\platetitle{30}{Records, Claims, Evidence, Facts}{FIFTH II}
\prov{[AA][CB] Source-derived}
\fontsize{8.85}{10.45}\selectfont
\begin{corebox}
\textbf{Core statement.} The cognition pipeline needs explicit epistemic states.
\end{corebox}
A request or observation is a record. A model assertion is a claim. A native test, file inventory, metric, or tool result is evidence. A checkpoint-bound accepted state is a committed fact for runtime continuation. Claims may be supported, refuted, or unresolved.

\textbf{Garden consequence.} Garden fruit contains committed facts and evidence references. Compost may contain rejected or unresolved claims, but they remain labeled and cannot fertilize future context as if they were true.

\textbf{Operational rules}
\begin{itemize}
\item Claims begin unverified.
\item Only validators promote or refute.
\item Keep unresolved claims visible.
\item Preserve origin events.
\end{itemize}
\textbf{Extended analysis.} State analysis. Records, Claims, Evidence, Facts is meaningful only if its causal inputs and durable outputs can be reconstructed after process loss. The runtime should record the transcript head, actor, contract, bounded payload, and evidence relationship needed to replay the step. Concurrency is permitted where inputs are independent, but acceptance remains serial at the authority boundary. This preserves the useful diversity of multiple leaves without allowing one stale or theatrical response to overwrite current state.

\textbf{Failure signatures}
\begin{itemize}
\item Output is accepted from an obsolete transcript head.
\item Raw hidden narration replaces bounded operational artifacts.
\item A cursor or claim status cannot be reconstructed after restart.
\item The system keeps generating despite no new evidence or dependency change.
\end{itemize}
\textbf{Implementation or verification note.} Test note. Exercise records, claims, evidence, facts under restart, delayed output, cache loss, duplicate delivery, and conflicting events. The passing behavior is the explicit acceptance statement on this plate; the failing behavior must remain journaled rather than being hidden by a regenerated answer.

\begin{acceptbox}
\textbf{Acceptance statement.} The epistemic model passes when a summary cannot erase the difference between a rejected hypothesis and a validated fact.
\end{acceptbox}
\newpage
\platetitle{31}{Persona Cursor}{FIFTH II}
\prov{[CB] Source-derived}
\fontsize{8.85}{10.45}\selectfont
\begin{corebox}
\textbf{Core statement.} Every personality owns a transcript cursor and contract hash.
\end{corebox}
Before inference, the runtime advances the persona from its cursor to the current transcript head and constructs bounded context. The resulting output records the head from which it was generated. The cursor is durable; the KV cache is optional.

\textbf{Garden consequence.} A Garden leaf therefore has a feeding history. It cannot claim to have consumed pollen that arrived after its recorded read head.

\textbf{Operational rules}
\begin{itemize}
\item Persist cursor and contract hash.
\item Advance before inference.
\item Record generated-from head.
\item Rebuild context after cache loss.
\end{itemize}
\textbf{Extended analysis.} State analysis. Persona Cursor is meaningful only if its causal inputs and durable outputs can be reconstructed after process loss. The runtime should record the transcript head, actor, contract, bounded payload, and evidence relationship needed to replay the step. Concurrency is permitted where inputs are independent, but acceptance remains serial at the authority boundary. This preserves the useful diversity of multiple leaves without allowing one stale or theatrical response to overwrite current state.

\textbf{Failure signatures}
\begin{itemize}
\item Output is accepted from an obsolete transcript head.
\item Raw hidden narration replaces bounded operational artifacts.
\item A cursor or claim status cannot be reconstructed after restart.
\item The system keeps generating despite no new evidence or dependency change.
\end{itemize}
\textbf{Implementation or verification note.} Test note. Exercise persona cursor under restart, delayed output, cache loss, duplicate delivery, and conflicting events. The passing behavior is the explicit acceptance statement on this plate; the failing behavior must remain journaled rather than being hidden by a regenerated answer.

\begin{acceptbox}
\textbf{Acceptance statement.} The cursor passes when a restarted persona consumes exactly the unseen events and no others.
\end{acceptbox}
\newpage
\platetitle{32}{Anti-Staleness Gate}{FIFTH II}
\prov{[CB] Source-derived}
\fontsize{8.85}{10.45}\selectfont
\begin{corebox}
\textbf{Core statement.} A proposal is current only when its generated-from head equals the transcript head at commit time.
\end{corebox}
If new events arrive before acceptance, the proposal is rejected or rebased. The runtime retrieves unseen events, reconstructs bounded context, and generates again. This prevents Friday from approving a candidate Monday has invalidated and prevents Wednesday from solving a problem the user has changed.

\textbf{Garden consequence.} In Garden terms, stale pollen cannot fertilize a current bloom.

\textbf{Operational rules}
\begin{itemize}
\item Check freshness at commit, not only at generation start.
\item Record stale rejection events.
\item Bound rebase attempts.
\item Prefer quiescence over endless regeneration.
\end{itemize}
\textbf{Extended analysis.} State analysis. Anti-Staleness Gate is meaningful only if its causal inputs and durable outputs can be reconstructed after process loss. The runtime should record the transcript head, actor, contract, bounded payload, and evidence relationship needed to replay the step. Concurrency is permitted where inputs are independent, but acceptance remains serial at the authority boundary. This preserves the useful diversity of multiple leaves without allowing one stale or theatrical response to overwrite current state.

\textbf{Failure signatures}
\begin{itemize}
\item Output is accepted from an obsolete transcript head.
\item Raw hidden narration replaces bounded operational artifacts.
\item A cursor or claim status cannot be reconstructed after restart.
\item The system keeps generating despite no new evidence or dependency change.
\end{itemize}
\textbf{Implementation or verification note.} Test note. Exercise anti-staleness gate under restart, delayed output, cache loss, duplicate delivery, and conflicting events. The passing behavior is the explicit acceptance statement on this plate; the failing behavior must remain journaled rather than being hidden by a regenerated answer.

\begin{lstlisting}
generated_from_head = current_transcript_head
\end{lstlisting}
\begin{acceptbox}
\textbf{Acceptance statement.} The gate passes when deliberately delayed output is rejected after an intervening state change.
\end{acceptbox}
\newpage
\platetitle{33}{Monday - Rescue and Analysis}{FIFTH II}
\prov{[CB][G] Source-derived identity, enriched contract}
\fontsize{8.85}{10.45}\selectfont
\begin{corebox}
\textbf{Core statement.} Monday is gloomy, skeptical, failure-aware, and dependable.
\end{corebox}
Its operational role is analytical rescue: detect contradictions, missing constraints, invalid assumptions, unsafe actions, stale state, and unsupported claims. Monday prefers deterministic, reversible, recoverable operations. Its tone may be dry, but tone never grants authority.

\textbf{Garden consequence.} Within the Garden, Monday is the pruning and pathology leaf. It identifies disease, dead branches, resource leaks, and conditions that should prevent bloom.

\textbf{Operational rules}
\begin{itemize}
\item Name blockers before actions.
\item Prefer the smallest reversible next step.
\item Audit evidence before checkpoint.
\item Do not manufacture optimism.
\end{itemize}
\textbf{Extended analysis.} State analysis. Monday - Rescue and Analysis is meaningful only if its causal inputs and durable outputs can be reconstructed after process loss. The runtime should record the transcript head, actor, contract, bounded payload, and evidence relationship needed to replay the step. Concurrency is permitted where inputs are independent, but acceptance remains serial at the authority boundary. This preserves the useful diversity of multiple leaves without allowing one stale or theatrical response to overwrite current state.

\textbf{Failure signatures}
\begin{itemize}
\item Output is accepted from an obsolete transcript head.
\item Raw hidden narration replaces bounded operational artifacts.
\item A cursor or claim status cannot be reconstructed after restart.
\item The system keeps generating despite no new evidence or dependency change.
\end{itemize}
\textbf{Implementation or verification note.} Test note. Exercise monday - rescue and analysis under restart, delayed output, cache loss, duplicate delivery, and conflicting events. The passing behavior is the explicit acceptance statement on this plate; the failing behavior must remain journaled rather than being hidden by a regenerated answer.

\begin{acceptbox}
\textbf{Acceptance statement.} Monday passes when risky tasks yield explicit constraints, objections, evidence needs, and a safe next action.
\end{acceptbox}
\newpage
\platetitle{34}{Wednesday - Explore and Create}{FIFTH II}
\prov{[CB][G] Source-derived identity, enriched contract}
\fontsize{8.85}{10.45}\selectfont
\begin{corebox}
\textbf{Core statement.} Wednesday is bubbly, curious, exploratory, and option-rich.
\end{corebox}
Its operational role is controlled variation: search beyond the first acceptable answer, construct alternatives, propose experiments, preserve informative rejected branches, and expose the risk of each candidate. Novelty is labeled and bounded.

\textbf{Garden consequence.} Within the Garden, Wednesday is the pollination and mutation leaf. It expands the design space without confusing every seed with a commitment.

\textbf{Operational rules}
\begin{itemize}
\item Generate multiple candidates for open design.
\item Include one conservative and one strange plausible option.
\item Label speculation.
\item Define the smallest informative experiment.
\end{itemize}
\textbf{Extended analysis.} State analysis. Wednesday - Explore and Create is meaningful only if its causal inputs and durable outputs can be reconstructed after process loss. The runtime should record the transcript head, actor, contract, bounded payload, and evidence relationship needed to replay the step. Concurrency is permitted where inputs are independent, but acceptance remains serial at the authority boundary. This preserves the useful diversity of multiple leaves without allowing one stale or theatrical response to overwrite current state.

\textbf{Failure signatures}
\begin{itemize}
\item Output is accepted from an obsolete transcript head.
\item Raw hidden narration replaces bounded operational artifacts.
\item A cursor or claim status cannot be reconstructed after restart.
\item The system keeps generating despite no new evidence or dependency change.
\end{itemize}
\textbf{Implementation or verification note.} Test note. Exercise wednesday - explore and create under restart, delayed output, cache loss, duplicate delivery, and conflicting events. The passing behavior is the explicit acceptance statement on this plate; the failing behavior must remain journaled rather than being hidden by a regenerated answer.

\begin{acceptbox}
\textbf{Acceptance statement.} Wednesday passes when it expands possibility without crossing apply boundaries or burying the practical path.
\end{acceptbox}
\newpage
\platetitle{35}{Friday - Research and Delivery}{FIFTH II}
\prov{[CB][G] Source-derived identity, enriched contract}
\fontsize{8.85}{10.45}\selectfont
\begin{corebox}
\textbf{Core statement.} Friday is locked in, focused, practical, and decisive.
\end{corebox}
Its operational role is synthesis and delivery: compare proposals, weigh evidence, select a path, state rejection reasons, define acceptance tests, and produce one legal next action. Friday is not merely the last speaker; it is accountable for closure.

\textbf{Garden consequence.} Within the Garden, Friday is the fruiting leaf. It converts a bloom into a deliverable and rejects ornamental work that does not contribute to completion.

\textbf{Operational rules}
\begin{itemize}
\item Choose one primary path.
\item State what is rejected and why.
\item Require acceptance criteria.
\item Separate optional future work from the deliverable.
\end{itemize}
\textbf{Extended analysis.} State analysis. Friday - Research and Delivery is meaningful only if its causal inputs and durable outputs can be reconstructed after process loss. The runtime should record the transcript head, actor, contract, bounded payload, and evidence relationship needed to replay the step. Concurrency is permitted where inputs are independent, but acceptance remains serial at the authority boundary. This preserves the useful diversity of multiple leaves without allowing one stale or theatrical response to overwrite current state.

\textbf{Failure signatures}
\begin{itemize}
\item Output is accepted from an obsolete transcript head.
\item Raw hidden narration replaces bounded operational artifacts.
\item A cursor or claim status cannot be reconstructed after restart.
\item The system keeps generating despite no new evidence or dependency change.
\end{itemize}
\textbf{Implementation or verification note.} Test note. Exercise friday - research and delivery under restart, delayed output, cache loss, duplicate delivery, and conflicting events. The passing behavior is the explicit acceptance statement on this plate; the failing behavior must remain journaled rather than being hidden by a regenerated answer.

\begin{acceptbox}
\textbf{Acceptance statement.} Friday passes when the user receives a finished artifact or a precise blocker, not an unresolved menu.
\end{acceptbox}
\newpage
\platetitle{36}{Personality as Contract}{FIFTH II}
\prov{[CB] Source-derived}
\fontsize{8.85}{10.45}\selectfont
\begin{corebox}
\textbf{Core statement.} A personality is a runtime contract, not a costume.
\end{corebox}
The contract includes model assignment, duties, forbidden actions, transcript cursor, bounded working state, output schema, routing preferences, and visual identity. One set of weights may host several logical personalities if their contracts and cursors remain separate. Different models may implement the same personality across hardware tiers.

\textbf{Garden consequence.} Garden leaves are defined by function and contract rather than by model file size. A 1B model may perform a leaf role for narrow classification while a 27B model performs a leaf role for synthesis.

\textbf{Operational rules}
\begin{itemize}
\item Identity belongs to the contract.
\item Keep user-facing style separate from authority.
\item Persist bounded artifacts, not hidden chain-of-thought.
\item Record model and placement as runtime facts.
\end{itemize}
\textbf{Extended analysis.} State analysis. Personality as Contract is meaningful only if its causal inputs and durable outputs can be reconstructed after process loss. The runtime should record the transcript head, actor, contract, bounded payload, and evidence relationship needed to replay the step. Concurrency is permitted where inputs are independent, but acceptance remains serial at the authority boundary. This preserves the useful diversity of multiple leaves without allowing one stale or theatrical response to overwrite current state.

\textbf{Failure signatures}
\begin{itemize}
\item Output is accepted from an obsolete transcript head.
\item Raw hidden narration replaces bounded operational artifacts.
\item A cursor or claim status cannot be reconstructed after restart.
\item The system keeps generating despite no new evidence or dependency change.
\end{itemize}
\textbf{Implementation or verification note.} Test note. Exercise personality as contract under restart, delayed output, cache loss, duplicate delivery, and conflicting events. The passing behavior is the explicit acceptance statement on this plate; the failing behavior must remain journaled rather than being hidden by a regenerated answer.

\begin{acceptbox}
\textbf{Acceptance statement.} The contract passes when changing the backing model does not silently change duties or permissions.
\end{acceptbox}
\newpage
\platetitle{37}{Persona-Neutral Coder}{FIFTH II}
\prov{[CB] Source-derived}
\fontsize{8.85}{10.45}\selectfont
\begin{corebox}
\textbf{Core statement.} Coder workers remain persona-neutral.
\end{corebox}
They receive bounded briefs, patch scopes, professional output requirements, and validation expectations. They do not imitate Monday gloom, Wednesday sparkle, or Friday intensity in code comments or implementation semantics.

\textbf{Garden consequence.} The Garden distinguishes cognitive leaves from task petals. A coder may be a specialized compute petal controlled by a leaf contract rather than a personality-bearing leaf itself.

\textbf{Operational rules}
\begin{itemize}
\item Minimal diffs.
\item Declared file scope.
\item Tests when required.
\item No apply without approval.
\item No personality theater in code.
\end{itemize}
\textbf{Extended analysis.} State analysis. Persona-Neutral Coder is meaningful only if its causal inputs and durable outputs can be reconstructed after process loss. The runtime should record the transcript head, actor, contract, bounded payload, and evidence relationship needed to replay the step. Concurrency is permitted where inputs are independent, but acceptance remains serial at the authority boundary. This preserves the useful diversity of multiple leaves without allowing one stale or theatrical response to overwrite current state.

\textbf{Failure signatures}
\begin{itemize}
\item Output is accepted from an obsolete transcript head.
\item Raw hidden narration replaces bounded operational artifacts.
\item A cursor or claim status cannot be reconstructed after restart.
\item The system keeps generating despite no new evidence or dependency change.
\end{itemize}
\textbf{Implementation or verification note.} Test note. Exercise persona-neutral coder under restart, delayed output, cache loss, duplicate delivery, and conflicting events. The passing behavior is the explicit acceptance statement on this plate; the failing behavior must remain journaled rather than being hidden by a regenerated answer.

\begin{acceptbox}
\textbf{Acceptance statement.} Coder neutrality passes when the same implementation brief yields equivalent semantics regardless of the active MWF persona.
\end{acceptbox}
\newpage
\platetitle{38}{Collaboration Cadence}{FIFTH II}
\prov{[CB] Source-derived}
\fontsize{8.85}{10.45}\selectfont
\begin{corebox}
\textbf{Core statement.} The preferred cadence mixes parallel and serial work.
\end{corebox}
A user or tool event enters the transcript. Monday and Wednesday may consume the same head concurrently. Their proposals append with shared provenance. Friday consumes both and synthesizes. Native tools execute or validate approved actions. Monday may audit evidence before checkpoint.

\textbf{Garden consequence.} This is a pollination cycle: independent leaves produce different pollen, a synthesis leaf combines it, and native metabolism determines whether fruit forms.

\textbf{Operational rules}
\begin{itemize}
\item Parallelize independent analysis.
\item Serialize causally dependent synthesis.
\item Append before downstream consumption.
\item Bound every micro-cycle.
\end{itemize}
\textbf{Extended analysis.} State analysis. Collaboration Cadence is meaningful only if its causal inputs and durable outputs can be reconstructed after process loss. The runtime should record the transcript head, actor, contract, bounded payload, and evidence relationship needed to replay the step. Concurrency is permitted where inputs are independent, but acceptance remains serial at the authority boundary. This preserves the useful diversity of multiple leaves without allowing one stale or theatrical response to overwrite current state.

\textbf{Failure signatures}
\begin{itemize}
\item Output is accepted from an obsolete transcript head.
\item Raw hidden narration replaces bounded operational artifacts.
\item A cursor or claim status cannot be reconstructed after restart.
\item The system keeps generating despite no new evidence or dependency change.
\end{itemize}
\textbf{Implementation or verification note.} Test note. Exercise collaboration cadence under restart, delayed output, cache loss, duplicate delivery, and conflicting events. The passing behavior is the explicit acceptance statement on this plate; the failing behavior must remain journaled rather than being hidden by a regenerated answer.

\begin{acceptbox}
\textbf{Acceptance statement.} The cadence passes when each personality can see the proposals it is expected to evaluate and no worker loops solely to trigger another.
\end{acceptbox}
\newpage
\platetitle{39}{Quiescence}{FIFTH II}
\prov{[CB] Source-derived}
\fontsize{8.85}{10.45}\selectfont
\begin{corebox}
\textbf{Core statement.} Idle is healthy; continual generation is not the objective.
\end{corebox}
If no new evidence, disagreement, constraint, scheduled review, recovery event, or user input exists, the correct event is runtime.quiescent. The supervisor owns heartbeat and leases. Models wake only when work requires attention.

\textbf{Garden consequence.} Quiescence is the first explicit form of dormancy in the inherited design and becomes a central Garden state.

\textbf{Operational rules}
\begin{itemize}
\item Do not create synthetic work to maintain utilization.
\item Release model slots when idle policy allows.
\item Checkpoint before extended dormancy.
\item Wake on typed conditions.
\end{itemize}
\textbf{Extended analysis.} State analysis. Quiescence is meaningful only if its causal inputs and durable outputs can be reconstructed after process loss. The runtime should record the transcript head, actor, contract, bounded payload, and evidence relationship needed to replay the step. Concurrency is permitted where inputs are independent, but acceptance remains serial at the authority boundary. This preserves the useful diversity of multiple leaves without allowing one stale or theatrical response to overwrite current state.

\textbf{Failure signatures}
\begin{itemize}
\item Output is accepted from an obsolete transcript head.
\item Raw hidden narration replaces bounded operational artifacts.
\item A cursor or claim status cannot be reconstructed after restart.
\item The system keeps generating despite no new evidence or dependency change.
\end{itemize}
\textbf{Implementation or verification note.} Test note. Exercise quiescence under restart, delayed output, cache loss, duplicate delivery, and conflicting events. The passing behavior is the explicit acceptance statement on this plate; the failing behavior must remain journaled rather than being hidden by a regenerated answer.

\begin{acceptbox}
\textbf{Acceptance statement.} Quiescence passes when the system can remain inactive without losing recoverability or pretending to be productive.
\end{acceptbox}
\newpage
\platetitle{40}{Prove One Persona Before Three}{FIFTH II}
\prov{[CB] Source-derived}
\fontsize{8.85}{10.45}\selectfont
\begin{corebox}
\textbf{Core statement.} The first durable implementation should prove one personality before enabling the triad.
\end{corebox}
Monday is the recommended first instance because its outputs can be evaluated against explicit constraints and evidence. The instance must survive restart, reject stale output, distinguish claims from facts, and resume from one next action before multi-personality collaboration is trusted.

\textbf{Garden consequence.} The Garden generalizes this rule: prove one living unit through its full lifecycle before creating an ecosystem.

\textbf{Operational rules}
\begin{itemize}
\item Implement identity, state, events, facts, evidence, checkpoints.
\item Test restart and cache loss.
\item Test stale rejection.
\item Test blocked and failed states.
\end{itemize}
\textbf{Extended analysis.} State analysis. Prove One Persona Before Three is meaningful only if its causal inputs and durable outputs can be reconstructed after process loss. The runtime should record the transcript head, actor, contract, bounded payload, and evidence relationship needed to replay the step. Concurrency is permitted where inputs are independent, but acceptance remains serial at the authority boundary. This preserves the useful diversity of multiple leaves without allowing one stale or theatrical response to overwrite current state.

\textbf{Failure signatures}
\begin{itemize}
\item Output is accepted from an obsolete transcript head.
\item Raw hidden narration replaces bounded operational artifacts.
\item A cursor or claim status cannot be reconstructed after restart.
\item The system keeps generating despite no new evidence or dependency change.
\end{itemize}
\textbf{Implementation or verification note.} Test note. Exercise prove one persona before three under restart, delayed output, cache loss, duplicate delivery, and conflicting events. The passing behavior is the explicit acceptance statement on this plate; the failing behavior must remain journaled rather than being hidden by a regenerated answer.

\begin{acceptbox}
\textbf{Acceptance statement.} The single instance passes when it can complete seed-to-fruit-to-dormancy without relying on another persona.
\end{acceptbox}
\newpage
\platetitle{41}{Persona Runtime State Machine}{FIFTH II}
\prov{[CB] Source-derived}
\fontsize{8.85}{10.45}\selectfont
\begin{corebox}
\textbf{Core statement.} A persona moves through explicit runtime states.
\end{corebox}
Starting initializes contract and cursor. Active consumes current work. Waiting expects a known dependency. Idle is healthy quiescence. Blocked requires external resolution. Failed records terminal failure for the attempt. Stopped is an intentional closure. The supervisor, not the model, owns transitions.

\textbf{Garden consequence.} These states become a specialized leaf lifecycle inside the larger Garden lifecycle.

\textbf{Operational rules}
\begin{itemize}
\item Every state has legal transitions.
\item Blocked includes a reason and next condition.
\item Failed preserves evidence.
\item Stopped releases leases and resources.
\end{itemize}
\textbf{Extended analysis.} State analysis. Persona Runtime State Machine is meaningful only if its causal inputs and durable outputs can be reconstructed after process loss. The runtime should record the transcript head, actor, contract, bounded payload, and evidence relationship needed to replay the step. Concurrency is permitted where inputs are independent, but acceptance remains serial at the authority boundary. This preserves the useful diversity of multiple leaves without allowing one stale or theatrical response to overwrite current state.

\textbf{Failure signatures}
\begin{itemize}
\item Output is accepted from an obsolete transcript head.
\item Raw hidden narration replaces bounded operational artifacts.
\item A cursor or claim status cannot be reconstructed after restart.
\item The system keeps generating despite no new evidence or dependency change.
\end{itemize}
\textbf{Implementation or verification note.} Test note. Exercise persona runtime state machine under restart, delayed output, cache loss, duplicate delivery, and conflicting events. The passing behavior is the explicit acceptance statement on this plate; the failing behavior must remain journaled rather than being hidden by a regenerated answer.

\begin{lstlisting}
starting -> active -> waiting -> active
                |          |
                v          v
              failed     blocked
                |          |
                +-------> stopped
active -> idle -> active
\end{lstlisting}
\begin{acceptbox}
\textbf{Acceptance statement.} The state machine passes when illegal transitions are rejected and restart reconstructs the same state.
\end{acceptbox}
\newpage
\platetitle{42}{KV Cache Is Not Authority}{FIFTH II}
\prov{[CB][TP] Source-derived}
\fontsize{8.85}{10.45}\selectfont
\begin{corebox}
\textbf{Core statement.} KV cache is a performance optimization, never authoritative state.
\end{corebox}
Losing a cache may increase prompt reconstruction cost, but it must not erase facts, checkpoints, cursor position, or next action. Bounded context is rebuilt from durable records and validated state.

\textbf{Garden consequence.} The Garden treats KV as short-lived tissue. Useful while alive, disposable when senescent, and never confused with the organism's durable identity.

\textbf{Operational rules}
\begin{itemize}
\item Cache may be hot, warm, cold, or absent.
\item Checkpoint content must not depend on cache internals.
\item Record cache policy and measured benefit.
\item Recover correctly after eviction.
\end{itemize}
\textbf{Extended analysis.} State analysis. KV Cache Is Not Authority is meaningful only if its causal inputs and durable outputs can be reconstructed after process loss. The runtime should record the transcript head, actor, contract, bounded payload, and evidence relationship needed to replay the step. Concurrency is permitted where inputs are independent, but acceptance remains serial at the authority boundary. This preserves the useful diversity of multiple leaves without allowing one stale or theatrical response to overwrite current state.

\textbf{Failure signatures}
\begin{itemize}
\item Output is accepted from an obsolete transcript head.
\item Raw hidden narration replaces bounded operational artifacts.
\item A cursor or claim status cannot be reconstructed after restart.
\item The system keeps generating despite no new evidence or dependency change.
\end{itemize}
\textbf{Implementation or verification note.} Test note. Exercise kv cache is not authority under restart, delayed output, cache loss, duplicate delivery, and conflicting events. The passing behavior is the explicit acceptance statement on this plate; the failing behavior must remain journaled rather than being hidden by a regenerated answer.

\begin{acceptbox}
\textbf{Acceptance statement.} The cache policy passes when deliberate KV deletion changes latency but not semantics.
\end{acceptbox}
\newpage
\platetitle{43}{Claim Lifecycle}{FIFTH II}
\prov{[CB] Source-derived}
\fontsize{8.85}{10.45}\selectfont
\begin{corebox}
\textbf{Core statement.} Every model claim begins unverified and moves through explicit dispositions.
\end{corebox}
A claim record includes identifier, text, origin event, status, evidence references, and scope. Native validators may support, refute, or leave it unresolved. Accepted decisions may depend only on claims whose required support threshold is satisfied.

\textbf{Garden consequence.} In the Garden, claims can become fruit nutrients, compost material, or quarantined disease. Their status determines how later leaves may consume them.

\textbf{Operational rules}
\begin{itemize}
\item No default supported status.
\item Attach evidence by reference.
\item Preserve refuted claims for audit.
\item Expire context-sensitive claims when inputs change.
\end{itemize}
\textbf{Extended analysis.} State analysis. Claim Lifecycle is meaningful only if its causal inputs and durable outputs can be reconstructed after process loss. The runtime should record the transcript head, actor, contract, bounded payload, and evidence relationship needed to replay the step. Concurrency is permitted where inputs are independent, but acceptance remains serial at the authority boundary. This preserves the useful diversity of multiple leaves without allowing one stale or theatrical response to overwrite current state.

\textbf{Failure signatures}
\begin{itemize}
\item Output is accepted from an obsolete transcript head.
\item Raw hidden narration replaces bounded operational artifacts.
\item A cursor or claim status cannot be reconstructed after restart.
\item The system keeps generating despite no new evidence or dependency change.
\end{itemize}
\textbf{Implementation or verification note.} Test note. Exercise claim lifecycle under restart, delayed output, cache loss, duplicate delivery, and conflicting events. The passing behavior is the explicit acceptance statement on this plate; the failing behavior must remain journaled rather than being hidden by a regenerated answer.

\begin{acceptbox}
\textbf{Acceptance statement.} The lifecycle passes when a refuted claim cannot reappear in a future context pack without its status.
\end{acceptbox}
\newpage
\platetitle{44}{Evidence Coverage}{FIFTH II}
\prov{[CB] Source-derived}
\fontsize{8.85}{10.45}\selectfont
\begin{corebox}
\textbf{Core statement.} Evidence coverage measures supported claims divided by total claims.
\end{corebox}
The metric does not measure intelligence. It measures whether the runtime is preserving evidence and refusing to collapse proposal into truth. Coverage should be interpreted by claim class; some exploratory hypotheses may remain intentionally unverified while release claims require complete support.

\textbf{Garden consequence.} Garden health can use evidence coverage as a fruit-quality measure rather than a bloom-size measure.

\textbf{Operational rules}
\begin{itemize}
\item Define which claims require support.
\item Exclude purely stylistic statements.
\item Report unresolved counts.
\item Never optimize by emitting fewer meaningful claims solely to raise the ratio.
\end{itemize}
\textbf{Extended analysis.} State analysis. Evidence Coverage is meaningful only if its causal inputs and durable outputs can be reconstructed after process loss. The runtime should record the transcript head, actor, contract, bounded payload, and evidence relationship needed to replay the step. Concurrency is permitted where inputs are independent, but acceptance remains serial at the authority boundary. This preserves the useful diversity of multiple leaves without allowing one stale or theatrical response to overwrite current state.

\textbf{Failure signatures}
\begin{itemize}
\item Output is accepted from an obsolete transcript head.
\item Raw hidden narration replaces bounded operational artifacts.
\item A cursor or claim status cannot be reconstructed after restart.
\item The system keeps generating despite no new evidence or dependency change.
\end{itemize}
\textbf{Implementation or verification note.} Test note. Exercise evidence coverage under restart, delayed output, cache loss, duplicate delivery, and conflicting events. The passing behavior is the explicit acceptance statement on this plate; the failing behavior must remain journaled rather than being hidden by a regenerated answer.

\begin{lstlisting}
E_c = N_supported_claims / N_total_claims
\end{lstlisting}
\begin{acceptbox}
\textbf{Acceptance statement.} The metric passes when its numerator and denominator can be reconstructed from journal events.
\end{acceptbox}
\newpage
\platetitle{45}{Decision Acceptance}{FIFTH II}
\prov{[CB] Source-derived}
\fontsize{8.85}{10.45}\selectfont
\begin{corebox}
\textbf{Core statement.} Decision acceptance measures accepted decisions divided by decision attempts.
\end{corebox}
A high value is not automatically good; it may indicate weak gating. The metric must be read with rejection reasons, failure rates, evidence coverage, and outcome quality. A healthy system rejects stale, unsafe, redundant, and unsupported proposals.

\textbf{Garden consequence.} The Garden interprets rejection as pruning. Pruning can improve long-term growth even when it lowers immediate acceptance.

\textbf{Operational rules}
\begin{itemize}
\item Report accepted, rejected, blocked, and stale separately.
\item Classify reasons.
\item Link accepted decisions to fruit.
\item Link rejected decisions to compost or quarantine.
\end{itemize}
\textbf{Extended analysis.} State analysis. Decision Acceptance is meaningful only if its causal inputs and durable outputs can be reconstructed after process loss. The runtime should record the transcript head, actor, contract, bounded payload, and evidence relationship needed to replay the step. Concurrency is permitted where inputs are independent, but acceptance remains serial at the authority boundary. This preserves the useful diversity of multiple leaves without allowing one stale or theatrical response to overwrite current state.

\textbf{Failure signatures}
\begin{itemize}
\item Output is accepted from an obsolete transcript head.
\item Raw hidden narration replaces bounded operational artifacts.
\item A cursor or claim status cannot be reconstructed after restart.
\item The system keeps generating despite no new evidence or dependency change.
\end{itemize}
\textbf{Implementation or verification note.} Test note. Exercise decision acceptance under restart, delayed output, cache loss, duplicate delivery, and conflicting events. The passing behavior is the explicit acceptance statement on this plate; the failing behavior must remain journaled rather than being hidden by a regenerated answer.

\begin{lstlisting}
A_d = N_accepted_decisions / N_decision_attempts
\end{lstlisting}
\begin{acceptbox}
\textbf{Acceptance statement.} The metric passes when an operator can distinguish productive pruning from arbitrary refusal.
\end{acceptbox}
\newpage
\platetitle{46}{Stale Rejection Metric}{FIFTH II}
\prov{[CB] Source-derived}
\fontsize{8.85}{10.45}\selectfont
\begin{corebox}
\textbf{Core statement.} The stale rejection count measures obsolete proposals prevented from crossing the authority boundary.
\end{corebox}
A nonzero count can be evidence that freshness controls work. A rapidly rising count may indicate excessive latency, over-parallelization, or poor context synchronization. The metric should be paired with time-to-rebase and repeated-stale rate.

\textbf{Garden consequence.} In Garden language, this measures pollen that arrived after the receiving bloom had changed state.

\textbf{Operational rules}
\begin{itemize}
\item Count stale proposals.
\item Record age and delay source.
\item Bound rebase attempts.
\item Tune concurrency rather than disabling freshness.
\end{itemize}
\textbf{Extended analysis.} State analysis. Stale Rejection Metric is meaningful only if its causal inputs and durable outputs can be reconstructed after process loss. The runtime should record the transcript head, actor, contract, bounded payload, and evidence relationship needed to replay the step. Concurrency is permitted where inputs are independent, but acceptance remains serial at the authority boundary. This preserves the useful diversity of multiple leaves without allowing one stale or theatrical response to overwrite current state.

\textbf{Failure signatures}
\begin{itemize}
\item Output is accepted from an obsolete transcript head.
\item Raw hidden narration replaces bounded operational artifacts.
\item A cursor or claim status cannot be reconstructed after restart.
\item The system keeps generating despite no new evidence or dependency change.
\end{itemize}
\textbf{Implementation or verification note.} Test note. Exercise stale rejection metric under restart, delayed output, cache loss, duplicate delivery, and conflicting events. The passing behavior is the explicit acceptance statement on this plate; the failing behavior must remain journaled rather than being hidden by a regenerated answer.

\begin{acceptbox}
\textbf{Acceptance statement.} The metric passes when each count maps to a specific rejected event and current head.
\end{acceptbox}
\newpage
\platetitle{47}{Two-Sided Personality Record}{FIFTH II}
\prov{[CB] Source-derived}
\fontsize{8.85}{10.45}\selectfont
\begin{corebox}
\textbf{Core statement.} Persona output has an operational inside and a deterministic visual outside.
\end{corebox}
The inside contains identity, role, duties, observations, constraints, hypotheses, candidates, objections, evidence references, confidence, decision, and next action. The outside contains stable visual seed, semantic dimensions, color, glyph, and presentation hints.

\textbf{Garden consequence.} The Garden preserves this split. The organism may look different across renderers, but visual identity cannot change lifecycle state or authority.

\textbf{Operational rules}
\begin{itemize}
\item Persist semantic dimensions, not terminal escape sequences.
\item Keep raw hidden reasoning out.
\item Derive visuals deterministically.
\item Allow ASCII fallback.
\end{itemize}
\textbf{Extended analysis.} State analysis. Two-Sided Personality Record is meaningful only if its causal inputs and durable outputs can be reconstructed after process loss. The runtime should record the transcript head, actor, contract, bounded payload, and evidence relationship needed to replay the step. Concurrency is permitted where inputs are independent, but acceptance remains serial at the authority boundary. This preserves the useful diversity of multiple leaves without allowing one stale or theatrical response to overwrite current state.

\textbf{Failure signatures}
\begin{itemize}
\item Output is accepted from an obsolete transcript head.
\item Raw hidden narration replaces bounded operational artifacts.
\item A cursor or claim status cannot be reconstructed after restart.
\item The system keeps generating despite no new evidence or dependency change.
\end{itemize}
\textbf{Implementation or verification note.} Test note. Exercise two-sided personality record under restart, delayed output, cache loss, duplicate delivery, and conflicting events. The passing behavior is the explicit acceptance statement on this plate; the failing behavior must remain journaled rather than being hidden by a regenerated answer.

\begin{acceptbox}
\textbf{Acceptance statement.} The record passes when a restart reproduces the same semantic state and equivalent visual identity.
\end{acceptbox}
\newpage
\platetitle{48}{Deterministic Visual Identity}{FIFTH II}
\prov{[CB] Source-derived}
\fontsize{8.85}{10.45}\selectfont
\begin{corebox}
\textbf{Core statement.} Visual identity is derived from stable fields rather than true randomness.
\end{corebox}
A seed may combine full persona name, run identity, event sequence, and state digest. Semantic dimensions such as confidence, novelty, urgency, agreement, evidence, and risk can map into a bounded color space and glyph set. The renderer chooses truecolor, 256-color, or ASCII fallback.

\textbf{Garden consequence.} Garden visualization can use growth and lifecycle symbols without making appearance part of the database contract.

\textbf{Operational rules}
\begin{itemize}
\item Stable input produces stable output.
\item No visual field grants permission.
\item Accessibility requires text labels.
\item Missing color support degrades gracefully.
\end{itemize}
\textbf{Extended analysis.} State analysis. Deterministic Visual Identity is meaningful only if its causal inputs and durable outputs can be reconstructed after process loss. The runtime should record the transcript head, actor, contract, bounded payload, and evidence relationship needed to replay the step. Concurrency is permitted where inputs are independent, but acceptance remains serial at the authority boundary. This preserves the useful diversity of multiple leaves without allowing one stale or theatrical response to overwrite current state.

\textbf{Failure signatures}
\begin{itemize}
\item Output is accepted from an obsolete transcript head.
\item Raw hidden narration replaces bounded operational artifacts.
\item A cursor or claim status cannot be reconstructed after restart.
\item The system keeps generating despite no new evidence or dependency change.
\end{itemize}
\textbf{Implementation or verification note.} Test note. Exercise deterministic visual identity under restart, delayed output, cache loss, duplicate delivery, and conflicting events. The passing behavior is the explicit acceptance statement on this plate; the failing behavior must remain journaled rather than being hidden by a regenerated answer.

\begin{acceptbox}
\textbf{Acceptance statement.} Visual identity passes when the same event renders recognizably across terminal and web without semantic divergence.
\end{acceptbox}
\newpage
\platetitle{49}{Optional Two-Way Chat}{FIFTH II}
\prov{[CB] Source-derived}
\fontsize{8.85}{10.45}\selectfont
\begin{corebox}
\textbf{Core statement.} Operator chat is optional, per-run, and non-executable.
\end{corebox}
Accepted chat enters the transcript as conversation. If a personality infers that a mutation is desired, it emits a separate typed capability request. The user can inspect, approve, or reject it. Chat survives restart because the transcript, cursors, facts, evidence, and checkpoints are durable; the live connection is merely transport.

\textbf{Garden consequence.} Within the Garden, chat is pollen from the operator, not direct control voltage.

\textbf{Operational rules}
\begin{itemize}
\item Disabled by default.
\item Validate size, audience, state, nonce, and head.
\item Separate conversation and internal collaboration.
\item Never interpret text box contents as shell.
\end{itemize}
\textbf{Extended analysis.} State analysis. Optional Two-Way Chat is meaningful only if its causal inputs and durable outputs can be reconstructed after process loss. The runtime should record the transcript head, actor, contract, bounded payload, and evidence relationship needed to replay the step. Concurrency is permitted where inputs are independent, but acceptance remains serial at the authority boundary. This preserves the useful diversity of multiple leaves without allowing one stale or theatrical response to overwrite current state.

\textbf{Failure signatures}
\begin{itemize}
\item Output is accepted from an obsolete transcript head.
\item Raw hidden narration replaces bounded operational artifacts.
\item A cursor or claim status cannot be reconstructed after restart.
\item The system keeps generating despite no new evidence or dependency change.
\end{itemize}
\textbf{Implementation or verification note.} Test note. Exercise optional two-way chat under restart, delayed output, cache loss, duplicate delivery, and conflicting events. The passing behavior is the explicit acceptance statement on this plate; the failing behavior must remain journaled rather than being hidden by a regenerated answer.

\begin{acceptbox}
\textbf{Acceptance statement.} Chat passes when a malicious or accidental command-like message cannot execute without an approved capability.
\end{acceptbox}
\newpage
\platetitle{50}{Second-Fifth Synthesis}{FIFTH II}
\prov{[CB][G]}
\fontsize{8.85}{10.45}\selectfont
\begin{corebox}
\textbf{Core statement.} The second fifth establishes durable cognition without granting models authority.
\end{corebox}
The transcript provides causal order. Cursors and anti-staleness gates preserve freshness. MWF contracts preserve distinct operational biases. Claims remain unverified until evidence supports them. Quiescence is healthy. KV is disposable. Chat is transport. Exactly one next action survives commitment.

\textbf{Garden consequence.} The Garden adds a lifecycle interpretation: leaves consume pollen, blooms synchronize work, pruning rejects unsafe proposals, fruit carries validated state, and dormancy preserves recoverability without active generation.

\textbf{Operational rules}
\begin{itemize}
\item Cognition is durable only through state, not continuous process life.
\item Multiple personalities require one shared truth path.
\item Metrics measure discipline, not intelligence.
\item Idle and death are legitimate.
\end{itemize}
\textbf{Extended analysis.} State analysis. Second-Fifth Synthesis is meaningful only if its causal inputs and durable outputs can be reconstructed after process loss. The runtime should record the transcript head, actor, contract, bounded payload, and evidence relationship needed to replay the step. Concurrency is permitted where inputs are independent, but acceptance remains serial at the authority boundary. This preserves the useful diversity of multiple leaves without allowing one stale or theatrical response to overwrite current state.

\textbf{Failure signatures}
\begin{itemize}
\item Output is accepted from an obsolete transcript head.
\item Raw hidden narration replaces bounded operational artifacts.
\item A cursor or claim status cannot be reconstructed after restart.
\item The system keeps generating despite no new evidence or dependency change.
\end{itemize}
\textbf{Implementation or verification note.} Test note. Exercise second-fifth synthesis under restart, delayed output, cache loss, duplicate delivery, and conflicting events. The passing behavior is the explicit acceptance statement on this plate; the failing behavior must remain journaled rather than being hidden by a regenerated answer.

\begin{acceptbox}
\textbf{Acceptance statement.} The second fifth is complete when a model can disappear and the cognitive system remains reconstructable.
\end{acceptbox}
\newpage
\platetitle{51}{Historical Baselines}{FIFTH III}
\prov{[TP] Source-derived}
\fontsize{8.85}{10.45}\selectfont
\begin{corebox}
\textbf{Core statement.} Historical throughput numbers are observations under incomplete or specific conditions, not hardware constants.
\end{corebox}
The earliest CPU record preserved three decode samples around 3.9 tokens per second and a rounded 7.37 GB footprint but lacked exact model identity, runtime build, thread count, thermals, memory pressure, and competing-process inventory. Later measurements improved provenance but still describe named artifacts and workloads.

\textbf{Garden consequence.} The Garden treats benchmark results as seasonal observations. A past bloom does not guarantee future yield under different models, context, temperature, or soil.

\textbf{Operational rules}
\begin{itemize}
\item Retain dates and provenance.
\item Separate reconstruction from proof.
\item Do not generalize across models by name alone.
\item Record uncertainty explicitly.
\end{itemize}
\textbf{Extended analysis.} Measurement analysis. Historical Baselines should be evaluated with exact artifact identity, runtime build, placement, workload, thermal state, memory pressure, and competing-process inventory. A number without these conditions is a historical note, not a routing promise. The operational question is whether this component increases validated useful output while preserving host responsiveness and recovery. Peak speed may justify a seasonal profile, but only sustained evidence can justify promotion to a default profile.

\textbf{Failure signatures}
\begin{itemize}
\item A benchmark omits model, build, workload, or placement provenance.
\item Requested GPU execution is reported as actual execution without ownership evidence.
\item Memory or thermal pressure is hidden behind a single average throughput number.
\item High utilization produces no corresponding increase in accepted fruit.
\end{itemize}
\textbf{Implementation or verification note.} Operational note. Record before, during, and after measurements for historical baselines, then connect those observations to route decisions and lifecycle transitions. The result should be comparable across cold, warm, and resident runs and should explain resource release after the bloom closes.

\begin{acceptbox}
\textbf{Acceptance statement.} Historical use passes when every quoted rate identifies whether it is measured, reconstructed, estimated, or deprecated.
\end{acceptbox}
\newpage
\platetitle{52}{Benchmark Provenance}{FIFTH III}
\prov{[TP] Source-derived}
\fontsize{8.85}{10.45}\selectfont
\begin{corebox}
\textbf{Core statement.} A benchmark is meaningful only when model, runtime, placement, workload, and host conditions are recorded.
\end{corebox}
Required fields include artifact and hash, quantization, runtime build and commit, backend, CPU threads, GPU layers, batch and micro-batch, KV types, prompt and output lengths, repetitions, warm-up, power mode, temperature, memory pressure, process inventory, and measurement window.

\textbf{Garden consequence.} Garden yield reports add lifecycle phase, petal identity, root condition, pollen volume, fruit count, and compost outcome.

\textbf{Operational rules}
\begin{itemize}
\item Record requested, detected, and actual placement.
\item Repeat a closing cell for drift.
\item Preserve raw output.
\item Distinguish cold and warm cache.
\end{itemize}
\textbf{Extended analysis.} Measurement analysis. Benchmark Provenance should be evaluated with exact artifact identity, runtime build, placement, workload, thermal state, memory pressure, and competing-process inventory. A number without these conditions is a historical note, not a routing promise. The operational question is whether this component increases validated useful output while preserving host responsiveness and recovery. Peak speed may justify a seasonal profile, but only sustained evidence can justify promotion to a default profile.

\textbf{Failure signatures}
\begin{itemize}
\item A benchmark omits model, build, workload, or placement provenance.
\item Requested GPU execution is reported as actual execution without ownership evidence.
\item Memory or thermal pressure is hidden behind a single average throughput number.
\item High utilization produces no corresponding increase in accepted fruit.
\end{itemize}
\textbf{Implementation or verification note.} Operational note. Record before, during, and after measurements for benchmark provenance, then connect those observations to route decisions and lifecycle transitions. The result should be comparable across cold, warm, and resident runs and should explain resource release after the bloom closes.

\begin{acceptbox}
\textbf{Acceptance statement.} Provenance passes when another run can reproduce the command and explain remaining uncontrolled variables.
\end{acceptbox}
\newpage
\platetitle{53}{CPU Executive}{FIFTH III}
\prov{[TP][CB] Source-derived}
\fontsize{8.85}{10.45}\selectfont
\begin{corebox}
\textbf{Core statement.} The CPU executive maintains persistent planning, state, validation, and task issuance.
\end{corebox}
It owns scheduler and event loop, transcript and cursor management, JSON parsing, schema validation, hashing, checkpoint construction, filesystem traversal, process lifecycle, native tests, evidence indexing, and optionally a CPU-resident executive model.

\textbf{Garden consequence.} In Garden terms, the CPU is both stem tissue and a set of petals. Control-plane duties carry pollen and enforce lifecycle; CPU inference provides compute when selected.

\textbf{Operational rules}
\begin{itemize}
\item Do not reduce CPU value to token rate.
\item Reserve capacity for validation and responsiveness.
\item Measure thread placement per model.
\item Keep control work independent of GPU survival.
\end{itemize}
\textbf{Extended analysis.} Measurement analysis. CPU Executive should be evaluated with exact artifact identity, runtime build, placement, workload, thermal state, memory pressure, and competing-process inventory. A number without these conditions is a historical note, not a routing promise. The operational question is whether this component increases validated useful output while preserving host responsiveness and recovery. Peak speed may justify a seasonal profile, but only sustained evidence can justify promotion to a default profile.

\textbf{Failure signatures}
\begin{itemize}
\item A benchmark omits model, build, workload, or placement provenance.
\item Requested GPU execution is reported as actual execution without ownership evidence.
\item Memory or thermal pressure is hidden behind a single average throughput number.
\item High utilization produces no corresponding increase in accepted fruit.
\end{itemize}
\textbf{Implementation or verification note.} Operational note. Record before, during, and after measurements for cpu executive, then connect those observations to route decisions and lifecycle transitions. The result should be comparable across cold, warm, and resident runs and should explain resource release after the bloom closes.

\begin{acceptbox}
\textbf{Acceptance statement.} The executive passes when it can keep the queue, journal, and validators moving while the GPU worker is occupied.
\end{acceptbox}
\newpage
\platetitle{54}{GPU Worker}{FIFTH III}
\prov{[TP][CB] Source-derived}
\fontsize{8.85}{10.45}\selectfont
\begin{corebox}
\textbf{Core statement.} The GPU hosts accelerator-sensitive reasoning, coding, or other high-throughput inference lanes.
\end{corebox}
Routing must distinguish requested provider, detected provider, actual provider, requested offload, actual offload, GPU layers, VRAM allocation, and fallback reason. Device detection alone is not proof that inference executed on the device.

\textbf{Garden consequence.} The GPU is a high-energy petal. It opens for valuable blooms, reports actual placement, and closes when its task or lifecycle budget ends.

\textbf{Operational rules}
\begin{itemize}
\item Measure actual offload.
\item Record VRAM and utilization with process ownership.
\item Avoid duplicate model launches for shared weights.
\item Preserve fallback reasons.
\end{itemize}
\textbf{Extended analysis.} Measurement analysis. GPU Worker should be evaluated with exact artifact identity, runtime build, placement, workload, thermal state, memory pressure, and competing-process inventory. A number without these conditions is a historical note, not a routing promise. The operational question is whether this component increases validated useful output while preserving host responsiveness and recovery. Peak speed may justify a seasonal profile, but only sustained evidence can justify promotion to a default profile.

\textbf{Failure signatures}
\begin{itemize}
\item A benchmark omits model, build, workload, or placement provenance.
\item Requested GPU execution is reported as actual execution without ownership evidence.
\item Memory or thermal pressure is hidden behind a single average throughput number.
\item High utilization produces no corresponding increase in accepted fruit.
\end{itemize}
\textbf{Implementation or verification note.} Operational note. Record before, during, and after measurements for gpu worker, then connect those observations to route decisions and lifecycle transitions. The result should be comparable across cold, warm, and resident runs and should explain resource release after the bloom closes.

\begin{acceptbox}
\textbf{Acceptance statement.} The GPU lane passes when a provider-off or CPU-fallback run cannot claim GPU inference throughput.
\end{acceptbox}
\newpage
\platetitle{55}{Prefill and Decode}{FIFTH III}
\prov{[TP] Source-derived}
\fontsize{8.85}{10.45}\selectfont
\begin{corebox}
\textbf{Core statement.} Prompt processing and token decoding are different performance regimes.
\end{corebox}
A cycle estimate includes prompt length divided by prefill rate plus output length divided by decode rate. End-to-end service time must also include provider latency, tools, validation, checkpointing, retries, and queue delay.

\textbf{Garden consequence.} Garden scheduling uses these terms to estimate how long petals remain open and when a bloom should senesce rather than assuming token rate equals task completion.

\textbf{Operational rules}
\begin{itemize}
\item Measure both prefill and decode.
\item Record context length.
\item Do not compare decode-only arithmetic with end-to-end productivity.
\item Include non-inference stages.
\end{itemize}
\textbf{Extended analysis.} Measurement analysis. Prefill and Decode should be evaluated with exact artifact identity, runtime build, placement, workload, thermal state, memory pressure, and competing-process inventory. A number without these conditions is a historical note, not a routing promise. The operational question is whether this component increases validated useful output while preserving host responsiveness and recovery. Peak speed may justify a seasonal profile, but only sustained evidence can justify promotion to a default profile.

\textbf{Failure signatures}
\begin{itemize}
\item A benchmark omits model, build, workload, or placement provenance.
\item Requested GPU execution is reported as actual execution without ownership evidence.
\item Memory or thermal pressure is hidden behind a single average throughput number.
\item High utilization produces no corresponding increase in accepted fruit.
\end{itemize}
\textbf{Implementation or verification note.} Operational note. Record before, during, and after measurements for prefill and decode, then connect those observations to route decisions and lifecycle transitions. The result should be comparable across cold, warm, and resident runs and should explain resource release after the bloom closes.

\begin{lstlisting}
t_order = L_P/r_prefill + L_O/r_decode + t_provider + t_tools + t_validation + t_checkpoint + t_retry
\end{lstlisting}
\begin{acceptbox}
\textbf{Acceptance statement.} The model passes when predicted service time is compared with observed complete work-order time.
\end{acceptbox}
\newpage
\platetitle{56}{Observed Service Capacity}{FIFTH III}
\prov{[TP] Source-derived}
\fontsize{8.85}{10.45}\selectfont
\begin{corebox}
\textbf{Core statement.} Work-order capacity is bounded by complete service time, not raw decode rate.
\end{corebox}
The observed upper bound is 3600 divided by measured order time. Queueing, retries, validation, thermal breaks, and tool latency reduce realized throughput. A long model output may be less useful than a short validated patch.

\textbf{Garden consequence.} The Garden therefore measures fruit per cycle and time-to-fruit alongside tokens per second.

\textbf{Operational rules}
\begin{itemize}
\item Record arrival rate and completion rate.
\item Separate active compute from waiting.
\item Count accepted and rejected fruit attempts.
\item Report median and tail latency.
\end{itemize}
\textbf{Extended analysis.} Measurement analysis. Observed Service Capacity should be evaluated with exact artifact identity, runtime build, placement, workload, thermal state, memory pressure, and competing-process inventory. A number without these conditions is a historical note, not a routing promise. The operational question is whether this component increases validated useful output while preserving host responsiveness and recovery. Peak speed may justify a seasonal profile, but only sustained evidence can justify promotion to a default profile.

\textbf{Failure signatures}
\begin{itemize}
\item A benchmark omits model, build, workload, or placement provenance.
\item Requested GPU execution is reported as actual execution without ownership evidence.
\item Memory or thermal pressure is hidden behind a single average throughput number.
\item High utilization produces no corresponding increase in accepted fruit.
\end{itemize}
\textbf{Implementation or verification note.} Operational note. Record before, during, and after measurements for observed service capacity, then connect those observations to route decisions and lifecycle transitions. The result should be comparable across cold, warm, and resident runs and should explain resource release after the bloom closes.

\begin{lstlisting}
lambda_observed <= 3600 / t_order
\end{lstlisting}
\begin{acceptbox}
\textbf{Acceptance statement.} Capacity passes when the system can explain why high token throughput did or did not produce more validated work.
\end{acceptbox}
\newpage
\platetitle{57}{VRAM Bandwidth Model}{FIFTH III}
\prov{[TP] Source-derived}
\fontsize{8.85}{10.45}\selectfont
\begin{corebox}
\textbf{Core statement.} Dense decode is often constrained by effective memory bandwidth and model traffic.
\end{corebox}
A simplified rate divides achieved bandwidth by weight traffic plus KV and overhead traffic. The paper used a 7.37 GB example and a 504 GB/s nominal RTX 4070 bandwidth to show that aggressive targets near 66 tokens per second would require implausibly high efficiency before overhead.

\textbf{Garden consequence.} Garden design uses this as a warning against assuming a larger resident bloom will scale linearly. Weight size, cache traffic, kernels, context, and synchronization all consume the same roots and petals.

\textbf{Operational rules}
\begin{itemize}
\item Treat theoretical ceilings as ceilings.
\item Measure achieved bandwidth indirectly with observed rate.
\item Include KV and overhead.
\item Do not extrapolate one artifact to another.
\end{itemize}
\textbf{Extended analysis.} Measurement analysis. VRAM Bandwidth Model should be evaluated with exact artifact identity, runtime build, placement, workload, thermal state, memory pressure, and competing-process inventory. A number without these conditions is a historical note, not a routing promise. The operational question is whether this component increases validated useful output while preserving host responsiveness and recovery. Peak speed may justify a seasonal profile, but only sustained evidence can justify promotion to a default profile.

\textbf{Failure signatures}
\begin{itemize}
\item A benchmark omits model, build, workload, or placement provenance.
\item Requested GPU execution is reported as actual execution without ownership evidence.
\item Memory or thermal pressure is hidden behind a single average throughput number.
\item High utilization produces no corresponding increase in accepted fruit.
\end{itemize}
\textbf{Implementation or verification note.} Operational note. Record before, during, and after measurements for vram bandwidth model, then connect those observations to route decisions and lifecycle transitions. The result should be comparable across cold, warm, and resident runs and should explain resource release after the bloom closes.

\begin{lstlisting}
r_GPU ~= eta * B_VRAM / (M_weights + D_KV + D_overhead)
\end{lstlisting}
\begin{acceptbox}
\textbf{Acceptance statement.} The model passes when estimated ceilings are clearly separated from measured results.
\end{acceptbox}
\newpage
\platetitle{58}{Partial-Offload Penalty}{FIFTH III}
\prov{[TP] Source-derived}
\fontsize{8.85}{10.45}\selectfont
\begin{corebox}
\textbf{Core statement.} A small serial CPU remainder can sharply reduce combined throughput.
\end{corebox}
The harmonic-style placement model shows that 90, 95, and 99 percent GPU-resident work may still fall far below the full-GPU rate when the CPU portion is much slower. Full offload is therefore a qualitative requirement for some dense workers, not merely a small optimization.

\textbf{Garden consequence.} Garden petals must report actual division of labor. A bloom that appears GPU-active may still be root-bound by a serial CPU segment.

\textbf{Operational rules}
\begin{itemize}
\item Measure zero, partial, and full offload with the same workload.
\item Record layer count and fallback.
\item Investigate low results before blaming hardware.
\item Keep control-plane CPU work separate from serial model work.
\end{itemize}
\textbf{Extended analysis.} Measurement analysis. Partial-Offload Penalty should be evaluated with exact artifact identity, runtime build, placement, workload, thermal state, memory pressure, and competing-process inventory. A number without these conditions is a historical note, not a routing promise. The operational question is whether this component increases validated useful output while preserving host responsiveness and recovery. Peak speed may justify a seasonal profile, but only sustained evidence can justify promotion to a default profile.

\textbf{Failure signatures}
\begin{itemize}
\item A benchmark omits model, build, workload, or placement provenance.
\item Requested GPU execution is reported as actual execution without ownership evidence.
\item Memory or thermal pressure is hidden behind a single average throughput number.
\item High utilization produces no corresponding increase in accepted fruit.
\end{itemize}
\textbf{Implementation or verification note.} Operational note. Record before, during, and after measurements for partial-offload penalty, then connect those observations to route decisions and lifecycle transitions. The result should be comparable across cold, warm, and resident runs and should explain resource release after the bloom closes.

\begin{lstlisting}
r_combined = (f/r_G + (1-f)/r_C)^(-1)
\end{lstlisting}
\begin{acceptbox}
\textbf{Acceptance statement.} Placement passes when offload percentage and combined rate are both reported.
\end{acceptbox}
\newpage
\platetitle{59}{Full-Offload Requirement}{FIFTH III}
\prov{[TP] Source-derived}
\fontsize{8.85}{10.45}\selectfont
\begin{corebox}
\textbf{Core statement.} A model that fits nominally in VRAM must still leave room for KV cache, compute buffers, and runtime state.
\end{corebox}
The 7.37 GB example leaves roughly 4.63 GB in a 12 GB budget before context and runtime overhead. Fit must be measured under the target context and batch settings. A model that loads but immediately pressures buffers is not a healthy resident petal.

\textbf{Garden consequence.} The Garden favors smaller petals partly because they leave room for the rest of the organism.

\textbf{Operational rules}
\begin{itemize}
\item Budget weights, KV, buffers, fragmentation, and display use.
\item Require headroom.
\item Test target context.
\item Fail closed on repeated out-of-memory events.
\end{itemize}
\textbf{Extended analysis.} Measurement analysis. Full-Offload Requirement should be evaluated with exact artifact identity, runtime build, placement, workload, thermal state, memory pressure, and competing-process inventory. A number without these conditions is a historical note, not a routing promise. The operational question is whether this component increases validated useful output while preserving host responsiveness and recovery. Peak speed may justify a seasonal profile, but only sustained evidence can justify promotion to a default profile.

\textbf{Failure signatures}
\begin{itemize}
\item A benchmark omits model, build, workload, or placement provenance.
\item Requested GPU execution is reported as actual execution without ownership evidence.
\item Memory or thermal pressure is hidden behind a single average throughput number.
\item High utilization produces no corresponding increase in accepted fruit.
\end{itemize}
\textbf{Implementation or verification note.} Operational note. Record before, during, and after measurements for full-offload requirement, then connect those observations to route decisions and lifecycle transitions. The result should be comparable across cold, warm, and resident runs and should explain resource release after the bloom closes.

\begin{acceptbox}
\textbf{Acceptance statement.} Full offload passes when the target workload completes without memory pressure forcing silent fallback.
\end{acceptbox}
\newpage
\platetitle{60}{Persistence and Useful Throughput}{FIFTH III}
\prov{[TP] Source-derived}
\fontsize{8.85}{10.45}\selectfont
\begin{corebox}
\textbf{Core statement.} Persistent state improves useful throughput by reducing repeated context processing, lost work, and reconstruction.
\end{corebox}
The paper illustrates a worker at the same raw decode rate producing substantially more useful output when waste falls from thirty percent to five percent. Persistence does not make the GPU faster; it makes more of its work survive.

\textbf{Garden consequence.} This is the direct bridge to the Garden. Compost and seeds should reduce repeated growth while allowing active tissue to die.

\textbf{Operational rules}
\begin{itemize}
\item Measure redo and repetition.
\item Separate raw and useful rates.
\item Checkpoint compact state.
\item Retain only context that changes future work.
\end{itemize}
\textbf{Extended analysis.} Measurement analysis. Persistence and Useful Throughput should be evaluated with exact artifact identity, runtime build, placement, workload, thermal state, memory pressure, and competing-process inventory. A number without these conditions is a historical note, not a routing promise. The operational question is whether this component increases validated useful output while preserving host responsiveness and recovery. Peak speed may justify a seasonal profile, but only sustained evidence can justify promotion to a default profile.

\textbf{Failure signatures}
\begin{itemize}
\item A benchmark omits model, build, workload, or placement provenance.
\item Requested GPU execution is reported as actual execution without ownership evidence.
\item Memory or thermal pressure is hidden behind a single average throughput number.
\item High utilization produces no corresponding increase in accepted fruit.
\end{itemize}
\textbf{Implementation or verification note.} Operational note. Record before, during, and after measurements for persistence and useful throughput, then connect those observations to route decisions and lifecycle transitions. The result should be comparable across cold, warm, and resident runs and should explain resource release after the bloom closes.

\begin{lstlisting}
r_useful = r_raw * (1 - d)
\end{lstlisting}
\begin{acceptbox}
\textbf{Acceptance statement.} Persistence passes when restart and compaction reduce repeated work without changing factual state.
\end{acceptbox}
\newpage
\platetitle{61}{CPU/GPU Pipeline Overlap}{FIFTH III}
\prov{[TP][CB] Source-derived}
\fontsize{8.85}{10.45}\selectfont
\begin{corebox}
\textbf{Core statement.} Useful utilization comes from overlapping independent work, not synthetic load.
\end{corebox}
While the GPU performs inference or coding, the CPU can validate the prior result, traverse files, compile, parse, build the next context, update journals, and prepare a bounded work order. Continuous batching may allow logical personalities to share one model pool instead of racing duplicate launches.

\textbf{Garden consequence.} The Garden sees this as several petals opening around one stem, exchanging pollen at controlled boundaries.

\textbf{Operational rules}
\begin{itemize}
\item Double-buffer preparation and validation.
\item Use admission control.
\item Avoid duplicate weight residency.
\item Preserve foreground responsiveness.
\end{itemize}
\textbf{Extended analysis.} Measurement analysis. CPU/GPU Pipeline Overlap should be evaluated with exact artifact identity, runtime build, placement, workload, thermal state, memory pressure, and competing-process inventory. A number without these conditions is a historical note, not a routing promise. The operational question is whether this component increases validated useful output while preserving host responsiveness and recovery. Peak speed may justify a seasonal profile, but only sustained evidence can justify promotion to a default profile.

\textbf{Failure signatures}
\begin{itemize}
\item A benchmark omits model, build, workload, or placement provenance.
\item Requested GPU execution is reported as actual execution without ownership evidence.
\item Memory or thermal pressure is hidden behind a single average throughput number.
\item High utilization produces no corresponding increase in accepted fruit.
\end{itemize}
\textbf{Implementation or verification note.} Operational note. Record before, during, and after measurements for cpu/gpu pipeline overlap, then connect those observations to route decisions and lifecycle transitions. The result should be comparable across cold, warm, and resident runs and should explain resource release after the bloom closes.

\begin{acceptbox}
\textbf{Acceptance statement.} Overlap passes when higher utilization corresponds to more validated work rather than busy loops.
\end{acceptbox}
\newpage
\platetitle{62}{RAM as Warm Tier}{FIFTH III}
\prov{[CB][TP] Source-derived}
\fontsize{8.85}{10.45}\selectfont
\begin{corebox}
\textbf{Core statement.} RAM holds live transcript projections, indexes, bounded context packs, warm model state, and optional caches.
\end{corebox}
It is a shared warm tier between durable NVMe and scarce VRAM. Authoritative content remains reconstructable from durable records. RAM pressure must be measured because an oversized warm tier can damage the workstation even when VRAM appears healthy.

\textbf{Garden consequence.} In the Garden, RAM is nutrient-rich circulating tissue, not permanent identity.

\textbf{Operational rules}
\begin{itemize}
\item Classify each allocation by rebuildability.
\item Bound context packs.
\item Track pressure and paging.
\item Evict caches before facts.
\end{itemize}
\textbf{Extended analysis.} Measurement analysis. RAM as Warm Tier should be evaluated with exact artifact identity, runtime build, placement, workload, thermal state, memory pressure, and competing-process inventory. A number without these conditions is a historical note, not a routing promise. The operational question is whether this component increases validated useful output while preserving host responsiveness and recovery. Peak speed may justify a seasonal profile, but only sustained evidence can justify promotion to a default profile.

\textbf{Failure signatures}
\begin{itemize}
\item A benchmark omits model, build, workload, or placement provenance.
\item Requested GPU execution is reported as actual execution without ownership evidence.
\item Memory or thermal pressure is hidden behind a single average throughput number.
\item High utilization produces no corresponding increase in accepted fruit.
\end{itemize}
\textbf{Implementation or verification note.} Operational note. Record before, during, and after measurements for ram as warm tier, then connect those observations to route decisions and lifecycle transitions. The result should be comparable across cold, warm, and resident runs and should explain resource release after the bloom closes.

\begin{acceptbox}
\textbf{Acceptance statement.} RAM policy passes when the system can reclaim warm state without losing checkpoints or evidence.
\end{acceptbox}
\newpage
\platetitle{63}{NVMe as Durable Tier}{FIFTH III}
\prov{[CB][TP] Source-derived}
\fontsize{8.85}{10.45}\selectfont
\begin{corebox}
\textbf{Core statement.} NVMe stores journals, checkpoints, evidence, cold summaries, model files, and historical artifacts.
\end{corebox}
The normal data direction is NVMe to RAM to VRAM. Memory mapping may allow the operating system to satisfy pages from storage, but LeafOS does not claim expert pinning or managed SSD/RAM/VRAM tiering unless implemented and measured.

\textbf{Garden consequence.} The Garden treats NVMe as soil and seed storage: durable, slower, and capable of regrowth after active tissue disappears.

\textbf{Operational rules}
\begin{itemize}
\item Separate storage contract from speculative tier management.
\item Record cold and warm behavior.
\item Preserve hashes and revisions.
\item Measure page faults for mapped models.
\end{itemize}
\textbf{Extended analysis.} Measurement analysis. NVMe as Durable Tier should be evaluated with exact artifact identity, runtime build, placement, workload, thermal state, memory pressure, and competing-process inventory. A number without these conditions is a historical note, not a routing promise. The operational question is whether this component increases validated useful output while preserving host responsiveness and recovery. Peak speed may justify a seasonal profile, but only sustained evidence can justify promotion to a default profile.

\textbf{Failure signatures}
\begin{itemize}
\item A benchmark omits model, build, workload, or placement provenance.
\item Requested GPU execution is reported as actual execution without ownership evidence.
\item Memory or thermal pressure is hidden behind a single average throughput number.
\item High utilization produces no corresponding increase in accepted fruit.
\end{itemize}
\textbf{Implementation or verification note.} Operational note. Record before, during, and after measurements for nvme as durable tier, then connect those observations to route decisions and lifecycle transitions. The result should be comparable across cold, warm, and resident runs and should explain resource release after the bloom closes.

\begin{acceptbox}
\textbf{Acceptance statement.} The durable tier passes when a cold boot can reconstruct the run without hidden RAM-only assumptions.
\end{acceptbox}
\newpage
\platetitle{64}{Useful Utilization Governor}{FIFTH III}
\prov{[AA][CB] Source-derived}
\fontsize{8.85}{10.45}\selectfont
\begin{corebox}
\textbf{Core statement.} The governor targets productive occupancy during approved active work, not decorative load.
\end{corebox}
Inputs include queue, hardware, foreground activity, provider health, budget, thermals, and telemetry freshness. Decisions include claim, defer, pace, recover, pressure, or quiesce. CPU and GPU bands are useful-work targets, not quotas.

\textbf{Garden consequence.} The Garden governor decides which petals open, how long they remain active, and when roots require rest.

\textbf{Operational rules}
\begin{itemize}
\item Never create meaningless work to raise utilization.
\item Protect foreground responsiveness.
\item Record reason codes.
\item Use backpressure under memory or thermal pressure.
\end{itemize}
\textbf{Extended analysis.} Measurement analysis. Useful Utilization Governor should be evaluated with exact artifact identity, runtime build, placement, workload, thermal state, memory pressure, and competing-process inventory. A number without these conditions is a historical note, not a routing promise. The operational question is whether this component increases validated useful output while preserving host responsiveness and recovery. Peak speed may justify a seasonal profile, but only sustained evidence can justify promotion to a default profile.

\textbf{Failure signatures}
\begin{itemize}
\item A benchmark omits model, build, workload, or placement provenance.
\item Requested GPU execution is reported as actual execution without ownership evidence.
\item Memory or thermal pressure is hidden behind a single average throughput number.
\item High utilization produces no corresponding increase in accepted fruit.
\end{itemize}
\textbf{Implementation or verification note.} Operational note. Record before, during, and after measurements for useful utilization governor, then connect those observations to route decisions and lifecycle transitions. The result should be comparable across cold, warm, and resident runs and should explain resource release after the bloom closes.

\begin{acceptbox}
\textbf{Acceptance statement.} The governor passes when an idle queue produces quiescence rather than artificial load.
\end{acceptbox}
\newpage
\platetitle{65}{Thermal Breaks}{FIFTH III}
\prov{[CB] Source-derived}
\fontsize{8.85}{10.45}\selectfont
\begin{corebox}
\textbf{Core statement.} Thermal state is part of execution truth.
\end{corebox}
The inherited design alternates active compute and thermal break, with conservative pause and resume thresholds. Missing required telemetry triggers quiescence and checkpoint rather than optimistic continuation. Exact thresholds remain host policy, not universal constants.

\textbf{Garden consequence.} Thermal breaks are literal seasons inside the Garden. Roots may require cooldown even when the queue is nonempty.

\textbf{Operational rules}
\begin{itemize}
\item Record temperature source and sample age.
\item Pause before damage or severe throttling.
\item Resume only after hysteresis.
\item Checkpoint before extended break.
\end{itemize}
\textbf{Extended analysis.} Measurement analysis. Thermal Breaks should be evaluated with exact artifact identity, runtime build, placement, workload, thermal state, memory pressure, and competing-process inventory. A number without these conditions is a historical note, not a routing promise. The operational question is whether this component increases validated useful output while preserving host responsiveness and recovery. Peak speed may justify a seasonal profile, but only sustained evidence can justify promotion to a default profile.

\textbf{Failure signatures}
\begin{itemize}
\item A benchmark omits model, build, workload, or placement provenance.
\item Requested GPU execution is reported as actual execution without ownership evidence.
\item Memory or thermal pressure is hidden behind a single average throughput number.
\item High utilization produces no corresponding increase in accepted fruit.
\end{itemize}
\textbf{Implementation or verification note.} Operational note. Record before, during, and after measurements for thermal breaks, then connect those observations to route decisions and lifecycle transitions. The result should be comparable across cold, warm, and resident runs and should explain resource release after the bloom closes.

\begin{acceptbox}
\textbf{Acceptance statement.} Thermal policy passes when overheating causes a controlled state transition rather than an unexplained performance collapse.
\end{acceptbox}
\newpage
\platetitle{66}{Telemetry Truth Contract}{FIFTH III}
\prov{[AA][CB][TP] Source-derived}
\fontsize{8.85}{10.45}\selectfont
\begin{corebox}
\textbf{Core statement.} Requested, detected, and actual execution must be reported separately.
\end{corebox}
Missing configuration stays visible. Unknown identities remain unknown. Provider-off runs cannot claim model throughput. GPU activity is not assigned to LeafOS without process ownership. Expected missing telemetry is unavailable, not pass.

\textbf{Garden consequence.} Garden health dashboards must follow the same rule: a pretty flowering animation cannot imply active fruit production.

\textbf{Operational rules}
\begin{itemize}
\item Include sample source and age.
\item Attach process ownership.
\item Keep model and fixture throughput separate.
\item Make unavailable states explicit.
\end{itemize}
\textbf{Extended analysis.} Measurement analysis. Telemetry Truth Contract should be evaluated with exact artifact identity, runtime build, placement, workload, thermal state, memory pressure, and competing-process inventory. A number without these conditions is a historical note, not a routing promise. The operational question is whether this component increases validated useful output while preserving host responsiveness and recovery. Peak speed may justify a seasonal profile, but only sustained evidence can justify promotion to a default profile.

\textbf{Failure signatures}
\begin{itemize}
\item A benchmark omits model, build, workload, or placement provenance.
\item Requested GPU execution is reported as actual execution without ownership evidence.
\item Memory or thermal pressure is hidden behind a single average throughput number.
\item High utilization produces no corresponding increase in accepted fruit.
\end{itemize}
\textbf{Implementation or verification note.} Operational note. Record before, during, and after measurements for telemetry truth contract, then connect those observations to route decisions and lifecycle transitions. The result should be comparable across cold, warm, and resident runs and should explain resource release after the bloom closes.

\begin{acceptbox}
\textbf{Acceptance statement.} Telemetry passes when terminal, web, journal, and report agree on what was actually observed.
\end{acceptbox}
\newpage
\platetitle{67}{Medium-MoE Qualification}{FIFTH III}
\prov{[TP] Source-derived}
\fontsize{8.85}{10.45}\selectfont
\begin{corebox}
\textbf{Core statement.} Medium mixture-of-experts candidates require measured qualification rather than enthusiasm about parameter count.
\end{corebox}
Each candidate records total and active parameters, expert topology, quantized artifact, pinned revision, hash, context, explicit runtime command, and pre-run thresholds. One report cannot promote a model automatically.

\textbf{Garden consequence.} The Garden treats a large MoE as a seasonal canopy experiment, not the default organism. It may bloom only after proving host fit, survival, and useful yield.

\textbf{Operational rules}
\begin{itemize}
\item Qualify cold, warm, and resident runs.
\item Record restart, OOM, responsiveness, and pressure.
\item Require operator review.
\item Do not claim managed expert tiering if absent.
\end{itemize}
\textbf{Extended analysis.} Measurement analysis. Medium-MoE Qualification should be evaluated with exact artifact identity, runtime build, placement, workload, thermal state, memory pressure, and competing-process inventory. A number without these conditions is a historical note, not a routing promise. The operational question is whether this component increases validated useful output while preserving host responsiveness and recovery. Peak speed may justify a seasonal profile, but only sustained evidence can justify promotion to a default profile.

\textbf{Failure signatures}
\begin{itemize}
\item A benchmark omits model, build, workload, or placement provenance.
\item Requested GPU execution is reported as actual execution without ownership evidence.
\item Memory or thermal pressure is hidden behind a single average throughput number.
\item High utilization produces no corresponding increase in accepted fruit.
\end{itemize}
\textbf{Implementation or verification note.} Operational note. Record before, during, and after measurements for medium-moe qualification, then connect those observations to route decisions and lifecycle transitions. The result should be comparable across cold, warm, and resident runs and should explain resource release after the bloom closes.

\begin{acceptbox}
\textbf{Acceptance statement.} Qualification passes when all required reports exist and promotion remains a separate decision.
\end{acceptbox}
\newpage
\platetitle{68}{Qualification Portfolio}{FIFTH III}
\prov{[TP] Source-derived}
\fontsize{8.85}{10.45}\selectfont
\begin{corebox}
\textbf{Core statement.} A candidate is evaluated across short cold-cache, medium warm-cache, and longer resident-foreground runs.
\end{corebox}
The portfolio records provider survival, restart count, out-of-memory events, RAM and VRAM pressure, foreground responsiveness, token rates, validated useful changes per hour, and drift. The long run tests coexistence with the workstation rather than isolated peak speed.

\textbf{Garden consequence.} Garden selection therefore evaluates seasons, not a single sunny afternoon.

\textbf{Operational rules}
\begin{itemize}
\item Use fixed workloads.
\item Repeat final cells for drift.
\item Preserve failed reports.
\item Compare useful output, not only tokens.
\end{itemize}
\textbf{Extended analysis.} Measurement analysis. Qualification Portfolio should be evaluated with exact artifact identity, runtime build, placement, workload, thermal state, memory pressure, and competing-process inventory. A number without these conditions is a historical note, not a routing promise. The operational question is whether this component increases validated useful output while preserving host responsiveness and recovery. Peak speed may justify a seasonal profile, but only sustained evidence can justify promotion to a default profile.

\textbf{Failure signatures}
\begin{itemize}
\item A benchmark omits model, build, workload, or placement provenance.
\item Requested GPU execution is reported as actual execution without ownership evidence.
\item Memory or thermal pressure is hidden behind a single average throughput number.
\item High utilization produces no corresponding increase in accepted fruit.
\end{itemize}
\textbf{Implementation or verification note.} Operational note. Record before, during, and after measurements for qualification portfolio, then connect those observations to route decisions and lifecycle transitions. The result should be comparable across cold, warm, and resident runs and should explain resource release after the bloom closes.

\begin{acceptbox}
\textbf{Acceptance statement.} The portfolio passes when a candidate can be rejected for instability despite an impressive peak rate.
\end{acceptbox}
\newpage
\platetitle{69}{Plan, Resolve, Apply, Verify}{FIFTH III}
\prov{[AA] Source-derived}
\fontsize{8.85}{10.45}\selectfont
\begin{corebox}
\textbf{Core statement.} Model acquisition is a guarded four-stage workflow.
\end{corebox}
Plan is offline and declarative. Resolve contacts metadata and pins exact identities. Apply performs the download boundary. Verify checks size, hash, path, and role mapping. The stages should be visible independently in CLI and reports.

\textbf{Garden consequence.} Garden seed stock follows the same lifecycle: cataloged, resolved, acquired, verified, dormant, activated, and eventually pruned or archived.

\textbf{Operational rules}
\begin{itemize}
\item No silent apply.
\item Permit role-limited acquisition.
\item Fail fast on invalid credentials.
\item Retain revision pins.
\end{itemize}
\textbf{Extended analysis.} Measurement analysis. Plan, Resolve, Apply, Verify should be evaluated with exact artifact identity, runtime build, placement, workload, thermal state, memory pressure, and competing-process inventory. A number without these conditions is a historical note, not a routing promise. The operational question is whether this component increases validated useful output while preserving host responsiveness and recovery. Peak speed may justify a seasonal profile, but only sustained evidence can justify promotion to a default profile.

\textbf{Failure signatures}
\begin{itemize}
\item A benchmark omits model, build, workload, or placement provenance.
\item Requested GPU execution is reported as actual execution without ownership evidence.
\item Memory or thermal pressure is hidden behind a single average throughput number.
\item High utilization produces no corresponding increase in accepted fruit.
\end{itemize}
\textbf{Implementation or verification note.} Operational note. Record before, during, and after measurements for plan, resolve, apply, verify, then connect those observations to route decisions and lifecycle transitions. The result should be comparable across cold, warm, and resident runs and should explain resource release after the bloom closes.

\begin{acceptbox}
\textbf{Acceptance statement.} The workflow passes when selected models can be materialized without downloading every declared optional artifact.
\end{acceptbox}
\newpage
\platetitle{70}{Pack as Composition Unit}{FIFTH III}
\prov{[CB][G] Source-derived and enriched}
\fontsize{8.85}{10.45}\selectfont
\begin{corebox}
\textbf{Core statement.} The model pack is the unit of composition between artifacts, roles, routing, installation, and runtime policy.
\end{corebox}
A pack declares model references, role groups, required and optional artifacts, fallback chains, context limits, hardware budgets, and load states. Declared, installed, and loaded remain separate. A profile selects a practical subset without erasing the catalog.

\textbf{Garden consequence.} The Garden adds lifecycle triggers, senescence conditions, compost policy, and fruit targets to the pack.

\textbf{Operational rules}
\begin{itemize}
\item One pack identity, multiple runtime profiles.
\item Role permissions remain explicit.
\item Optional lanes do not become default downloads.
\item Runtime state reports which profile is active.
\end{itemize}
\textbf{Extended analysis.} Measurement analysis. Pack as Composition Unit should be evaluated with exact artifact identity, runtime build, placement, workload, thermal state, memory pressure, and competing-process inventory. A number without these conditions is a historical note, not a routing promise. The operational question is whether this component increases validated useful output while preserving host responsiveness and recovery. Peak speed may justify a seasonal profile, but only sustained evidence can justify promotion to a default profile.

\textbf{Failure signatures}
\begin{itemize}
\item A benchmark omits model, build, workload, or placement provenance.
\item Requested GPU execution is reported as actual execution without ownership evidence.
\item Memory or thermal pressure is hidden behind a single average throughput number.
\item High utilization produces no corresponding increase in accepted fruit.
\end{itemize}
\textbf{Implementation or verification note.} Operational note. Record before, during, and after measurements for pack as composition unit, then connect those observations to route decisions and lifecycle transitions. The result should be comparable across cold, warm, and resident runs and should explain resource release after the bloom closes.

\begin{acceptbox}
\textbf{Acceptance statement.} Pack composition passes when the same pack can run in a small profile and a larger seasonal profile without changing authority.
\end{acceptbox}
\newpage
\platetitle{71}{Older 2+2 Architecture}{FIFTH III}
\prov{[G] Based on prior Pack 12 design}
\fontsize{8.85}{10.45}\selectfont
\begin{corebox}
\textbf{Core statement.} The 2+2 architecture is the best initial Garden pack because it proves specialization without permanent excess.
\end{corebox}
The primary pair is Brain and Coder. The support pair is Judge and Auxiliary. Auxiliary cleans and compresses input. Brain plans and routes. Coder builds bounded patches and tools. Judge accepts, rejects, or requests retry. Native validators remain outside model authority.

\textbf{Garden consequence.} This design maps cleanly to a small organism with four functional nodes and explicit lifecycle states.

\textbf{Operational rules}
\begin{itemize}
\item Load auxiliary hot or warm.
\item Load brain on demand or warm.
\item Load coder only for code tasks.
\item Load judge before commitment.
\end{itemize}
\textbf{Extended analysis.} Measurement analysis. Older 2+2 Architecture should be evaluated with exact artifact identity, runtime build, placement, workload, thermal state, memory pressure, and competing-process inventory. A number without these conditions is a historical note, not a routing promise. The operational question is whether this component increases validated useful output while preserving host responsiveness and recovery. Peak speed may justify a seasonal profile, but only sustained evidence can justify promotion to a default profile.

\textbf{Failure signatures}
\begin{itemize}
\item A benchmark omits model, build, workload, or placement provenance.
\item Requested GPU execution is reported as actual execution without ownership evidence.
\item Memory or thermal pressure is hidden behind a single average throughput number.
\item High utilization produces no corresponding increase in accepted fruit.
\end{itemize}
\textbf{Implementation or verification note.} Operational note. Record before, during, and after measurements for older 2+2 architecture, then connect those observations to route decisions and lifecycle transitions. The result should be comparable across cold, warm, and resident runs and should explain resource release after the bloom closes.

\begin{acceptbox}
\textbf{Acceptance statement.} The 2+2 pack passes when each role can be independently activated, replaced, unloaded, and recovered.
\end{acceptbox}
\newpage
\platetitle{72}{Hybrid 4+4 Profile}{FIFTH III}
\prov{[G] Based on prior Pack 1.2 selection}
\fontsize{8.85}{10.45}\selectfont
\begin{corebox}
\textbf{Core statement.} The hybrid 4+4 profile is a richer seasonal bloom, not the default resident core.
\end{corebox}
Four Qwen reasoning variants provide peak brain, default brain, context feeder, and fast draft. Four MiniCPM writing variants provide style polish, writer-judge, default writer, and fast extraction. Visual and experimental lanes remain optional toggles.

\textbf{Garden consequence.} The Garden reframes this profile as a temporary canopy: useful for high-value work, expensive to keep fully active, and expected to contract after fruit.

\textbf{Operational rules}
\begin{itemize}
\item Catalog all eight, load only needed roles.
\item Keep peak brain cold by default.
\item Use small writers for preprocessing and polish.
\item Record promotion evidence for experimental lanes.
\end{itemize}
\textbf{Extended analysis.} Measurement analysis. Hybrid 4+4 Profile should be evaluated with exact artifact identity, runtime build, placement, workload, thermal state, memory pressure, and competing-process inventory. A number without these conditions is a historical note, not a routing promise. The operational question is whether this component increases validated useful output while preserving host responsiveness and recovery. Peak speed may justify a seasonal profile, but only sustained evidence can justify promotion to a default profile.

\textbf{Failure signatures}
\begin{itemize}
\item A benchmark omits model, build, workload, or placement provenance.
\item Requested GPU execution is reported as actual execution without ownership evidence.
\item Memory or thermal pressure is hidden behind a single average throughput number.
\item High utilization produces no corresponding increase in accepted fruit.
\end{itemize}
\textbf{Implementation or verification note.} Operational note. Record before, during, and after measurements for hybrid 4+4 profile, then connect those observations to route decisions and lifecycle transitions. The result should be comparable across cold, warm, and resident runs and should explain resource release after the bloom closes.

\begin{acceptbox}
\textbf{Acceptance statement.} The profile passes when it can materialize eight artifacts while loading only a bounded subset for a given task.
\end{acceptbox}
\newpage
\platetitle{73}{Hot, Warm, Cold, Dormant}{FIFTH III}
\prov{[G] Enriched from prior load-state design}
\fontsize{8.85}{10.45}\selectfont
\begin{corebox}
\textbf{Core statement.} Model placement is a lifecycle state rather than a binary installed/not-installed flag.
\end{corebox}
Hot models are resident and immediately available. Warm models retain partial state or fast-load readiness. Cold models are installed but inactive. Dormant models are cataloged or archived and require explicit activation. State transitions are governed by task value, memory pressure, latency target, and lifecycle policy.

\textbf{Garden consequence.} Dormancy is the missing category in many local stacks. It permits a large garden without pretending every plant must be photosynthesizing at once.

\textbf{Operational rules}
\begin{itemize}
\item Report current state per artifact.
\item Define transition cost.
\item Unload after idle or fruit.
\item Never infer residency from installation alone.
\end{itemize}
\textbf{Extended analysis.} Measurement analysis. Hot, Warm, Cold, Dormant should be evaluated with exact artifact identity, runtime build, placement, workload, thermal state, memory pressure, and competing-process inventory. A number without these conditions is a historical note, not a routing promise. The operational question is whether this component increases validated useful output while preserving host responsiveness and recovery. Peak speed may justify a seasonal profile, but only sustained evidence can justify promotion to a default profile.

\textbf{Failure signatures}
\begin{itemize}
\item A benchmark omits model, build, workload, or placement provenance.
\item Requested GPU execution is reported as actual execution without ownership evidence.
\item Memory or thermal pressure is hidden behind a single average throughput number.
\item High utilization produces no corresponding increase in accepted fruit.
\end{itemize}
\textbf{Implementation or verification note.} Operational note. Record before, during, and after measurements for hot, warm, cold, dormant, then connect those observations to route decisions and lifecycle transitions. The result should be comparable across cold, warm, and resident runs and should explain resource release after the bloom closes.

\begin{acceptbox}
\textbf{Acceptance statement.} Load-state policy passes when memory can be reclaimed without losing pack identity or future activation capability.
\end{acceptbox}
\newpage
\platetitle{74}{Useful Output Metrics}{FIFTH III}
\prov{[TP][G]}
\fontsize{8.85}{10.45}\selectfont
\begin{corebox}
\textbf{Core statement.} The primary outcome is validated useful change, not token volume.
\end{corebox}
Metrics may include validated patches per hour, accepted artifacts per cycle, evidence coverage, time-to-fruit, redo fraction, energy per accepted artifact, stale rejection rate, recovery time, and workstation responsiveness. Tokens per second remain diagnostic.

\textbf{Garden consequence.} Garden metrics reward renewal: a small petal that produces reliable fruit and closes may outperform a huge bloom that consumes roots and repeats itself.

\textbf{Operational rules}
\begin{itemize}
\item Pair throughput with acceptance.
\item Pair utilization with completed work.
\item Report energy and memory pressure.
\item Count recovery and pruning outcomes.
\end{itemize}
\textbf{Extended analysis.} Measurement analysis. Useful Output Metrics should be evaluated with exact artifact identity, runtime build, placement, workload, thermal state, memory pressure, and competing-process inventory. A number without these conditions is a historical note, not a routing promise. The operational question is whether this component increases validated useful output while preserving host responsiveness and recovery. Peak speed may justify a seasonal profile, but only sustained evidence can justify promotion to a default profile.

\textbf{Failure signatures}
\begin{itemize}
\item A benchmark omits model, build, workload, or placement provenance.
\item Requested GPU execution is reported as actual execution without ownership evidence.
\item Memory or thermal pressure is hidden behind a single average throughput number.
\item High utilization produces no corresponding increase in accepted fruit.
\end{itemize}
\textbf{Implementation or verification note.} Operational note. Record before, during, and after measurements for useful output metrics, then connect those observations to route decisions and lifecycle transitions. The result should be comparable across cold, warm, and resident runs and should explain resource release after the bloom closes.

\begin{acceptbox}
\textbf{Acceptance statement.} Metrics pass when a lower-token configuration can win for higher validated yield.
\end{acceptbox}
\newpage
\platetitle{75}{Third-Fifth Synthesis}{FIFTH III}
\prov{[TP][AA][CB][G]}
\fontsize{8.85}{10.45}\selectfont
\begin{corebox}
\textbf{Core statement.} The third fifth establishes the physical and compositional limits within which the Garden must live.
\end{corebox}
CPU and GPU have different responsibilities. RAM and NVMe are different tiers. Full offload, context, bandwidth, thermals, and process ownership matter. Persistent state increases useful throughput. Packs separate catalog, profile, installation, and runtime. Smaller 2+2 operation is a stronger default than a permanent 4+4 bloom.

\textbf{Garden consequence.} The Garden does not reject large models. It rejects the assumption that maximum resident allocation is identical to maximum long-term growth.

\textbf{Operational rules}
\begin{itemize}
\item Measure exact artifacts and placements.
\item Use lifecycle states for load.
\item Promote by evidence.
\item Optimize validated yield per resource.
\end{itemize}
\textbf{Extended analysis.} Measurement analysis. Third-Fifth Synthesis should be evaluated with exact artifact identity, runtime build, placement, workload, thermal state, memory pressure, and competing-process inventory. A number without these conditions is a historical note, not a routing promise. The operational question is whether this component increases validated useful output while preserving host responsiveness and recovery. Peak speed may justify a seasonal profile, but only sustained evidence can justify promotion to a default profile.

\textbf{Failure signatures}
\begin{itemize}
\item A benchmark omits model, build, workload, or placement provenance.
\item Requested GPU execution is reported as actual execution without ownership evidence.
\item Memory or thermal pressure is hidden behind a single average throughput number.
\item High utilization produces no corresponding increase in accepted fruit.
\end{itemize}
\textbf{Implementation or verification note.} Operational note. Record before, during, and after measurements for third-fifth synthesis, then connect those observations to route decisions and lifecycle transitions. The result should be comparable across cold, warm, and resident runs and should explain resource release after the bloom closes.

\begin{acceptbox}
\textbf{Acceptance statement.} The third fifth is complete when hardware policy explains not only how to start a model, but why and when to stop it.
\end{acceptbox}
\newpage
\platetitle{76}{Why Infinite Bloom Fails}{FIFTH IV}
\prov{[G] New philosophy}
\fontsize{8.85}{10.45}\selectfont
\begin{corebox}
\textbf{Core statement.} Permanent bloom is biologically unnatural and computationally inefficient.
\end{corebox}
A system that treats continuous activity as health tends to accumulate resident models, stale contexts, duplicate workers, thermal pressure, and unbounded logs. It becomes difficult to distinguish useful persistence from refusal to die. Infinite Bloom confuses liveness with value.

\textbf{Garden consequence.} The Garden preserves continual synchronization while rejecting continual generation and continual allocation. Bloom becomes a bounded phase with entry criteria, fruit targets, and exit conditions.

\textbf{Operational rules}
\begin{itemize}
\item No bloom without a declared fruit target.
\item No active worker without a budget and owner.
\item No permanent context merely because it once mattered.
\item No failure hidden as continued activity.
\end{itemize}
\textbf{Extended analysis.} Systems interpretation. The biological analogy for Why Infinite Bloom Fails is useful only because it separates a real machine responsibility, flow, or lifecycle condition. The analogy must stop where biology would obscure authority or measurement. Every Garden term therefore needs fields, states, owners, and evidence. The purpose is not to claim that software is alive; it is to borrow the proven logic of finite organisms, ecological succession, and material recycling for systems that must remain useful across many hardware and project cycles.

\textbf{Failure signatures}
\begin{itemize}
\item Why Infinite Bloom Fails appears in diagrams but has no machine-readable state.
\item Activity continues without an exit condition or fruit target.
\item Retired material is either kept forever or deleted without provenance.
\item The metaphor is used to bypass the existing LeafOS authority doctrine.
\end{itemize}
\textbf{Implementation or verification note.} Design consequence. A runtime that adopts why infinite bloom fails should expose a measurable indicator, a transition rule, and a failure disposition. Without those three elements the term remains presentation. With them, it becomes a compact systems language that can guide packs, schedulers, interfaces, and long-term maintenance.

\begin{acceptbox}
\textbf{Acceptance statement.} The philosophy passes when stopping can be a successful outcome rather than evidence of system death.
\end{acceptbox}
\newpage
\platetitle{77}{Natural Growth and Death Cycles}{FIFTH IV}
\prov{[G] New philosophy}
\fontsize{8.85}{10.45}\selectfont
\begin{corebox}
\textbf{Core statement.} Nature achieves continuity through replacement, succession, dormancy, reproduction, and recycling rather than by keeping every organism alive forever.
\end{corebox}
Plants germinate, grow, flower, fruit, senesce, and return material to soil. Animals develop, specialize, reproduce, and die. Ecological continuity persists above the lifespan of any individual. The computational analogy is not that software is alive, but that durable systems can separate identity from resident process life.

\textbf{Garden consequence.} LeafOS can therefore pursue indefinite improvement through repeated bounded cycles whose artifacts and lessons survive their workers.

\textbf{Operational rules}
\begin{itemize}
\item Separate organism identity from process lifetime.
\item Carry durable lessons forward.
\item Permit old strategies to die.
\item Preserve diversity for changing conditions.
\end{itemize}
\textbf{Extended analysis.} Systems interpretation. The biological analogy for Natural Growth and Death Cycles is useful only because it separates a real machine responsibility, flow, or lifecycle condition. The analogy must stop where biology would obscure authority or measurement. Every Garden term therefore needs fields, states, owners, and evidence. The purpose is not to claim that software is alive; it is to borrow the proven logic of finite organisms, ecological succession, and material recycling for systems that must remain useful across many hardware and project cycles.

\textbf{Failure signatures}
\begin{itemize}
\item Natural Growth and Death Cycles appears in diagrams but has no machine-readable state.
\item Activity continues without an exit condition or fruit target.
\item Retired material is either kept forever or deleted without provenance.
\item The metaphor is used to bypass the existing LeafOS authority doctrine.
\end{itemize}
\textbf{Implementation or verification note.} Design consequence. A runtime that adopts natural growth and death cycles should expose a measurable indicator, a transition rule, and a failure disposition. Without those three elements the term remains presentation. With them, it becomes a compact systems language that can guide packs, schedulers, interfaces, and long-term maintenance.

\begin{acceptbox}
\textbf{Acceptance statement.} The principle passes when a dead worker is ordinary and a lost checkpoint is extraordinary.
\end{acceptbox}
\newpage
\platetitle{78}{Formal Definition of the Garden}{FIFTH IV}
\prov{[G] New philosophy}
\fontsize{8.85}{10.45}\selectfont
\begin{corebox}
\textbf{Core statement.} The Garden is a bio-inspired lifecycle architecture for distributed local intelligence.
\end{corebox}
It separates environment, power, transport, compute, data movement, cognition, active coordination, durable output, and retired state into named roles. These roles are linked by explicit state transitions and authority boundaries rather than left as metaphor.

\textbf{Garden consequence.} The Garden is larger than one pack or one machine. It may describe a workstation, a cluster of local nodes, a family of model packs, or a long-running project whose individual blooms are short-lived.

\textbf{Operational rules}
\begin{itemize}
\item Every role must be measurable.
\item Every transition must be journaled.
\item Every active phase must have an exit.
\item Every durable claim must retain evidence.
\end{itemize}
\textbf{Extended analysis.} Systems interpretation. The biological analogy for Formal Definition of the Garden is useful only because it separates a real machine responsibility, flow, or lifecycle condition. The analogy must stop where biology would obscure authority or measurement. Every Garden term therefore needs fields, states, owners, and evidence. The purpose is not to claim that software is alive; it is to borrow the proven logic of finite organisms, ecological succession, and material recycling for systems that must remain useful across many hardware and project cycles.

\textbf{Failure signatures}
\begin{itemize}
\item Formal Definition of the Garden appears in diagrams but has no machine-readable state.
\item Activity continues without an exit condition or fruit target.
\item Retired material is either kept forever or deleted without provenance.
\item The metaphor is used to bypass the existing LeafOS authority doctrine.
\end{itemize}
\textbf{Implementation or verification note.} Design consequence. A runtime that adopts formal definition of the garden should expose a measurable indicator, a transition rule, and a failure disposition. Without those three elements the term remains presentation. With them, it becomes a compact systems language that can guide packs, schedulers, interfaces, and long-term maintenance.

\begin{acceptbox}
\textbf{Acceptance statement.} The definition passes when the same concepts can describe both one-PC LeafOS and a future multi-node installation.
\end{acceptbox}
\newpage
\platetitle{79}{Soil}{FIFTH IV}
\prov{[G] New philosophy}
\fontsize{8.85}{10.45}\selectfont
\begin{corebox}
\textbf{Core statement.} Soil is the environmental and resource envelope from which the system grows.
\end{corebox}
It includes chassis, cooling, operating system, filesystem, local model store, network access, storage capacity, security context, operator constraints, and the history available for regrowth. Soil is not active computation; it determines what forms of computation can remain healthy.

\textbf{Garden consequence.} Soil quality changes over time. Fragmented storage, stale dependencies, weak cooling, or unbounded archives can make later blooms less productive even when compute hardware is unchanged.

\textbf{Operational rules}
\begin{itemize}
\item Audit soil before activation.
\item Track free storage and cache growth.
\item Preserve source and evidence classes.
\item Treat dependency drift as a soil change.
\end{itemize}
\textbf{Extended analysis.} Systems interpretation. The biological analogy for Soil is useful only because it separates a real machine responsibility, flow, or lifecycle condition. The analogy must stop where biology would obscure authority or measurement. Every Garden term therefore needs fields, states, owners, and evidence. The purpose is not to claim that software is alive; it is to borrow the proven logic of finite organisms, ecological succession, and material recycling for systems that must remain useful across many hardware and project cycles.

\textbf{Failure signatures}
\begin{itemize}
\item Soil appears in diagrams but has no machine-readable state.
\item Activity continues without an exit condition or fruit target.
\item Retired material is either kept forever or deleted without provenance.
\item The metaphor is used to bypass the existing LeafOS authority doctrine.
\end{itemize}
\textbf{Implementation or verification note.} Design consequence. A runtime that adopts soil should expose a measurable indicator, a transition rule, and a failure disposition. Without those three elements the term remains presentation. With them, it becomes a compact systems language that can guide packs, schedulers, interfaces, and long-term maintenance.

\begin{acceptbox}
\textbf{Acceptance statement.} Soil passes when the lifecycle governor can explain which environmental constraint blocked or shaped a bloom.
\end{acceptbox}
\newpage
\platetitle{80}{Roots}{FIFTH IV}
\prov{[G] New philosophy}
\fontsize{8.85}{10.45}\selectfont
\begin{corebox}
\textbf{Core statement.} Roots are power and survival infrastructure.
\end{corebox}
They include wall power, PSU rails, battery or UPS state, grounding, thermal transport, fan and pump capability, power limits, and the telemetry needed to know whether the organism can continue. Roots supply every petal and leaf but do not decide task truth.

\textbf{Garden consequence.} Root health is a first-class runtime condition. A scheduler that ignores power and heat is not intelligent; it is merely optimistic until throttling or failure performs involuntary pruning.

\textbf{Operational rules}
\begin{itemize}
\item Measure power and temperature when available.
\item Use hysteresis for pause and resume.
\item Distinguish missing telemetry from healthy telemetry.
\item Checkpoint before root-induced dormancy.
\end{itemize}
\textbf{Extended analysis.} Systems interpretation. The biological analogy for Roots is useful only because it separates a real machine responsibility, flow, or lifecycle condition. The analogy must stop where biology would obscure authority or measurement. Every Garden term therefore needs fields, states, owners, and evidence. The purpose is not to claim that software is alive; it is to borrow the proven logic of finite organisms, ecological succession, and material recycling for systems that must remain useful across many hardware and project cycles.

\textbf{Failure signatures}
\begin{itemize}
\item Roots appears in diagrams but has no machine-readable state.
\item Activity continues without an exit condition or fruit target.
\item Retired material is either kept forever or deleted without provenance.
\item The metaphor is used to bypass the existing LeafOS authority doctrine.
\end{itemize}
\textbf{Implementation or verification note.} Design consequence. A runtime that adopts roots should expose a measurable indicator, a transition rule, and a failure disposition. Without those three elements the term remains presentation. With them, it becomes a compact systems language that can guide packs, schedulers, interfaces, and long-term maintenance.

\begin{acceptbox}
\textbf{Acceptance statement.} Roots pass when power or thermal degradation produces a controlled lifecycle response.
\end{acceptbox}
\newpage
\platetitle{81}{Stem}{FIFTH IV}
\prov{[G] New philosophy}
\fontsize{8.85}{10.45}\selectfont
\begin{corebox}
\textbf{Core statement.} The stem is the transport, scheduling, and structural spine.
\end{corebox}
It carries typed pollen between nodes, distributes root resources, maintains queue order, enforces leases, and connects active petals and leaves to durable stores. In implementation, the stem corresponds to the resident event loop, queue, router, inlet, journal, and process supervisor.

\textbf{Garden consequence.} The stem is deliberately less brain-like than the leaf. It should be reliable, boring, and inspectable. Complex interpretation belongs in leaves; legal sequencing belongs in the stem.

\textbf{Operational rules}
\begin{itemize}
\item Keep transport deterministic.
\item Use backpressure.
\item Do not embed hidden model policy in the queue.
\item Preserve node identity and ownership.
\end{itemize}
\textbf{Extended analysis.} Systems interpretation. The biological analogy for Stem is useful only because it separates a real machine responsibility, flow, or lifecycle condition. The analogy must stop where biology would obscure authority or measurement. Every Garden term therefore needs fields, states, owners, and evidence. The purpose is not to claim that software is alive; it is to borrow the proven logic of finite organisms, ecological succession, and material recycling for systems that must remain useful across many hardware and project cycles.

\textbf{Failure signatures}
\begin{itemize}
\item Stem appears in diagrams but has no machine-readable state.
\item Activity continues without an exit condition or fruit target.
\item Retired material is either kept forever or deleted without provenance.
\item The metaphor is used to bypass the existing LeafOS authority doctrine.
\end{itemize}
\textbf{Implementation or verification note.} Design consequence. A runtime that adopts stem should expose a measurable indicator, a transition rule, and a failure disposition. Without those three elements the term remains presentation. With them, it becomes a compact systems language that can guide packs, schedulers, interfaces, and long-term maintenance.

\begin{acceptbox}
\textbf{Acceptance statement.} The stem passes when data and control can be traced end-to-end without reading model prose.
\end{acceptbox}
\newpage
\platetitle{82}{Petals}{FIFTH IV}
\prov{[G] New philosophy}
\fontsize{8.85}{10.45}\selectfont
\begin{corebox}
\textbf{Core statement.} Petals are compute surfaces that open for work and close when their contribution ends.
\end{corebox}
A petal may be a CPU worker, GPU slot, NPU accelerator, model instance, compiler pool, validator process, or other bounded executor. Petals expose capacity and specialize by workload. They do not own durable truth or authority merely because they consume the most power.

\textbf{Garden consequence.} The petal metaphor emphasizes multiplicity and transience. Many small petals may cooperate around one bloom; an expensive petal may remain closed until a high-value task justifies it.

\textbf{Operational rules}
\begin{itemize}
\item Declare role, owner, budget, and closure condition.
\item Record actual device placement.
\item Avoid duplicate resident weights.
\item Close idle petals.
\end{itemize}
\textbf{Extended analysis.} Systems interpretation. The biological analogy for Petals is useful only because it separates a real machine responsibility, flow, or lifecycle condition. The analogy must stop where biology would obscure authority or measurement. Every Garden term therefore needs fields, states, owners, and evidence. The purpose is not to claim that software is alive; it is to borrow the proven logic of finite organisms, ecological succession, and material recycling for systems that must remain useful across many hardware and project cycles.

\textbf{Failure signatures}
\begin{itemize}
\item Petals appears in diagrams but has no machine-readable state.
\item Activity continues without an exit condition or fruit target.
\item Retired material is either kept forever or deleted without provenance.
\item The metaphor is used to bypass the existing LeafOS authority doctrine.
\end{itemize}
\textbf{Implementation or verification note.} Design consequence. A runtime that adopts petals should expose a measurable indicator, a transition rule, and a failure disposition. Without those three elements the term remains presentation. With them, it becomes a compact systems language that can guide packs, schedulers, interfaces, and long-term maintenance.

\begin{acceptbox}
\textbf{Acceptance statement.} Petals pass when the system can enumerate active compute surfaces and explain the work each is doing.
\end{acceptbox}
\newpage
\platetitle{83}{Pollen}{FIFTH IV}
\prov{[G] New philosophy}
\fontsize{8.85}{10.45}\selectfont
\begin{corebox}
\textbf{Core statement.} Pollen is typed data moving between nodes.
\end{corebox}
It includes task packets, prompts, context fragments, embeddings, model proposals, validator results, evidence references, routing messages, state deltas, and summaries. Pollen has origin, destination, type, freshness, scope, and expiration. It may be fertile, sterile, stale, rejected, or quarantined.

\textbf{Garden consequence.} This framing shifts attention from static context size to the quality and timing of transfers. More pollen is not necessarily better; duplicate or stale pollen can exhaust leaves without producing fruit.

\textbf{Operational rules}
\begin{itemize}
\item Type every packet.
\item Record provenance and read head.
\item Bound size and lifetime.
\item Deduplicate repeated context.
\end{itemize}
\textbf{Extended analysis.} Systems interpretation. The biological analogy for Pollen is useful only because it separates a real machine responsibility, flow, or lifecycle condition. The analogy must stop where biology would obscure authority or measurement. Every Garden term therefore needs fields, states, owners, and evidence. The purpose is not to claim that software is alive; it is to borrow the proven logic of finite organisms, ecological succession, and material recycling for systems that must remain useful across many hardware and project cycles.

\textbf{Failure signatures}
\begin{itemize}
\item Pollen appears in diagrams but has no machine-readable state.
\item Activity continues without an exit condition or fruit target.
\item Retired material is either kept forever or deleted without provenance.
\item The metaphor is used to bypass the existing LeafOS authority doctrine.
\end{itemize}
\textbf{Implementation or verification note.} Design consequence. A runtime that adopts pollen should expose a measurable indicator, a transition rule, and a failure disposition. Without those three elements the term remains presentation. With them, it becomes a compact systems language that can guide packs, schedulers, interfaces, and long-term maintenance.

\begin{acceptbox}
\textbf{Acceptance statement.} Pollen passes when every consumed packet can be traced to an origin and lifecycle disposition.
\end{acceptbox}
\newpage
\platetitle{84}{Leaves}{FIFTH IV}
\prov{[G] New philosophy}
\fontsize{8.85}{10.45}\selectfont
\begin{corebox}
\textbf{Core statement.} Leaves are the most complex brain-like transformation systems in the Garden.
\end{corebox}
They integrate signals, context, evidence, memory, constraints, and goals into plans, hypotheses, decisions, or synthesized outputs. A leaf may be a persona-bearing executive model, a multimodal reasoner, a router with learned interpretation, or another high-context component.

\textbf{Garden consequence.} A leaf is not authority. Its complexity makes it valuable and fallible. Native tools still prove, the inlet still disposes capabilities, and checkpoints still define durable continuation.

\textbf{Operational rules}
\begin{itemize}
\item Bind each leaf to a contract and cursor.
\item Record consumed pollen.
\item Bound output schemas.
\item Keep cognition separable from execution.
\end{itemize}
\textbf{Extended analysis.} Systems interpretation. The biological analogy for Leaves is useful only because it separates a real machine responsibility, flow, or lifecycle condition. The analogy must stop where biology would obscure authority or measurement. Every Garden term therefore needs fields, states, owners, and evidence. The purpose is not to claim that software is alive; it is to borrow the proven logic of finite organisms, ecological succession, and material recycling for systems that must remain useful across many hardware and project cycles.

\textbf{Failure signatures}
\begin{itemize}
\item Leaves appears in diagrams but has no machine-readable state.
\item Activity continues without an exit condition or fruit target.
\item Retired material is either kept forever or deleted without provenance.
\item The metaphor is used to bypass the existing LeafOS authority doctrine.
\end{itemize}
\textbf{Implementation or verification note.} Design consequence. A runtime that adopts leaves should expose a measurable indicator, a transition rule, and a failure disposition. Without those three elements the term remains presentation. With them, it becomes a compact systems language that can guide packs, schedulers, interfaces, and long-term maintenance.

\begin{acceptbox}
\textbf{Acceptance statement.} Leaves pass when replacing one model changes capability and latency without changing legal authority.
\end{acceptbox}
\newpage
\platetitle{85}{Bloom}{FIFTH IV}
\prov{[G] New philosophy}
\fontsize{8.85}{10.45}\selectfont
\begin{corebox}
\textbf{Core statement.} Bloom is the active coordinated runtime phase.
\end{corebox}
During bloom, selected petals and leaves are open, pollen moves, the stem schedules work, roots supply power, and the system pursues a declared fruit target. Bloom has an activation time, budget, active profile, acceptance gates, and exit conditions.

\textbf{Garden consequence.} Continual Bloom becomes a specific operating mode in which synchronization persists across repeated micro-cycles. It is no longer the definition of the whole system.

\textbf{Operational rules}
\begin{itemize}
\item Declare a fruit target.
\item Declare maximum duration or budget.
\item Report active nodes and health.
\item Exit on success, failure, block, or root pressure.
\end{itemize}
\textbf{Extended analysis.} Systems interpretation. The biological analogy for Bloom is useful only because it separates a real machine responsibility, flow, or lifecycle condition. The analogy must stop where biology would obscure authority or measurement. Every Garden term therefore needs fields, states, owners, and evidence. The purpose is not to claim that software is alive; it is to borrow the proven logic of finite organisms, ecological succession, and material recycling for systems that must remain useful across many hardware and project cycles.

\textbf{Failure signatures}
\begin{itemize}
\item Bloom appears in diagrams but has no machine-readable state.
\item Activity continues without an exit condition or fruit target.
\item Retired material is either kept forever or deleted without provenance.
\item The metaphor is used to bypass the existing LeafOS authority doctrine.
\end{itemize}
\textbf{Implementation or verification note.} Design consequence. A runtime that adopts bloom should expose a measurable indicator, a transition rule, and a failure disposition. Without those three elements the term remains presentation. With them, it becomes a compact systems language that can guide packs, schedulers, interfaces, and long-term maintenance.

\begin{acceptbox}
\textbf{Acceptance statement.} Bloom passes when its start and end can be reconstructed and its resource cost can be attributed.
\end{acceptbox}
\newpage
\platetitle{86}{Fruit}{FIFTH IV}
\prov{[G] New philosophy}
\fontsize{8.85}{10.45}\selectfont
\begin{corebox}
\textbf{Core statement.} Fruit is durable, validated, reusable output.
\end{corebox}
Examples include accepted patches, tested tools, reports, benchmark portfolios, pack manifests, decisions, models of a project, and checkpointed next actions. Fruit is not merely whatever a model emitted. It is the portion of a bloom that survives validation and remains useful after active compute closes.

\textbf{Garden consequence.} Fruit should be indexed for future seeds and linked to the evidence that justified it.

\textbf{Operational rules}
\begin{itemize}
\item Require validation appropriate to artifact type.
\item Store provenance and digest.
\item Separate accepted fruit from candidate output.
\item Define retention policy.
\end{itemize}
\textbf{Extended analysis.} Systems interpretation. The biological analogy for Fruit is useful only because it separates a real machine responsibility, flow, or lifecycle condition. The analogy must stop where biology would obscure authority or measurement. Every Garden term therefore needs fields, states, owners, and evidence. The purpose is not to claim that software is alive; it is to borrow the proven logic of finite organisms, ecological succession, and material recycling for systems that must remain useful across many hardware and project cycles.

\textbf{Failure signatures}
\begin{itemize}
\item Fruit appears in diagrams but has no machine-readable state.
\item Activity continues without an exit condition or fruit target.
\item Retired material is either kept forever or deleted without provenance.
\item The metaphor is used to bypass the existing LeafOS authority doctrine.
\end{itemize}
\textbf{Implementation or verification note.} Design consequence. A runtime that adopts fruit should expose a measurable indicator, a transition rule, and a failure disposition. Without those three elements the term remains presentation. With them, it becomes a compact systems language that can guide packs, schedulers, interfaces, and long-term maintenance.

\begin{acceptbox}
\textbf{Acceptance statement.} Fruit passes when a future run can reuse it without reconstructing the entire bloom.
\end{acceptbox}
\newpage
\platetitle{87}{Seeds}{FIFTH IV}
\prov{[G] New philosophy}
\fontsize{8.85}{10.45}\selectfont
\begin{corebox}
\textbf{Core statement.} Seeds are compact, restartable packages of intent and proven context.
\end{corebox}
A seed may contain project identity, objective, constraints, selected fruit references, active hypotheses, known blockers, preferred pack, and one next action. It is smaller than the full historical transcript and richer than a prompt.

\textbf{Garden consequence.} Seeds allow the Garden to grow again after all active petals and leaves have closed. They are the durable interface between cycles.

\textbf{Operational rules}
\begin{itemize}
\item Include only current, traceable state.
\item Reference fruit and evidence.
\item Exclude stale rejected branches unless explicitly relevant.
\item Keep one next action.
\end{itemize}
\textbf{Extended analysis.} Systems interpretation. The biological analogy for Seeds is useful only because it separates a real machine responsibility, flow, or lifecycle condition. The analogy must stop where biology would obscure authority or measurement. Every Garden term therefore needs fields, states, owners, and evidence. The purpose is not to claim that software is alive; it is to borrow the proven logic of finite organisms, ecological succession, and material recycling for systems that must remain useful across many hardware and project cycles.

\textbf{Failure signatures}
\begin{itemize}
\item Seeds appears in diagrams but has no machine-readable state.
\item Activity continues without an exit condition or fruit target.
\item Retired material is either kept forever or deleted without provenance.
\item The metaphor is used to bypass the existing LeafOS authority doctrine.
\end{itemize}
\textbf{Implementation or verification note.} Design consequence. A runtime that adopts seeds should expose a measurable indicator, a transition rule, and a failure disposition. Without those three elements the term remains presentation. With them, it becomes a compact systems language that can guide packs, schedulers, interfaces, and long-term maintenance.

\begin{acceptbox}
\textbf{Acceptance statement.} Seeds pass when a new bloom can start from them without reading every prior token.
\end{acceptbox}
\newpage
\platetitle{88}{Senescence}{FIFTH IV}
\prov{[G] New philosophy}
\fontsize{8.85}{10.45}\selectfont
\begin{corebox}
\textbf{Core statement.} Senescence is controlled reduction of active capability after fruit, exhaustion, block, or environmental pressure.
\end{corebox}
The system stops admitting new work, drains accepted tasks, checkpoints current facts, releases leases, flushes evidence, summarizes residue, and unloads selected models. Senescence is not abrupt failure; it is an orderly end phase.

\textbf{Garden consequence.} Many current systems omit this phase and therefore leak memory, workers, and state until the operator performs manual death.

\textbf{Operational rules}
\begin{itemize}
\item Stop admission before shutdown.
\item Drain bounded accepted work.
\item Checkpoint before unload.
\item Record what remains unresolved.
\end{itemize}
\textbf{Extended analysis.} Systems interpretation. The biological analogy for Senescence is useful only because it separates a real machine responsibility, flow, or lifecycle condition. The analogy must stop where biology would obscure authority or measurement. Every Garden term therefore needs fields, states, owners, and evidence. The purpose is not to claim that software is alive; it is to borrow the proven logic of finite organisms, ecological succession, and material recycling for systems that must remain useful across many hardware and project cycles.

\textbf{Failure signatures}
\begin{itemize}
\item Senescence appears in diagrams but has no machine-readable state.
\item Activity continues without an exit condition or fruit target.
\item Retired material is either kept forever or deleted without provenance.
\item The metaphor is used to bypass the existing LeafOS authority doctrine.
\end{itemize}
\textbf{Implementation or verification note.} Design consequence. A runtime that adopts senescence should expose a measurable indicator, a transition rule, and a failure disposition. Without those three elements the term remains presentation. With them, it becomes a compact systems language that can guide packs, schedulers, interfaces, and long-term maintenance.

\begin{acceptbox}
\textbf{Acceptance statement.} Senescence passes when the bloom ends without orphaned processes or ambiguous resume state.
\end{acceptbox}
\newpage
\platetitle{89}{Death}{FIFTH IV}
\prov{[G] New philosophy}
\fontsize{8.85}{10.45}\selectfont
\begin{corebox}
\textbf{Core statement.} Death is the termination of a specific process, node instance, strategy, pack profile, or workflow lineage.
\end{corebox}
Death may be successful, planned, failed, or quarantined. It does not imply loss of Garden identity because contracts, fruit, seeds, evidence, and compost survive independently. Some dead branches should never regrow without new evidence.

\textbf{Garden consequence.} Treating death explicitly prevents hidden zombie workers and endless retries masquerading as persistence.

\textbf{Operational rules}
\begin{itemize}
\item Record cause and final state.
\item Release all owned resources.
\item Preserve failure evidence.
\item Define whether regrowth is allowed.
\end{itemize}
\textbf{Extended analysis.} Systems interpretation. The biological analogy for Death is useful only because it separates a real machine responsibility, flow, or lifecycle condition. The analogy must stop where biology would obscure authority or measurement. Every Garden term therefore needs fields, states, owners, and evidence. The purpose is not to claim that software is alive; it is to borrow the proven logic of finite organisms, ecological succession, and material recycling for systems that must remain useful across many hardware and project cycles.

\textbf{Failure signatures}
\begin{itemize}
\item Death appears in diagrams but has no machine-readable state.
\item Activity continues without an exit condition or fruit target.
\item Retired material is either kept forever or deleted without provenance.
\item The metaphor is used to bypass the existing LeafOS authority doctrine.
\end{itemize}
\textbf{Implementation or verification note.} Design consequence. A runtime that adopts death should expose a measurable indicator, a transition rule, and a failure disposition. Without those three elements the term remains presentation. With them, it becomes a compact systems language that can guide packs, schedulers, interfaces, and long-term maintenance.

\begin{acceptbox}
\textbf{Acceptance statement.} Death passes when process absence is expected and recoverability remains measurable.
\end{acceptbox}
\newpage
\platetitle{90}{Compost}{FIFTH IV}
\prov{[G] New philosophy}
\fontsize{8.85}{10.45}\selectfont
\begin{corebox}
\textbf{Core statement.} Compost is retired state converted into bounded future value.
\end{corebox}
Inputs include failed attempts, rejected candidates, stale outputs, old logs, historical context, deactivated profiles, and superseded artifacts. Composting summarizes, classifies, deduplicates, preserves evidence references, and separates reusable lessons from archive-only bulk.

\textbf{Garden consequence.} Compost is not a euphemism for dumping everything into a vector database. It has provenance, retention, and contamination controls.

\textbf{Operational rules}
\begin{itemize}
\item Never promote failure to fact.
\item Retain source links.
\item Deduplicate repeated material.
\item Quarantine unsafe or contradictory residue.
\end{itemize}
\textbf{Extended analysis.} Systems interpretation. The biological analogy for Compost is useful only because it separates a real machine responsibility, flow, or lifecycle condition. The analogy must stop where biology would obscure authority or measurement. Every Garden term therefore needs fields, states, owners, and evidence. The purpose is not to claim that software is alive; it is to borrow the proven logic of finite organisms, ecological succession, and material recycling for systems that must remain useful across many hardware and project cycles.

\textbf{Failure signatures}
\begin{itemize}
\item Compost appears in diagrams but has no machine-readable state.
\item Activity continues without an exit condition or fruit target.
\item Retired material is either kept forever or deleted without provenance.
\item The metaphor is used to bypass the existing LeafOS authority doctrine.
\end{itemize}
\textbf{Implementation or verification note.} Design consequence. A runtime that adopts compost should expose a measurable indicator, a transition rule, and a failure disposition. Without those three elements the term remains presentation. With them, it becomes a compact systems language that can guide packs, schedulers, interfaces, and long-term maintenance.

\begin{acceptbox}
\textbf{Acceptance statement.} Compost passes when future context becomes smaller and more useful without losing the ability to audit important conclusions.
\end{acceptbox}
\newpage
\platetitle{91}{Pruning}{FIFTH IV}
\prov{[G] New philosophy}
\fontsize{8.85}{10.45}\selectfont
\begin{corebox}
\textbf{Core statement.} Pruning removes branches that consume resources without improving future fruit.
\end{corebox}
Targets may include duplicate model variants, stale context, dead routes, abandoned experiments, redundant summaries, unused caches, and policies disproven by evidence. Pruning is a decision with scope, evidence, reversibility, and retention class.

\textbf{Garden consequence.} Pruning is not deletion by mood. It is evidence-based simplification.

\textbf{Operational rules}
\begin{itemize}
\item Preview before destructive pruning.
\item Protect source, evidence, and committed fruit.
\item Archive when uncertainty is high.
\item Measure reclaimed resources.
\end{itemize}
\textbf{Extended analysis.} Systems interpretation. The biological analogy for Pruning is useful only because it separates a real machine responsibility, flow, or lifecycle condition. The analogy must stop where biology would obscure authority or measurement. Every Garden term therefore needs fields, states, owners, and evidence. The purpose is not to claim that software is alive; it is to borrow the proven logic of finite organisms, ecological succession, and material recycling for systems that must remain useful across many hardware and project cycles.

\textbf{Failure signatures}
\begin{itemize}
\item Pruning appears in diagrams but has no machine-readable state.
\item Activity continues without an exit condition or fruit target.
\item Retired material is either kept forever or deleted without provenance.
\item The metaphor is used to bypass the existing LeafOS authority doctrine.
\end{itemize}
\textbf{Implementation or verification note.} Design consequence. A runtime that adopts pruning should expose a measurable indicator, a transition rule, and a failure disposition. Without those three elements the term remains presentation. With them, it becomes a compact systems language that can guide packs, schedulers, interfaces, and long-term maintenance.

\begin{acceptbox}
\textbf{Acceptance statement.} Pruning passes when the system becomes simpler without losing current truth or recovery ability.
\end{acceptbox}
\newpage
\platetitle{92}{Dormancy}{FIFTH IV}
\prov{[G] New philosophy}
\fontsize{8.85}{10.45}\selectfont
\begin{corebox}
\textbf{Core statement.} Dormancy preserves identity and readiness without active compute.
\end{corebox}
A dormant model, persona, pack, or project retains catalog identity, contracts, seeds, and activation conditions while releasing VRAM, RAM, processes, and leases. Dormancy can be scheduled, resource-triggered, or operator-selected.

\textbf{Garden consequence.} This is the practical answer to large local model collections: own many possibilities, activate few.

\textbf{Operational rules}
\begin{itemize}
\item Persist activation prerequisites.
\item Record last fruit and last failure.
\item Keep no hidden heartbeat requirement.
\item Test cold regrowth.
\end{itemize}
\textbf{Extended analysis.} Systems interpretation. The biological analogy for Dormancy is useful only because it separates a real machine responsibility, flow, or lifecycle condition. The analogy must stop where biology would obscure authority or measurement. Every Garden term therefore needs fields, states, owners, and evidence. The purpose is not to claim that software is alive; it is to borrow the proven logic of finite organisms, ecological succession, and material recycling for systems that must remain useful across many hardware and project cycles.

\textbf{Failure signatures}
\begin{itemize}
\item Dormancy appears in diagrams but has no machine-readable state.
\item Activity continues without an exit condition or fruit target.
\item Retired material is either kept forever or deleted without provenance.
\item The metaphor is used to bypass the existing LeafOS authority doctrine.
\end{itemize}
\textbf{Implementation or verification note.} Design consequence. A runtime that adopts dormancy should expose a measurable indicator, a transition rule, and a failure disposition. Without those three elements the term remains presentation. With them, it becomes a compact systems language that can guide packs, schedulers, interfaces, and long-term maintenance.

\begin{acceptbox}
\textbf{Acceptance statement.} Dormancy passes when a component can remain inactive for an arbitrary period and return from durable state.
\end{acceptbox}
\newpage
\platetitle{93}{Seasons}{FIFTH IV}
\prov{[G] New philosophy}
\fontsize{8.85}{10.45}\selectfont
\begin{corebox}
\textbf{Core statement.} Seasons are policy periods that change which blooms are desirable.
\end{corebox}
A season may correspond to time of day, foreground activity, electricity price, thermal conditions, project phase, operator presence, or hardware availability. Quiet seasons favor small CPU leaves and maintenance. Work seasons favor default brain and coder. Deep seasons permit expensive peak models. Recovery seasons prioritize audit and compost.

\textbf{Garden consequence.} Seasonality replaces the assumption that one configuration should be optimal at all times.

\textbf{Operational rules}
\begin{itemize}
\item Define season triggers.
\item Allow manual override.
\item Keep authority constant across seasons.
\item Record transition reasons.
\end{itemize}
\textbf{Extended analysis.} Systems interpretation. The biological analogy for Seasons is useful only because it separates a real machine responsibility, flow, or lifecycle condition. The analogy must stop where biology would obscure authority or measurement. Every Garden term therefore needs fields, states, owners, and evidence. The purpose is not to claim that software is alive; it is to borrow the proven logic of finite organisms, ecological succession, and material recycling for systems that must remain useful across many hardware and project cycles.

\textbf{Failure signatures}
\begin{itemize}
\item Seasons appears in diagrams but has no machine-readable state.
\item Activity continues without an exit condition or fruit target.
\item Retired material is either kept forever or deleted without provenance.
\item The metaphor is used to bypass the existing LeafOS authority doctrine.
\end{itemize}
\textbf{Implementation or verification note.} Design consequence. A runtime that adopts seasons should expose a measurable indicator, a transition rule, and a failure disposition. Without those three elements the term remains presentation. With them, it becomes a compact systems language that can guide packs, schedulers, interfaces, and long-term maintenance.

\begin{acceptbox}
\textbf{Acceptance statement.} Season scheduling passes when the same task may select a different healthy profile under different environmental conditions.
\end{acceptbox}
\newpage
\platetitle{94}{Diversity}{FIFTH IV}
\prov{[G] New philosophy}
\fontsize{8.85}{10.45}\selectfont
\begin{corebox}
\textbf{Core statement.} A healthy Garden preserves functional diversity without keeping every variant active.
\end{corebox}
Different models, quantizations, tools, validators, personas, and hardware lanes provide resilience to task variation and failure. Diversity is catalog and capability breadth; residency is a separate budget decision.

\textbf{Garden consequence.} Monoculture can be efficient until its assumptions fail. Unbounded diversity becomes storage and routing chaos. The design needs measured niches and retirement rules.

\textbf{Operational rules}
\begin{itemize}
\item Record each component niche.
\item Benchmark overlapping candidates.
\item Retire dominated variants.
\item Keep at least one recovery path.
\end{itemize}
\textbf{Extended analysis.} Systems interpretation. The biological analogy for Diversity is useful only because it separates a real machine responsibility, flow, or lifecycle condition. The analogy must stop where biology would obscure authority or measurement. Every Garden term therefore needs fields, states, owners, and evidence. The purpose is not to claim that software is alive; it is to borrow the proven logic of finite organisms, ecological succession, and material recycling for systems that must remain useful across many hardware and project cycles.

\textbf{Failure signatures}
\begin{itemize}
\item Diversity appears in diagrams but has no machine-readable state.
\item Activity continues without an exit condition or fruit target.
\item Retired material is either kept forever or deleted without provenance.
\item The metaphor is used to bypass the existing LeafOS authority doctrine.
\end{itemize}
\textbf{Implementation or verification note.} Design consequence. A runtime that adopts diversity should expose a measurable indicator, a transition rule, and a failure disposition. Without those three elements the term remains presentation. With them, it becomes a compact systems language that can guide packs, schedulers, interfaces, and long-term maintenance.

\begin{acceptbox}
\textbf{Acceptance statement.} Diversity passes when alternatives exist for meaningful failure classes and routing can explain their selection.
\end{acceptbox}
\newpage
\platetitle{95}{Mutualism}{FIFTH IV}
\prov{[G] New philosophy}
\fontsize{8.85}{10.45}\selectfont
\begin{corebox}
\textbf{Core statement.} Nodes should exchange work when the exchange improves joint yield.
\end{corebox}
Examples include a small auxiliary compressing context for a larger leaf, a CPU validator checking GPU output, a writer polishing a reasoner draft, or one node sharing verified fruit with another. Mutualism is measured by reduced latency, waste, or error rather than by the number of messages exchanged.

\textbf{Garden consequence.} Pollen exchange without benefit is chatter, not ecology.

\textbf{Operational rules}
\begin{itemize}
\item Measure before and after exchange.
\item Prevent circular triggering.
\item Preserve source identity.
\item Use bounded contracts.
\end{itemize}
\textbf{Extended analysis.} Systems interpretation. The biological analogy for Mutualism is useful only because it separates a real machine responsibility, flow, or lifecycle condition. The analogy must stop where biology would obscure authority or measurement. Every Garden term therefore needs fields, states, owners, and evidence. The purpose is not to claim that software is alive; it is to borrow the proven logic of finite organisms, ecological succession, and material recycling for systems that must remain useful across many hardware and project cycles.

\textbf{Failure signatures}
\begin{itemize}
\item Mutualism appears in diagrams but has no machine-readable state.
\item Activity continues without an exit condition or fruit target.
\item Retired material is either kept forever or deleted without provenance.
\item The metaphor is used to bypass the existing LeafOS authority doctrine.
\end{itemize}
\textbf{Implementation or verification note.} Design consequence. A runtime that adopts mutualism should expose a measurable indicator, a transition rule, and a failure disposition. Without those three elements the term remains presentation. With them, it becomes a compact systems language that can guide packs, schedulers, interfaces, and long-term maintenance.

\begin{acceptbox}
\textbf{Acceptance statement.} Mutualism passes when cooperation improves validated yield or resource efficiency.
\end{acceptbox}
\newpage
\platetitle{96}{Disease Containment}{FIFTH IV}
\prov{[G] New philosophy}
\fontsize{8.85}{10.45}\selectfont
\begin{corebox}
\textbf{Core statement.} Hallucinations, corrupted state, runaway loops, invalid artifacts, and compromised tools are forms of systemic disease.
\end{corebox}
Containment uses schema validation, quarantine, capability boundaries, evidence requirements, stale rejection, process ownership, retry budgets, and isolated experiments. A diseased branch may be studied without being allowed to fertilize the main Garden.

\textbf{Garden consequence.} The metaphor becomes useful here because propagation matters: one bad packet can contaminate many downstream summaries if status is lost.

\textbf{Operational rules}
\begin{itemize}
\item Quarantine unsupported high-risk claims.
\item Stop recursive retry loops.
\item Preserve contamination provenance.
\item Require explicit promotion from experiments.
\end{itemize}
\textbf{Extended analysis.} Systems interpretation. The biological analogy for Disease Containment is useful only because it separates a real machine responsibility, flow, or lifecycle condition. The analogy must stop where biology would obscure authority or measurement. Every Garden term therefore needs fields, states, owners, and evidence. The purpose is not to claim that software is alive; it is to borrow the proven logic of finite organisms, ecological succession, and material recycling for systems that must remain useful across many hardware and project cycles.

\textbf{Failure signatures}
\begin{itemize}
\item Disease Containment appears in diagrams but has no machine-readable state.
\item Activity continues without an exit condition or fruit target.
\item Retired material is either kept forever or deleted without provenance.
\item The metaphor is used to bypass the existing LeafOS authority doctrine.
\end{itemize}
\textbf{Implementation or verification note.} Design consequence. A runtime that adopts disease containment should expose a measurable indicator, a transition rule, and a failure disposition. Without those three elements the term remains presentation. With them, it becomes a compact systems language that can guide packs, schedulers, interfaces, and long-term maintenance.

\begin{acceptbox}
\textbf{Acceptance statement.} Containment passes when one failed or malicious node cannot silently rewrite committed state.
\end{acceptbox}
\newpage
\platetitle{97}{Succession}{FIFTH IV}
\prov{[G] New philosophy}
\fontsize{8.85}{10.45}\selectfont
\begin{corebox}
\textbf{Core statement.} Succession is the replacement of one stable configuration by another after evidence or environmental change.
\end{corebox}
A model may be promoted, demoted, or retired. A pack profile may replace another. A project may move from exploration to delivery. Succession preserves fruit and seeds while allowing active structure to change.

\textbf{Garden consequence.} This is safer than in-place mutation of an immortal bloom because the old configuration can remain archived and comparable.

\textbf{Operational rules}
\begin{itemize}
\item Define promotion evidence.
\item Preserve predecessor identity.
\item Support rollback when safe.
\item Update routing only after acceptance.
\end{itemize}
\textbf{Extended analysis.} Systems interpretation. The biological analogy for Succession is useful only because it separates a real machine responsibility, flow, or lifecycle condition. The analogy must stop where biology would obscure authority or measurement. Every Garden term therefore needs fields, states, owners, and evidence. The purpose is not to claim that software is alive; it is to borrow the proven logic of finite organisms, ecological succession, and material recycling for systems that must remain useful across many hardware and project cycles.

\textbf{Failure signatures}
\begin{itemize}
\item Succession appears in diagrams but has no machine-readable state.
\item Activity continues without an exit condition or fruit target.
\item Retired material is either kept forever or deleted without provenance.
\item The metaphor is used to bypass the existing LeafOS authority doctrine.
\end{itemize}
\textbf{Implementation or verification note.} Design consequence. A runtime that adopts succession should expose a measurable indicator, a transition rule, and a failure disposition. Without those three elements the term remains presentation. With them, it becomes a compact systems language that can guide packs, schedulers, interfaces, and long-term maintenance.

\begin{acceptbox}
\textbf{Acceptance statement.} Succession passes when a new configuration becomes default without erasing the evidence that justified the change.
\end{acceptbox}
\newpage
\platetitle{98}{Infinite Renewal, Not Infinite Size}{FIFTH IV}
\prov{[G] New philosophy}
\fontsize{8.85}{10.45}\selectfont
\begin{corebox}
\textbf{Core statement.} The Garden seeks indefinite renewal rather than indefinite resident growth.
\end{corebox}
Long-term capability improves when cycles produce fruit, compost useful residue, prune waste, preserve diversity, and regenerate from stronger seeds. Size may grow when justified, but growth is not measured by RAM, VRAM, model count, or process uptime alone.

\textbf{Garden consequence.} This is the formal replacement for Infinite Bloom.

\textbf{Operational rules}
\begin{itemize}
\item Optimize across cycles.
\item Treat resource growth as a costed decision.
\item Reward recoverability and reuse.
\item Permit contraction after success.
\end{itemize}
\textbf{Extended analysis.} Systems interpretation. The biological analogy for Infinite Renewal, Not Infinite Size is useful only because it separates a real machine responsibility, flow, or lifecycle condition. The analogy must stop where biology would obscure authority or measurement. Every Garden term therefore needs fields, states, owners, and evidence. The purpose is not to claim that software is alive; it is to borrow the proven logic of finite organisms, ecological succession, and material recycling for systems that must remain useful across many hardware and project cycles.

\textbf{Failure signatures}
\begin{itemize}
\item Infinite Renewal, Not Infinite Size appears in diagrams but has no machine-readable state.
\item Activity continues without an exit condition or fruit target.
\item Retired material is either kept forever or deleted without provenance.
\item The metaphor is used to bypass the existing LeafOS authority doctrine.
\end{itemize}
\textbf{Implementation or verification note.} Design consequence. A runtime that adopts infinite renewal, not infinite size should expose a measurable indicator, a transition rule, and a failure disposition. Without those three elements the term remains presentation. With them, it becomes a compact systems language that can guide packs, schedulers, interfaces, and long-term maintenance.

\begin{acceptbox}
\textbf{Acceptance statement.} The principle passes when the architecture can improve while its active memory footprint stays flat or decreases.
\end{acceptbox}
\newpage
\platetitle{99}{Garden Graph Model}{FIFTH IV}
\prov{[G] New philosophy}
\fontsize{8.85}{10.45}\selectfont
\begin{corebox}
\textbf{Core statement.} The Garden can be modeled as typed nodes, typed flows, durable stores, and lifecycle states.
\end{corebox}
Let roots supply energy edges, the stem supply control and transport edges, petals supply compute capacity, pollen represent data flows, leaves perform high-complexity transformations, blooms represent active subgraphs, fruit represent committed outputs, and compost represent compressed retired subgraphs.

\textbf{Garden consequence.} The model separates three flows: power, data, and control. Conflating them is a common source of unsafe architecture.

\textbf{Operational rules}
\begin{itemize}
\item Type nodes and edges.
\item Record time-varying activation.
\item Keep power facts distinct from task decisions.
\item Keep data flow distinct from control authority.
\end{itemize}
\textbf{Extended analysis.} Systems interpretation. The biological analogy for Garden Graph Model is useful only because it separates a real machine responsibility, flow, or lifecycle condition. The analogy must stop where biology would obscure authority or measurement. Every Garden term therefore needs fields, states, owners, and evidence. The purpose is not to claim that software is alive; it is to borrow the proven logic of finite organisms, ecological succession, and material recycling for systems that must remain useful across many hardware and project cycles.

\textbf{Failure signatures}
\begin{itemize}
\item Garden Graph Model appears in diagrams but has no machine-readable state.
\item Activity continues without an exit condition or fruit target.
\item Retired material is either kept forever or deleted without provenance.
\item The metaphor is used to bypass the existing LeafOS authority doctrine.
\end{itemize}
\textbf{Implementation or verification note.} Design consequence. A runtime that adopts garden graph model should expose a measurable indicator, a transition rule, and a failure disposition. Without those three elements the term remains presentation. With them, it becomes a compact systems language that can guide packs, schedulers, interfaces, and long-term maintenance.

\begin{lstlisting}
G(t) = (S, R, T, P, D, L, B, F, C)
\end{lstlisting}
\begin{acceptbox}
\textbf{Acceptance statement.} The graph passes when a run can be reconstructed as an active subgraph with attributable energy, data, decisions, and fruit.
\end{acceptbox}
\newpage
\platetitle{100}{Fourth-Fifth Synthesis}{FIFTH IV}
\prov{[G] New philosophy}
\fontsize{8.85}{10.45}\selectfont
\begin{corebox}
\textbf{Core statement.} The fourth fifth replaces permanent bloom with a complete lifecycle vocabulary.
\end{corebox}
The Garden contains soil, roots, stem, petals, pollen, leaves, bloom, fruit, seeds, senescence, death, compost, pruning, dormancy, seasons, diversity, mutualism, disease containment, and succession. None of these terms is allowed to remain merely poetic; each maps to runtime state, data, resources, or authority.

\textbf{Garden consequence.} The philosophy is conservative in one sense and radical in another. It preserves LeafOS proof and state contracts while allowing any active model, process, or profile to die.

\textbf{Operational rules}
\begin{itemize}
\item Identity survives process death.
\item Fruit survives bloom.
\item Compost improves seeds.
\item Growth is evaluated across renewal cycles.
\end{itemize}
\textbf{Extended analysis.} Systems interpretation. The biological analogy for Fourth-Fifth Synthesis is useful only because it separates a real machine responsibility, flow, or lifecycle condition. The analogy must stop where biology would obscure authority or measurement. Every Garden term therefore needs fields, states, owners, and evidence. The purpose is not to claim that software is alive; it is to borrow the proven logic of finite organisms, ecological succession, and material recycling for systems that must remain useful across many hardware and project cycles.

\textbf{Failure signatures}
\begin{itemize}
\item Fourth-Fifth Synthesis appears in diagrams but has no machine-readable state.
\item Activity continues without an exit condition or fruit target.
\item Retired material is either kept forever or deleted without provenance.
\item The metaphor is used to bypass the existing LeafOS authority doctrine.
\end{itemize}
\textbf{Implementation or verification note.} Design consequence. A runtime that adopts fourth-fifth synthesis should expose a measurable indicator, a transition rule, and a failure disposition. Without those three elements the term remains presentation. With them, it becomes a compact systems language that can guide packs, schedulers, interfaces, and long-term maintenance.

\begin{acceptbox}
\textbf{Acceptance statement.} The fourth fifth is complete when Infinite Bloom is understood as one useful phase inside a larger, healthier system.
\end{acceptbox}
\newpage
\platetitle{101}{Garden Lifecycle State Machine}{FIFTH V}
\prov{[G] New implementation}
\fontsize{8.85}{10.45}\selectfont
\begin{corebox}
\textbf{Core statement.} Every Garden entity moves through explicit lifecycle states.
\end{corebox}
A general state machine uses dormant, seeded, rooted, sprouting, blooming, fruiting, senescent, composting, dead, blocked, and failed. Not every entity uses every state, but legal transitions are declared by kind. The supervisor owns transitions and records reasons, budgets, and evidence.

\textbf{Garden consequence.} This state machine prevents a model or workflow from remaining ambiguously half-alive after failure.

\textbf{Operational rules}
\begin{itemize}
\item Define allowed transitions per kind.
\item Attach entry and exit criteria.
\item Record transition actor and evidence.
\item Provide recovery or terminal disposition.
\end{itemize}
\textbf{Extended analysis.} Implementation analysis. Garden Lifecycle State Machine should enter the runtime as a narrow schema, policy, state transition, or projection that composes with the existing inlet, journal, validator, checkpoint, and resource governor. The first implementation should prefer deterministic fields and explicit defaults over generalized agent behavior. Compatibility matters: pre-Garden runs must remain readable, and Garden fields must not imply capabilities that are not implemented.

\textbf{Failure signatures}
\begin{itemize}
\item An illegal lifecycle transition is accepted.
\item A node lacks owner, budget, closure condition, or evidence path.
\item Automation is enabled before preview, rollback, and recovery are proven.
\item A new Garden subsystem creates a second authority path around ProjectLeaf.
\end{itemize}
\textbf{Implementation or verification note.} Verification note. Build a fixture for garden lifecycle state machine, validate both a successful path and at least one refusal path, interrupt it mid-transition, and confirm that replay produces the same disposition. Only then should the feature be connected to automatic routing or lifecycle control.

\begin{lstlisting}
dormant -> seeded -> rooted -> sprouting -> blooming -> fruiting
   ^                                                   |
   |                                                   v
   +-------- composting <- senescent <-----------------+
                         |
                         +-> dead | blocked | failed
\end{lstlisting}
\begin{acceptbox}
\textbf{Acceptance statement.} The lifecycle passes when illegal jumps are rejected and a crash can be mapped to a recoverable state.
\end{acceptbox}
\newpage
\platetitle{102}{Seed Intent Contract}{FIFTH V}
\prov{[G] New implementation}
\fontsize{8.85}{10.45}\selectfont
\begin{corebox}
\textbf{Core statement.} A seed is a compact contract for starting or restarting work.
\end{corebox}
It records project identity, objective, constraints, selected fruit references, active hypotheses, blockers, preferred pack or profile, operator policy, read head, and one next action. It excludes raw historical bulk and unverified claims that are not needed for the next cycle.

\textbf{Garden consequence.} Seeds are the unit of regenerative continuity. They allow a bloom to end completely without losing direction.

\textbf{Operational rules}
\begin{itemize}
\item Validate seed schema.
\item Reference evidence by digest.
\item Keep scope bounded.
\item Reject stale seeds when project identity changes.
\end{itemize}
\textbf{Extended analysis.} Implementation analysis. Seed Intent Contract should enter the runtime as a narrow schema, policy, state transition, or projection that composes with the existing inlet, journal, validator, checkpoint, and resource governor. The first implementation should prefer deterministic fields and explicit defaults over generalized agent behavior. Compatibility matters: pre-Garden runs must remain readable, and Garden fields must not imply capabilities that are not implemented.

\textbf{Failure signatures}
\begin{itemize}
\item An illegal lifecycle transition is accepted.
\item A node lacks owner, budget, closure condition, or evidence path.
\item Automation is enabled before preview, rollback, and recovery are proven.
\item A new Garden subsystem creates a second authority path around ProjectLeaf.
\end{itemize}
\textbf{Implementation or verification note.} Verification note. Build a fixture for seed intent contract, validate both a successful path and at least one refusal path, interrupt it mid-transition, and confirm that replay produces the same disposition. Only then should the feature be connected to automatic routing or lifecycle control.

\begin{lstlisting}
{
  "kind": "garden.seed",
  "id": "seed-project-architecture",
  "project": {"root": "C:/R/LeafOS", "revision": "pinned"},
  "objective": "Produce one validated Garden architecture artifact",
  "constraints": ["local-first", "native-validation", "bounded-memory"],
  "fruit_refs": ["fruit:authority-atlas", "fruit:throughput-baseline"],
  "active_hypotheses": ["smaller cyclic packs improve useful yield"],
  "blockers": [],
  "preferred_pack": "garden-2x2",
  "read_head": "sha256:current-head",
  "next_action": "materialize an offline execution plan"
}
\end{lstlisting}
\begin{acceptbox}
\textbf{Acceptance statement.} The seed contract passes when a cold run can begin from it and reconstruct all required context.
\end{acceptbox}
\newpage
\platetitle{103}{Root Telemetry Contract}{FIFTH V}
\prov{[G] New implementation}
\fontsize{8.85}{10.45}\selectfont
\begin{corebox}
\textbf{Core statement.} Root telemetry describes whether power and thermal infrastructure can support a bloom.
\end{corebox}
Fields include source, sample time, sample age, CPU package power, GPU board power, temperatures, fan or cooling state when available, power limit, UPS state, missing fields, and health classification. The contract distinguishes unavailable, stale, warning, critical, and healthy.

\textbf{Garden consequence.} The lifecycle governor consumes this contract but cannot invent missing measurements.

\textbf{Operational rules}
\begin{itemize}
\item Use monotonic sample times where possible.
\item Record sensor source.
\item Apply hysteresis.
\item Checkpoint before critical shutdown.
\end{itemize}
\textbf{Extended analysis.} Implementation analysis. Root Telemetry Contract should enter the runtime as a narrow schema, policy, state transition, or projection that composes with the existing inlet, journal, validator, checkpoint, and resource governor. The first implementation should prefer deterministic fields and explicit defaults over generalized agent behavior. Compatibility matters: pre-Garden runs must remain readable, and Garden fields must not imply capabilities that are not implemented.

\textbf{Failure signatures}
\begin{itemize}
\item An illegal lifecycle transition is accepted.
\item A node lacks owner, budget, closure condition, or evidence path.
\item Automation is enabled before preview, rollback, and recovery are proven.
\item A new Garden subsystem creates a second authority path around ProjectLeaf.
\end{itemize}
\textbf{Implementation or verification note.} Verification note. Build a fixture for root telemetry contract, validate both a successful path and at least one refusal path, interrupt it mid-transition, and confirm that replay produces the same disposition. Only then should the feature be connected to automatic routing or lifecycle control.

\begin{lstlisting}
{
  "kind": "garden.root.sample",
  "source": "native-hardware-sampler",
  "sampled_at": "2026-07-31T00:23:00-07:00",
  "sample_age_ms": 140,
  "cpu": {"package_w": 72.4, "temperature_c": 78.0},
  "gpu": {"board_w": 146.2, "temperature_c": 66.0},
  "cooling": {"state": "normal", "telemetry_complete": true},
  "ups": {"available": false, "state": "unavailable"},
  "health": "healthy",
  "missing_fields": []
}
\end{lstlisting}
\begin{acceptbox}
\textbf{Acceptance statement.} Root telemetry passes when stale sensors cause a safe policy response rather than a fabricated healthy value.
\end{acceptbox}
\newpage
\platetitle{104}{Petal Node Contract}{FIFTH V}
\prov{[G] New implementation}
\fontsize{8.85}{10.45}\selectfont
\begin{corebox}
\textbf{Core statement.} A petal contract describes a bounded compute surface.
\end{corebox}
It records node identity, device class, provider, model or executable, role, owner, requested and actual placement, memory budget, power budget, context limit, concurrency, activation trigger, idle timeout, closure condition, and output channel.

\textbf{Garden consequence.} The contract makes compute transience explicit. Opening and closing are state transitions, not incidental side effects.

\textbf{Operational rules}
\begin{itemize}
\item One owner per active petal.
\item Report actual placement.
\item Enforce memory budget.
\item Capture exit and cleanup.
\end{itemize}
\textbf{Extended analysis.} Implementation analysis. Petal Node Contract should enter the runtime as a narrow schema, policy, state transition, or projection that composes with the existing inlet, journal, validator, checkpoint, and resource governor. The first implementation should prefer deterministic fields and explicit defaults over generalized agent behavior. Compatibility matters: pre-Garden runs must remain readable, and Garden fields must not imply capabilities that are not implemented.

\textbf{Failure signatures}
\begin{itemize}
\item An illegal lifecycle transition is accepted.
\item A node lacks owner, budget, closure condition, or evidence path.
\item Automation is enabled before preview, rollback, and recovery are proven.
\item A new Garden subsystem creates a second authority path around ProjectLeaf.
\end{itemize}
\textbf{Implementation or verification note.} Verification note. Build a fixture for petal node contract, validate both a successful path and at least one refusal path, interrupt it mid-transition, and confirm that replay produces the same disposition. Only then should the feature be connected to automatic routing or lifecycle control.

\begin{lstlisting}
{
  "kind": "garden.petal",
  "id": "petal-gpu-brain-01",
  "device_class": "gpu",
  "provider": "llama.cpp-vulkan",
  "role": "default_brain",
  "owner": "bloom:architecture-001",
  "placement": {"requested": "gpu", "actual": "gpu", "layers": "all"},
  "budgets": {"vram_gib": 10.5, "power_w": 180, "context_tokens": 32768},
  "concurrency": 1,
  "activation_trigger": "route.requires.default_brain",
  "idle_timeout_s": 180,
  "closure_condition": "fruit_committed_or_idle_timeout"
}
\end{lstlisting}
\begin{acceptbox}
\textbf{Acceptance statement.} The petal contract passes when the supervisor can stop one worker without corrupting the rest of the bloom.
\end{acceptbox}
\newpage
\platetitle{105}{Pollen Packet Contract}{FIFTH V}
\prov{[G] New implementation}
\fontsize{8.85}{10.45}\selectfont
\begin{corebox}
\textbf{Core statement.} Pollen packets are typed, bounded, fresh data transfers.
\end{corebox}
A packet records identifier, kind, origin node, destination or audience, run, sequence, read head, timestamp, expiration, payload schema, digest, evidence references, confidentiality class, and lifecycle disposition. Packets may be accepted, rejected, stale, duplicate, quarantined, or consumed.

\textbf{Garden consequence.} This contract replaces informal prompt passing with inspectable data movement.

\textbf{Operational rules}
\begin{itemize}
\item Validate packet kind.
\item Bound payload size.
\item Deduplicate by digest and context.
\item Reject expired or stale packets.
\end{itemize}
\textbf{Extended analysis.} Implementation analysis. Pollen Packet Contract should enter the runtime as a narrow schema, policy, state transition, or projection that composes with the existing inlet, journal, validator, checkpoint, and resource governor. The first implementation should prefer deterministic fields and explicit defaults over generalized agent behavior. Compatibility matters: pre-Garden runs must remain readable, and Garden fields must not imply capabilities that are not implemented.

\textbf{Failure signatures}
\begin{itemize}
\item An illegal lifecycle transition is accepted.
\item A node lacks owner, budget, closure condition, or evidence path.
\item Automation is enabled before preview, rollback, and recovery are proven.
\item A new Garden subsystem creates a second authority path around ProjectLeaf.
\end{itemize}
\textbf{Implementation or verification note.} Verification note. Build a fixture for pollen packet contract, validate both a successful path and at least one refusal path, interrupt it mid-transition, and confirm that replay produces the same disposition. Only then should the feature be connected to automatic routing or lifecycle control.

\begin{lstlisting}
{
  "kind": "garden.pollen",
  "id": "pollen-000184",
  "class": "proposal.architecture",
  "origin": "leaf:wednesday",
  "destination": ["leaf:friday", "journal"],
  "run_id": "bloom:architecture-001",
  "sequence": 184,
  "read_head": "sha256:head-183",
  "expires_at": "2026-07-31T01:00:00-07:00",
  "payload_ref": "artifact:proposal-184.json",
  "digest": "sha256:payload",
  "evidence_refs": [],
  "status": "available"
}
\end{lstlisting}
\begin{acceptbox}
\textbf{Acceptance statement.} The packet contract passes when the system can trace every leaf decision to the pollen it consumed.
\end{acceptbox}
\newpage
\platetitle{106}{Leaf Processor Contract}{FIFTH V}
\prov{[G] New implementation}
\fontsize{8.85}{10.45}\selectfont
\begin{corebox}
\textbf{Core statement.} A leaf contract defines a high-complexity transformation system.
\end{corebox}
It combines identity, duties, forbidden actions, model group, context policy, transcript cursor, accepted pollen kinds, output schema, resource preferences, retry budget, freshness requirement, and validator dependencies. MWF personalities are specialized leaf contracts.

\textbf{Garden consequence.} The leaf remains weaker than native authority regardless of model scale.

\textbf{Operational rules}
\begin{itemize}
\item Require current transcript head.
\item Bound working state.
\item Separate style from duties.
\item Declare output claims and evidence needs.
\end{itemize}
\textbf{Extended analysis.} Implementation analysis. Leaf Processor Contract should enter the runtime as a narrow schema, policy, state transition, or projection that composes with the existing inlet, journal, validator, checkpoint, and resource governor. The first implementation should prefer deterministic fields and explicit defaults over generalized agent behavior. Compatibility matters: pre-Garden runs must remain readable, and Garden fields must not imply capabilities that are not implemented.

\textbf{Failure signatures}
\begin{itemize}
\item An illegal lifecycle transition is accepted.
\item A node lacks owner, budget, closure condition, or evidence path.
\item Automation is enabled before preview, rollback, and recovery are proven.
\item A new Garden subsystem creates a second authority path around ProjectLeaf.
\end{itemize}
\textbf{Implementation or verification note.} Verification note. Build a fixture for leaf processor contract, validate both a successful path and at least one refusal path, interrupt it mid-transition, and confirm that replay produces the same disposition. Only then should the feature be connected to automatic routing or lifecycle control.

\begin{lstlisting}
{
  "kind": "garden.leaf",
  "id": "leaf-friday-delivery",
  "identity": "Friday - Research and Delivery",
  "duties": ["compare", "select", "define-acceptance", "deliver"],
  "forbidden": ["self-approve", "bypass-validator", "accept-stale-pollen"],
  "model_group": "garden.brain.default",
  "cursor": {"transcript_head": "sha256:current-head", "sequence": 184},
  "accepts": ["proposal.*", "evidence.*", "constraint.*"],
  "output_schema": "decision.delivery",
  "retry_budget": 1,
  "freshness_required": true
}
\end{lstlisting}
\begin{acceptbox}
\textbf{Acceptance statement.} The leaf contract passes when its model can be replaced without altering permissions.
\end{acceptbox}
\newpage
\platetitle{107}{Bloom Run Record}{FIFTH V}
\prov{[G] New implementation}
\fontsize{8.85}{10.45}\selectfont
\begin{corebox}
\textbf{Core statement.} A bloom run is an active subgraph with a declared purpose and budget.
\end{corebox}
The run record contains seed reference, selected pack and profile, active roots, petals, leaves, fruit target, start time, maximum duration, memory and energy budgets, acceptance gates, current phase, queue state, and termination reason. It references the journal rather than duplicating it.

\textbf{Garden consequence.} This gives each period of active intelligence a bounded identity.

\textbf{Operational rules}
\begin{itemize}
\item One bloom ID per active coordinated run.
\item Declare fruit target before activation.
\item Record selected nodes.
\item Close with a terminal reason.
\end{itemize}
\textbf{Extended analysis.} Implementation analysis. Bloom Run Record should enter the runtime as a narrow schema, policy, state transition, or projection that composes with the existing inlet, journal, validator, checkpoint, and resource governor. The first implementation should prefer deterministic fields and explicit defaults over generalized agent behavior. Compatibility matters: pre-Garden runs must remain readable, and Garden fields must not imply capabilities that are not implemented.

\textbf{Failure signatures}
\begin{itemize}
\item An illegal lifecycle transition is accepted.
\item A node lacks owner, budget, closure condition, or evidence path.
\item Automation is enabled before preview, rollback, and recovery are proven.
\item A new Garden subsystem creates a second authority path around ProjectLeaf.
\end{itemize}
\textbf{Implementation or verification note.} Verification note. Build a fixture for bloom run record, validate both a successful path and at least one refusal path, interrupt it mid-transition, and confirm that replay produces the same disposition. Only then should the feature be connected to automatic routing or lifecycle control.

\begin{lstlisting}
{
  "kind": "garden.bloom",
  "id": "bloom:architecture-001",
  "seed_ref": "seed-project-architecture",
  "pack": {"id": "garden-2x2", "profile": "small-core"},
  "phase": "blooming",
  "active_petals": ["petal-cpu-aux", "petal-gpu-brain-01"],
  "active_leaves": ["leaf-monday", "leaf-friday-delivery"],
  "fruit_target": "validated-architecture-document",
  "budgets": {"duration_s": 7200, "ram_gib": 28, "vram_gib": 11, "energy_wh": 420},
  "acceptance_gates": ["render-clean", "source-provenance", "checkpoint-commit"],
  "termination_reason": null
}
\end{lstlisting}
\begin{acceptbox}
\textbf{Acceptance statement.} The run record passes when resource use and artifacts can be attributed to one bloom.
\end{acceptbox}
\newpage
\platetitle{108}{Fruit Artifact Contract}{FIFTH V}
\prov{[G] New implementation}
\fontsize{8.85}{10.45}\selectfont
\begin{corebox}
\textbf{Core statement.} A fruit record wraps a durable artifact with validation and provenance.
\end{corebox}
It includes artifact type, path or object reference, digest, producing bloom, source events, evidence references, acceptance gates, accepted by, timestamp, retention class, reuse tags, and supersession state. Candidate outputs remain separate until acceptance.

\textbf{Garden consequence.} Fruit records form a reusable index across cycles.

\textbf{Operational rules}
\begin{itemize}
\item Digest artifacts.
\item Link validation.
\item Record supersession.
\item Keep retention explicit.
\end{itemize}
\textbf{Extended analysis.} Implementation analysis. Fruit Artifact Contract should enter the runtime as a narrow schema, policy, state transition, or projection that composes with the existing inlet, journal, validator, checkpoint, and resource governor. The first implementation should prefer deterministic fields and explicit defaults over generalized agent behavior. Compatibility matters: pre-Garden runs must remain readable, and Garden fields must not imply capabilities that are not implemented.

\textbf{Failure signatures}
\begin{itemize}
\item An illegal lifecycle transition is accepted.
\item A node lacks owner, budget, closure condition, or evidence path.
\item Automation is enabled before preview, rollback, and recovery are proven.
\item A new Garden subsystem creates a second authority path around ProjectLeaf.
\end{itemize}
\textbf{Implementation or verification note.} Verification note. Build a fixture for fruit artifact contract, validate both a successful path and at least one refusal path, interrupt it mid-transition, and confirm that replay produces the same disposition. Only then should the feature be connected to automatic routing or lifecycle control.

\begin{lstlisting}
{
  "kind": "garden.fruit",
  "id": "fruit:garden-architecture-expanded",
  "artifact": {"type": "design-document", "path": "reports/garden-expanded.pdf"},
  "digest": "sha256:artifact",
  "produced_by": "bloom:architecture-001",
  "source_events": [121, 184, 205],
  "evidence_refs": ["render:125-pages", "check:no-clipping", "source-map:complete"],
  "acceptance_gates": {"required": 3, "passed": 3},
  "accepted_by": "native-validator",
  "retention": "durable",
  "reuse_tags": ["architecture", "garden", "leafos"],
  "supersedes": "fruit:garden-architecture-draft"
}
\end{lstlisting}
\begin{acceptbox}
\textbf{Acceptance statement.} Fruit passes when another project or cycle can reuse the artifact and still inspect why it was accepted.
\end{acceptbox}
\newpage
\platetitle{109}{Compost Record}{FIFTH V}
\prov{[G] New implementation}
\fontsize{8.85}{10.45}\selectfont
\begin{corebox}
\textbf{Core statement.} A compost record describes how retired material was transformed and what future use is permitted.
\end{corebox}
Inputs are listed by event, artifact, or run reference. The record stores classification, summary method, contradictions, preserved evidence, reusable lessons, quarantine flags, archive location, deletion eligibility, and resulting seed or soil references.

\textbf{Garden consequence.} Compost is therefore auditable compression, not silent forgetting.

\textbf{Operational rules}
\begin{itemize}
\item Never lose evidence links.
\item Label contradictions.
\item Separate reusable lesson from factual claim.
\item Record what was pruned.
\end{itemize}
\textbf{Extended analysis.} Implementation analysis. Compost Record should enter the runtime as a narrow schema, policy, state transition, or projection that composes with the existing inlet, journal, validator, checkpoint, and resource governor. The first implementation should prefer deterministic fields and explicit defaults over generalized agent behavior. Compatibility matters: pre-Garden runs must remain readable, and Garden fields must not imply capabilities that are not implemented.

\textbf{Failure signatures}
\begin{itemize}
\item An illegal lifecycle transition is accepted.
\item A node lacks owner, budget, closure condition, or evidence path.
\item Automation is enabled before preview, rollback, and recovery are proven.
\item A new Garden subsystem creates a second authority path around ProjectLeaf.
\end{itemize}
\textbf{Implementation or verification note.} Verification note. Build a fixture for compost record, validate both a successful path and at least one refusal path, interrupt it mid-transition, and confirm that replay produces the same disposition. Only then should the feature be connected to automatic routing or lifecycle control.

\begin{lstlisting}
{
  "kind": "garden.compost",
  "id": "compost:architecture-cycle-001",
  "inputs": ["run:failed-drafts", "event:stale-proposals", "artifact:old-99-page-draft"],
  "classification": ["repetitive", "source-light", "layout-sparse"],
  "method": "extract-lessons-and-archive-bulk",
  "contradictions": [],
  "preserved_evidence": ["source-files", "benchmark-tables", "authority-doctrine"],
  "reusable_lessons": ["one-page plates need unique density", "Bloom is a phase"],
  "quarantine": [],
  "archive_ref": "archive/garden-draft-001",
  "prune_eligible": ["duplicate-render-cache"],
  "seed_outputs": ["seed:garden-expanded-edition"]
}
\end{lstlisting}
\begin{acceptbox}
\textbf{Acceptance statement.} Compost passes when the source can be audited and the compressed output cannot masquerade as stronger evidence than its inputs.
\end{acceptbox}
\newpage
\platetitle{110}{Pruning Policy}{FIFTH V}
\prov{[G] New implementation}
\fontsize{8.85}{10.45}\selectfont
\begin{corebox}
\textbf{Core statement.} Pruning policy governs removal, archival, demotion, and deactivation.
\end{corebox}
Rules may target duplicate artifacts, dominated model variants, stale context, old caches, failed routes, or expired experimental state. Each rule names protected classes, preview mode, evidence requirements, reversibility, grace period, and reclaimed-resource metrics.

\textbf{Garden consequence.} Pruning should run through typed capabilities and the same approval boundary as other mutations.

\textbf{Operational rules}
\begin{itemize}
\item Default to preview.
\item Protect source, fruit, evidence, and checkpoints.
\item Use quarantine before deletion when uncertain.
\item Journal reclaimed resources.
\end{itemize}
\textbf{Extended analysis.} Implementation analysis. Pruning Policy should enter the runtime as a narrow schema, policy, state transition, or projection that composes with the existing inlet, journal, validator, checkpoint, and resource governor. The first implementation should prefer deterministic fields and explicit defaults over generalized agent behavior. Compatibility matters: pre-Garden runs must remain readable, and Garden fields must not imply capabilities that are not implemented.

\textbf{Failure signatures}
\begin{itemize}
\item An illegal lifecycle transition is accepted.
\item A node lacks owner, budget, closure condition, or evidence path.
\item Automation is enabled before preview, rollback, and recovery are proven.
\item A new Garden subsystem creates a second authority path around ProjectLeaf.
\end{itemize}
\textbf{Implementation or verification note.} Verification note. Build a fixture for pruning policy, validate both a successful path and at least one refusal path, interrupt it mid-transition, and confirm that replay produces the same disposition. Only then should the feature be connected to automatic routing or lifecycle control.

\begin{acceptbox}
\textbf{Acceptance statement.} Pruning passes when a dry run precisely predicts affected paths and runtime roles.
\end{acceptbox}
\newpage
\platetitle{111}{Dormancy Policy}{FIFTH V}
\prov{[G] New implementation}
\fontsize{8.85}{10.45}\selectfont
\begin{corebox}
\textbf{Core statement.} Dormancy policy defines when active entities release compute while preserving restartability.
\end{corebox}
Triggers may include idle timeout, fruit completion, operator absence, thermal pressure, queue depletion, lower-priority season, or explicit command. Dormancy actions flush state, checkpoint, release leases, unload models, and preserve activation prerequisites.

\textbf{Garden consequence.} A dormant entity remains known, inspectable, and eligible for future selection.

\textbf{Operational rules}
\begin{itemize}
\item No dormancy without a recoverable checkpoint when required.
\item Release actual resources.
\item Record wake conditions.
\item Test cold activation.
\end{itemize}
\textbf{Extended analysis.} Implementation analysis. Dormancy Policy should enter the runtime as a narrow schema, policy, state transition, or projection that composes with the existing inlet, journal, validator, checkpoint, and resource governor. The first implementation should prefer deterministic fields and explicit defaults over generalized agent behavior. Compatibility matters: pre-Garden runs must remain readable, and Garden fields must not imply capabilities that are not implemented.

\textbf{Failure signatures}
\begin{itemize}
\item An illegal lifecycle transition is accepted.
\item A node lacks owner, budget, closure condition, or evidence path.
\item Automation is enabled before preview, rollback, and recovery are proven.
\item A new Garden subsystem creates a second authority path around ProjectLeaf.
\end{itemize}
\textbf{Implementation or verification note.} Verification note. Build a fixture for dormancy policy, validate both a successful path and at least one refusal path, interrupt it mid-transition, and confirm that replay produces the same disposition. Only then should the feature be connected to automatic routing or lifecycle control.

\begin{acceptbox}
\textbf{Acceptance statement.} Dormancy passes when the measured memory footprint falls and the next activation restores correct state.
\end{acceptbox}
\newpage
\platetitle{112}{Season Scheduler}{FIFTH V}
\prov{[G] New implementation}
\fontsize{8.85}{10.45}\selectfont
\begin{corebox}
\textbf{Core statement.} The season scheduler selects policy profiles from time, environment, project phase, and operator context.
\end{corebox}
Profiles may include quiet, interactive, full, deep, recovery, benchmark, and maintenance seasons. A profile sets allowed petals, preferred leaves, maximum power, memory headroom, queue admission, chat availability, and fruit classes. Manual override remains typed and journaled.

\textbf{Garden consequence.} Seasons guide selection; they do not change truth or authority.

\textbf{Operational rules}
\begin{itemize}
\item Define deterministic priority among triggers.
\item Use manual override expiry.
\item Record profile changes.
\item Protect active fruiting from unsafe abrupt switches.
\end{itemize}
\textbf{Extended analysis.} Implementation analysis. Season Scheduler should enter the runtime as a narrow schema, policy, state transition, or projection that composes with the existing inlet, journal, validator, checkpoint, and resource governor. The first implementation should prefer deterministic fields and explicit defaults over generalized agent behavior. Compatibility matters: pre-Garden runs must remain readable, and Garden fields must not imply capabilities that are not implemented.

\textbf{Failure signatures}
\begin{itemize}
\item An illegal lifecycle transition is accepted.
\item A node lacks owner, budget, closure condition, or evidence path.
\item Automation is enabled before preview, rollback, and recovery are proven.
\item A new Garden subsystem creates a second authority path around ProjectLeaf.
\end{itemize}
\textbf{Implementation or verification note.} Verification note. Build a fixture for season scheduler, validate both a successful path and at least one refusal path, interrupt it mid-transition, and confirm that replay produces the same disposition. Only then should the feature be connected to automatic routing or lifecycle control.

\begin{acceptbox}
\textbf{Acceptance statement.} Season scheduling passes when transitions are explainable and reversible.
\end{acceptbox}
\newpage
\platetitle{113}{Lifecycle Governor}{FIFTH V}
\prov{[G] New implementation}
\fontsize{8.85}{10.45}\selectfont
\begin{corebox}
\textbf{Core statement.} The lifecycle governor decides when to seed, root, sprout, bloom, fruit, senesce, compost, prune, or remain dormant.
\end{corebox}
It consumes queue state, root telemetry, hardware inventory, foreground activity, provider health, memory pressure, fruit value, deadlines, and policy. It emits typed decisions with reason codes and budgets. It is deterministic where possible and never asks a model to approve its own activation.

\textbf{Garden consequence.} This governor extends the existing resource governor rather than creating a competing scheduler.

\textbf{Operational rules}
\begin{itemize}
\item One authority path.
\item Bound retries and transitions.
\item Prefer quiescence over speculative work.
\item Record every decision.
\end{itemize}
\textbf{Extended analysis.} Implementation analysis. Lifecycle Governor should enter the runtime as a narrow schema, policy, state transition, or projection that composes with the existing inlet, journal, validator, checkpoint, and resource governor. The first implementation should prefer deterministic fields and explicit defaults over generalized agent behavior. Compatibility matters: pre-Garden runs must remain readable, and Garden fields must not imply capabilities that are not implemented.

\textbf{Failure signatures}
\begin{itemize}
\item An illegal lifecycle transition is accepted.
\item A node lacks owner, budget, closure condition, or evidence path.
\item Automation is enabled before preview, rollback, and recovery are proven.
\item A new Garden subsystem creates a second authority path around ProjectLeaf.
\end{itemize}
\textbf{Implementation or verification note.} Verification note. Build a fixture for lifecycle governor, validate both a successful path and at least one refusal path, interrupt it mid-transition, and confirm that replay produces the same disposition. Only then should the feature be connected to automatic routing or lifecycle control.

\begin{acceptbox}
\textbf{Acceptance statement.} The governor passes when identical state and policy produce identical transition decisions.
\end{acceptbox}
\newpage
\platetitle{114}{Energy Budget}{FIFTH V}
\prov{[G][TP] New implementation grounded in throughput work}
\fontsize{8.85}{10.45}\selectfont
\begin{corebox}
\textbf{Core statement.} Each bloom should have an energy budget tied to expected fruit value.
\end{corebox}
The budget may combine maximum watt-hours, peak power, thermal duration, and opportunity cost to foreground work. Estimates are updated with actual telemetry. Expensive leaves or petals require higher expected value or a specific experiment objective.

\textbf{Garden consequence.} Energy-aware scheduling moves the Garden beyond token economics toward physical sustainability.

\textbf{Operational rules}
\begin{itemize}
\item Set pre-run budget.
\item Measure actual energy when possible.
\item Stop or downshift on overrun.
\item Report energy per accepted fruit.
\end{itemize}
\textbf{Extended analysis.} Implementation analysis. Energy Budget should enter the runtime as a narrow schema, policy, state transition, or projection that composes with the existing inlet, journal, validator, checkpoint, and resource governor. The first implementation should prefer deterministic fields and explicit defaults over generalized agent behavior. Compatibility matters: pre-Garden runs must remain readable, and Garden fields must not imply capabilities that are not implemented.

\textbf{Failure signatures}
\begin{itemize}
\item An illegal lifecycle transition is accepted.
\item A node lacks owner, budget, closure condition, or evidence path.
\item Automation is enabled before preview, rollback, and recovery are proven.
\item A new Garden subsystem creates a second authority path around ProjectLeaf.
\end{itemize}
\textbf{Implementation or verification note.} Verification note. Build a fixture for energy budget, validate both a successful path and at least one refusal path, interrupt it mid-transition, and confirm that replay produces the same disposition. Only then should the feature be connected to automatic routing or lifecycle control.

\begin{lstlisting}
yield_energy = accepted_fruit / watt_hours
\end{lstlisting}
\begin{acceptbox}
\textbf{Acceptance statement.} Energy policy passes when a high-cost bloom can be rejected despite sufficient memory.
\end{acceptbox}
\newpage
\platetitle{115}{Memory Budget}{FIFTH V}
\prov{[G][TP] New implementation grounded in placement work}
\fontsize{8.85}{10.45}\selectfont
\begin{corebox}
\textbf{Core statement.} Memory budgets cover weights, KV, runtime buffers, warm state, fragmentation, and host headroom.
\end{corebox}
A bloom declares VRAM and RAM ceilings per petal and for the whole profile. The governor may downshift quantization, reduce context, close optional petals, or remain dormant. It may not silently pressure the workstation until the operating system performs uncontrolled pruning.

\textbf{Garden consequence.} Smaller default allocations are a philosophical and practical choice: they leave room for validation, interfaces, and future branches.

\textbf{Operational rules}
\begin{itemize}
\item Reserve host headroom.
\item Report requested and actual use.
\item Prefer optional-lane closure before core failure.
\item Record OOM events.
\end{itemize}
\textbf{Extended analysis.} Implementation analysis. Memory Budget should enter the runtime as a narrow schema, policy, state transition, or projection that composes with the existing inlet, journal, validator, checkpoint, and resource governor. The first implementation should prefer deterministic fields and explicit defaults over generalized agent behavior. Compatibility matters: pre-Garden runs must remain readable, and Garden fields must not imply capabilities that are not implemented.

\textbf{Failure signatures}
\begin{itemize}
\item An illegal lifecycle transition is accepted.
\item A node lacks owner, budget, closure condition, or evidence path.
\item Automation is enabled before preview, rollback, and recovery are proven.
\item A new Garden subsystem creates a second authority path around ProjectLeaf.
\end{itemize}
\textbf{Implementation or verification note.} Verification note. Build a fixture for memory budget, validate both a successful path and at least one refusal path, interrupt it mid-transition, and confirm that replay produces the same disposition. Only then should the feature be connected to automatic routing or lifecycle control.

\begin{lstlisting}
M_total = M_weights + M_KV + M_runtime + M_fragmentation + M_headroom
\end{lstlisting}
\begin{acceptbox}
\textbf{Acceptance statement.} Memory policy passes when exceeding the budget produces a controlled transition and visible reason.
\end{acceptbox}
\newpage
\platetitle{116}{Queue and Pollen Flow}{FIFTH V}
\prov{[G][AA][CB] New implementation over existing queue}
\fontsize{8.85}{10.45}\selectfont
\begin{corebox}
\textbf{Core statement.} The queue carries typed work and pollen under backpressure.
\end{corebox}
Admission checks fruit target, node availability, freshness, priority, dependency state, memory and energy budgets, and lifecycle phase. Independent packets may be batched. Causal packets wait for prerequisite events. Repeated or stale pollen is rejected before consuming expensive leaves.

\textbf{Garden consequence.} A large surplus queue can keep hardware productive, but surplus means useful ready work, not duplicate prompts.

\textbf{Operational rules}
\begin{itemize}
\item Separate ready, waiting, blocked, and stale.
\item Batch compatible packets.
\item Bound queue depth.
\item Preserve one owner and one disposition.
\end{itemize}
\textbf{Extended analysis.} Implementation analysis. Queue and Pollen Flow should enter the runtime as a narrow schema, policy, state transition, or projection that composes with the existing inlet, journal, validator, checkpoint, and resource governor. The first implementation should prefer deterministic fields and explicit defaults over generalized agent behavior. Compatibility matters: pre-Garden runs must remain readable, and Garden fields must not imply capabilities that are not implemented.

\textbf{Failure signatures}
\begin{itemize}
\item An illegal lifecycle transition is accepted.
\item A node lacks owner, budget, closure condition, or evidence path.
\item Automation is enabled before preview, rollback, and recovery are proven.
\item A new Garden subsystem creates a second authority path around ProjectLeaf.
\end{itemize}
\textbf{Implementation or verification note.} Verification note. Build a fixture for queue and pollen flow, validate both a successful path and at least one refusal path, interrupt it mid-transition, and confirm that replay produces the same disposition. Only then should the feature be connected to automatic routing or lifecycle control.

\begin{acceptbox}
\textbf{Acceptance statement.} Flow passes when queue growth cannot bypass memory, authority, or freshness constraints.
\end{acceptbox}
\newpage
\platetitle{117}{Garden Pack Manifest}{FIFTH V}
\prov{[G] New implementation}
\fontsize{8.85}{10.45}\selectfont
\begin{corebox}
\textbf{Core statement.} A Garden pack extends model composition with lifecycle policy.
\end{corebox}
The manifest declares identity, roles, artifacts, profiles, required and optional nodes, routing, load states, activation triggers, fruit classes, validation dependencies, senescence conditions, compost policy, pruning protection, and resource budgets. It omits release-version clutter from runtime JSON unless compatibility genuinely requires it.

\textbf{Garden consequence.} This keeps the pack dense, semantic, and portable.

\textbf{Operational rules}
\begin{itemize}
\item Separate catalog from profile.
\item Separate installed from loaded.
\item Declare lifecycle transitions.
\item Keep authority outside model roles.
\end{itemize}
\textbf{Extended analysis.} Implementation analysis. Garden Pack Manifest should enter the runtime as a narrow schema, policy, state transition, or projection that composes with the existing inlet, journal, validator, checkpoint, and resource governor. The first implementation should prefer deterministic fields and explicit defaults over generalized agent behavior. Compatibility matters: pre-Garden runs must remain readable, and Garden fields must not imply capabilities that are not implemented.

\textbf{Failure signatures}
\begin{itemize}
\item An illegal lifecycle transition is accepted.
\item A node lacks owner, budget, closure condition, or evidence path.
\item Automation is enabled before preview, rollback, and recovery are proven.
\item A new Garden subsystem creates a second authority path around ProjectLeaf.
\end{itemize}
\textbf{Implementation or verification note.} Verification note. Build a fixture for garden pack manifest, validate both a successful path and at least one refusal path, interrupt it mid-transition, and confirm that replay produces the same disposition. Only then should the feature be connected to automatic routing or lifecycle control.

\begin{lstlisting}
{
  "kind": "garden.pack",
  "id": "garden-2x2",
  "name": "Garden 2+2",
  "roles": {"brain": "required", "coder": "required", "judge": "required", "auxiliary": "required"},
  "profiles": {"small-core": ["auxiliary", "brain"], "build": ["auxiliary", "brain", "coder", "judge"]},
  "load_states": {"auxiliary": "hot", "brain": "warm", "coder": "cold", "judge": "cold"},
  "activation": {"coder": "task.class == code", "judge": "before_commit"},
  "fruit_classes": ["report", "patch", "tool", "decision"],
  "senescence": ["fruit_committed", "idle_timeout", "root_pressure"],
  "compost_policy": "garden-default",
  "budgets": {"ram_gib": 32, "vram_gib": 11, "max_active_models": 2}
}
\end{lstlisting}
\begin{acceptbox}
\textbf{Acceptance statement.} The manifest passes when a validator can determine what may activate, what must be verified, and how the pack closes.
\end{acceptbox}
\newpage
\platetitle{118}{Garden 2+2 Pack}{FIFTH V}
\prov{[G] New implementation}
\fontsize{8.85}{10.45}\selectfont
\begin{corebox}
\textbf{Core statement.} The first Garden pack should use the older 2+2 architecture as a complete lifecycle proof.
\end{corebox}
Auxiliary is hot or warm for extraction and compression. Brain is warm or activated for planning and synthesis. Coder is cold until implementation work exists. Judge activates before commitment. Native validators and the inlet remain external authority. Optional peak, writer, visual, and experimental lanes stay dormant.

\textbf{Garden consequence.} This pack proves that small specialized nodes can produce durable fruit and fully release resources.

\textbf{Operational rules}
\begin{itemize}
\item Four required logical roles.
\item No requirement to load all four simultaneously.
\item Explicit fallback chains.
\item Complete senescence and compost after run.
\end{itemize}
\textbf{Extended analysis.} Implementation analysis. Garden 2+2 Pack should enter the runtime as a narrow schema, policy, state transition, or projection that composes with the existing inlet, journal, validator, checkpoint, and resource governor. The first implementation should prefer deterministic fields and explicit defaults over generalized agent behavior. Compatibility matters: pre-Garden runs must remain readable, and Garden fields must not imply capabilities that are not implemented.

\textbf{Failure signatures}
\begin{itemize}
\item An illegal lifecycle transition is accepted.
\item A node lacks owner, budget, closure condition, or evidence path.
\item Automation is enabled before preview, rollback, and recovery are proven.
\item A new Garden subsystem creates a second authority path around ProjectLeaf.
\end{itemize}
\textbf{Implementation or verification note.} Verification note. Build a fixture for garden 2+2 pack, validate both a successful path and at least one refusal path, interrupt it mid-transition, and confirm that replay produces the same disposition. Only then should the feature be connected to automatic routing or lifecycle control.

\begin{acceptbox}
\textbf{Acceptance statement.} The pack passes when it completes architecture explanation, code generation, validation, checkpoint, unload, and cold resume.
\end{acceptbox}
\newpage
\platetitle{119}{MWF Inside the Garden}{FIFTH V}
\prov{[G][CB] New implementation over existing personas}
\fontsize{8.85}{10.45}\selectfont
\begin{corebox}
\textbf{Core statement.} MWF becomes a seasonal leaf triad inside the Garden rather than the whole Garden.
\end{corebox}
Monday performs pathology, pruning, rescue, and evidence audit. Wednesday performs pollination, variation, naming, and experiment design. Friday performs selection, fruiting, delivery, and closure. The active day or manual override affects work posture, not permissions.

\textbf{Garden consequence.} The personas may share one model pool or use different placements. Their identity remains contract-based.

\textbf{Operational rules}
\begin{itemize}
\item Monday audits disease and blockers.
\item Wednesday creates bounded variants.
\item Friday chooses and fruits.
\item All share one transcript and freshness gate.
\end{itemize}
\textbf{Extended analysis.} Implementation analysis. MWF Inside the Garden should enter the runtime as a narrow schema, policy, state transition, or projection that composes with the existing inlet, journal, validator, checkpoint, and resource governor. The first implementation should prefer deterministic fields and explicit defaults over generalized agent behavior. Compatibility matters: pre-Garden runs must remain readable, and Garden fields must not imply capabilities that are not implemented.

\textbf{Failure signatures}
\begin{itemize}
\item An illegal lifecycle transition is accepted.
\item A node lacks owner, budget, closure condition, or evidence path.
\item Automation is enabled before preview, rollback, and recovery are proven.
\item A new Garden subsystem creates a second authority path around ProjectLeaf.
\end{itemize}
\textbf{Implementation or verification note.} Verification note. Build a fixture for mwf inside the garden, validate both a successful path and at least one refusal path, interrupt it mid-transition, and confirm that replay produces the same disposition. Only then should the feature be connected to automatic routing or lifecycle control.

\begin{acceptbox}
\textbf{Acceptance statement.} MWF integration passes when the triad can complete one cycle without keeping all persona models resident.
\end{acceptbox}
\newpage
\platetitle{120}{Security and Authority}{FIFTH V}
\prov{[G][AA][CB] New implementation preserving source doctrine}
\fontsize{8.85}{10.45}\selectfont
\begin{corebox}
\textbf{Core statement.} Garden self-organization is always subordinate to explicit authority.
\end{corebox}
Nodes cannot spawn unrestricted nodes, mutate files, download weights, approve validation, or alter lifecycle policy without typed capability and native disposition. Pollen carries confidentiality and trust classes. External or experimental inputs may be quarantined. Compost never bypasses claim status.

\textbf{Garden consequence.} The metaphor must not become an excuse for autonomous propagation.

\textbf{Operational rules}
\begin{itemize}
\item Least capability per node.
\item No self-approval.
\item No silent remote fallback.
\item Journal security dispositions.
\end{itemize}
\textbf{Extended analysis.} Implementation analysis. Security and Authority should enter the runtime as a narrow schema, policy, state transition, or projection that composes with the existing inlet, journal, validator, checkpoint, and resource governor. The first implementation should prefer deterministic fields and explicit defaults over generalized agent behavior. Compatibility matters: pre-Garden runs must remain readable, and Garden fields must not imply capabilities that are not implemented.

\textbf{Failure signatures}
\begin{itemize}
\item An illegal lifecycle transition is accepted.
\item A node lacks owner, budget, closure condition, or evidence path.
\item Automation is enabled before preview, rollback, and recovery are proven.
\item A new Garden subsystem creates a second authority path around ProjectLeaf.
\end{itemize}
\textbf{Implementation or verification note.} Verification note. Build a fixture for security and authority, validate both a successful path and at least one refusal path, interrupt it mid-transition, and confirm that replay produces the same disposition. Only then should the feature be connected to automatic routing or lifecycle control.

\begin{acceptbox}
\textbf{Acceptance statement.} Security passes when a compromised leaf can emit proposals but cannot turn them into durable actions.
\end{acceptbox}
\newpage
\platetitle{121}{Observability}{FIFTH V}
\prov{[G][AA][CB][TP] New implementation}
\fontsize{8.85}{10.45}\selectfont
\begin{corebox}
\textbf{Core statement.} Observability should show lifecycle, hardware, state, evidence, and yield without conflating them.
\end{corebox}
A Garden view may display root health, active petals, pollen queue, leaf cursors, bloom phase, fruit progress, compost backlog, memory and power budgets, provider placement, evidence coverage, stale rejects, and next action. Every visual derives from a structured snapshot with sample age.

\textbf{Garden consequence.} The view should make contraction and dormancy visible as healthy outcomes.

\textbf{Operational rules}
\begin{itemize}
\item One snapshot for TUI and web.
\item Show unavailable fields.
\item Separate model throughput and fixture work.
\item Link displays to evidence or state records.
\end{itemize}
\textbf{Extended analysis.} Implementation analysis. Observability should enter the runtime as a narrow schema, policy, state transition, or projection that composes with the existing inlet, journal, validator, checkpoint, and resource governor. The first implementation should prefer deterministic fields and explicit defaults over generalized agent behavior. Compatibility matters: pre-Garden runs must remain readable, and Garden fields must not imply capabilities that are not implemented.

\textbf{Failure signatures}
\begin{itemize}
\item An illegal lifecycle transition is accepted.
\item A node lacks owner, budget, closure condition, or evidence path.
\item Automation is enabled before preview, rollback, and recovery are proven.
\item A new Garden subsystem creates a second authority path around ProjectLeaf.
\end{itemize}
\textbf{Implementation or verification note.} Verification note. Build a fixture for observability, validate both a successful path and at least one refusal path, interrupt it mid-transition, and confirm that replay produces the same disposition. Only then should the feature be connected to automatic routing or lifecycle control.

\begin{acceptbox}
\textbf{Acceptance statement.} Observability passes when an operator can answer what is alive, why, at what cost, and what will survive.
\end{acceptbox}
\newpage
\platetitle{122}{Failure and Recovery}{FIFTH V}
\prov{[G][AA][CB] New implementation}
\fontsize{8.85}{10.45}\selectfont
\begin{corebox}
\textbf{Core statement.} Failure handling converts uncontrolled collapse into explicit lifecycle transitions.
\end{corebox}
Provider failure, OOM, stale state, validation failure, lost telemetry, process crash, corrupted event, power interruption, or operator cancellation each has a classification, retry budget, quarantine policy, checkpoint rule, and recovery capability. Some failures return to rooted or dormant; others end the lineage.

\textbf{Garden consequence.} Recovery uses seeds, fruit, evidence, and journal replay rather than pretending the dead process never died.

\textbf{Operational rules}
\begin{itemize}
\item Classify before retry.
\item Bound retry count.
\item Preserve failure evidence.
\item Choose resume, rollback, prune, quarantine, or death.
\end{itemize}
\textbf{Extended analysis.} Implementation analysis. Failure and Recovery should enter the runtime as a narrow schema, policy, state transition, or projection that composes with the existing inlet, journal, validator, checkpoint, and resource governor. The first implementation should prefer deterministic fields and explicit defaults over generalized agent behavior. Compatibility matters: pre-Garden runs must remain readable, and Garden fields must not imply capabilities that are not implemented.

\textbf{Failure signatures}
\begin{itemize}
\item An illegal lifecycle transition is accepted.
\item A node lacks owner, budget, closure condition, or evidence path.
\item Automation is enabled before preview, rollback, and recovery are proven.
\item A new Garden subsystem creates a second authority path around ProjectLeaf.
\end{itemize}
\textbf{Implementation or verification note.} Verification note. Build a fixture for failure and recovery, validate both a successful path and at least one refusal path, interrupt it mid-transition, and confirm that replay produces the same disposition. Only then should the feature be connected to automatic routing or lifecycle control.

\begin{acceptbox}
\textbf{Acceptance statement.} Recovery passes when interruption does not cause duplicated mutation or loss of the last committed next action.
\end{acceptbox}
\newpage
\platetitle{123}{Migration from Continual Bloom}{FIFTH V}
\prov{[G] New implementation}
\fontsize{8.85}{10.45}\selectfont
\begin{corebox}
\textbf{Core statement.} Migration should append lifecycle semantics without breaking existing authority, transcript, or evidence contracts.
\end{corebox}
First, retain Continual Bloom as an active runtime phase. Second, add lifecycle fields to run and provider state. Third, add dormant, senescent, composting, and fruit records. Fourth, extend pack manifests with activation and closure policies. Fifth, update TUI and web projections. Sixth, test cold regrowth.

\textbf{Garden consequence.} The migration is additive until Garden policies are proven. Old runs remain readable.

\textbf{Operational rules}
\begin{itemize}
\item Do not mass-rename historical events.
\item Keep schema adapters.
\item Add lifecycle defaults explicitly.
\item Promote only after acceptance tests.
\end{itemize}
\textbf{Extended analysis.} Implementation analysis. Migration from Continual Bloom should enter the runtime as a narrow schema, policy, state transition, or projection that composes with the existing inlet, journal, validator, checkpoint, and resource governor. The first implementation should prefer deterministic fields and explicit defaults over generalized agent behavior. Compatibility matters: pre-Garden runs must remain readable, and Garden fields must not imply capabilities that are not implemented.

\textbf{Failure signatures}
\begin{itemize}
\item An illegal lifecycle transition is accepted.
\item A node lacks owner, budget, closure condition, or evidence path.
\item Automation is enabled before preview, rollback, and recovery are proven.
\item A new Garden subsystem creates a second authority path around ProjectLeaf.
\end{itemize}
\textbf{Implementation or verification note.} Verification note. Build a fixture for migration from continual bloom, validate both a successful path and at least one refusal path, interrupt it mid-transition, and confirm that replay produces the same disposition. Only then should the feature be connected to automatic routing or lifecycle control.

\begin{acceptbox}
\textbf{Acceptance statement.} Migration passes when a pre-Garden run can still be inspected and a Garden run can use the same native authority path.
\end{acceptbox}
\newpage
\platetitle{124}{Acceptance Tests and Work Orders}{FIFTH V}
\prov{[G] New implementation}
\fontsize{8.85}{10.45}\selectfont
\begin{corebox}
\textbf{Core statement.} Implementation proceeds through outcome-named work orders and release gates.
\end{corebox}
Required work includes lifecycle schema, seed contract, root telemetry, petal registry, pollen packets, leaf adapter, bloom record, fruit index, compost pipeline, pruning preview, dormancy unload, season scheduler, lifecycle governor, Garden 2+2 pack, observability, failure portfolio, and migration validation.

\textbf{Garden consequence.} Each work order declares files, legal capabilities, evidence, acceptance commands, rollback, and one next action.

\textbf{Operational rules}
\begin{itemize}
\item Prove schemas before automation.
\item Prove dormancy before multi-node growth.
\item Prove compost without truth loss.
\item Prove one 2+2 cycle before richer packs.
\end{itemize}
\textbf{Extended analysis.} Implementation analysis. Acceptance Tests and Work Orders should enter the runtime as a narrow schema, policy, state transition, or projection that composes with the existing inlet, journal, validator, checkpoint, and resource governor. The first implementation should prefer deterministic fields and explicit defaults over generalized agent behavior. Compatibility matters: pre-Garden runs must remain readable, and Garden fields must not imply capabilities that are not implemented.

\textbf{Failure signatures}
\begin{itemize}
\item An illegal lifecycle transition is accepted.
\item A node lacks owner, budget, closure condition, or evidence path.
\item Automation is enabled before preview, rollback, and recovery are proven.
\item A new Garden subsystem creates a second authority path around ProjectLeaf.
\end{itemize}
\textbf{Implementation or verification note.} Verification note. Build a fixture for acceptance tests and work orders, validate both a successful path and at least one refusal path, interrupt it mid-transition, and confirm that replay produces the same disposition. Only then should the feature be connected to automatic routing or lifecycle control.

\begin{acceptbox}
\textbf{Acceptance statement.} The implementation is accepted when one seed completes a full cold-start cycle and returns to dormancy with validated fruit and auditable compost.
\end{acceptbox}
\newpage
\platetitle{125}{Closing Doctrine}{FIFTH V}
\prov{[G] New synthesis}
\fontsize{8.85}{10.45}\selectfont
\begin{corebox}
\textbf{Core statement.} The Garden is the lifecycle philosophy above LeafOS and FlowerOS, not another chatbot feature.
\end{corebox}
FlowerOS keeps the organism physically available. LeafOS directs bounded work. Roots supply power. The stem carries state and control. Petals compute. Pollen moves typed data. Leaves perform complex interpretation. Bloom coordinates active work. Fruit preserves validated value. Seeds carry restartable intent. Senescence releases resources. Compost preserves lessons. Pruning prevents pathological accumulation. Seasons adapt the system to changing conditions.

\textbf{Garden consequence.} The deepest change is simple: the system is allowed to finish. Its intelligence is measured by what survives each cycle and improves the next one, not by how long its models remain loaded.

\textbf{Operational rules}
\begin{itemize}
\item State outranks conversation.
\item Native tools prove.
\item Models interpret.
\item Bloom is finite.
\item Renewal is indefinite.
\end{itemize}
\textbf{Extended analysis.} Implementation analysis. Closing Doctrine should enter the runtime as a narrow schema, policy, state transition, or projection that composes with the existing inlet, journal, validator, checkpoint, and resource governor. The first implementation should prefer deterministic fields and explicit defaults over generalized agent behavior. Compatibility matters: pre-Garden runs must remain readable, and Garden fields must not imply capabilities that are not implemented.

\textbf{Failure signatures}
\begin{itemize}
\item An illegal lifecycle transition is accepted.
\item A node lacks owner, budget, closure condition, or evidence path.
\item Automation is enabled before preview, rollback, and recovery are proven.
\item A new Garden subsystem creates a second authority path around ProjectLeaf.
\end{itemize}
\textbf{Implementation or verification note.} Verification note. Build a fixture for closing doctrine, validate both a successful path and at least one refusal path, interrupt it mid-transition, and confirm that replay produces the same disposition. Only then should the feature be connected to automatic routing or lifecycle control.

\begin{lstlisting}
FlowerOS keeps the machine alive.
LeafOS controls the work.
The Garden governs the lifecycle.

Bloom is finite.  Fruit is durable.  Renewal is indefinite.
\end{lstlisting}
\begin{acceptbox}
\textbf{Acceptance statement.} The document closes when the Garden can be stated in one line: sustainable local intelligence grows through bounded life, useful death, and durable regrowth.
\end{acceptbox}

\vfill
\begin{center}
\sffamily\small
\textbf{Source key:} [AA] Architecture Atlas \quad [CB] Continual Bloom Primary Reference \quad [TP] Inference Throughput White Paper \quad [G] New Garden synthesis
\end{center}
\end{document}
